\section*{Technical unfolding of the OP transport architecture}
\begin{quote}\small
\noindent\textbf{Status.} \emph{Active proof-unfolding block.}\par
\noindent\textbf{Block function.} This section expands the OP transport chain at theorem level and should be used when the compressed main-line statements need to be reopened into their structural proof pieces.
\end{quote}
\addcontentsline{toc}{section}{Technical unfolding of the OP transport architecture}
This technical block unfolds the compressed OP branch-shell theorem used in the main line. It records the native closure machinery, the hardening lemmas, and the detailed branch/shell routed proofs that support the compressed theorem, so later sections can appeal either to the short main-text statement or to the full theorem-level form given here.

\paragraph{Dependency guide for the unfolding.}\label{par:op-unfolding-dependency-guide}
The technical order is now explicit: Theorem~\ref{thm:op15-closure} fixes the OP~1/OP~5 source-return pair, Theorem~\ref{thm:admissible-branch} promotes admissible branch transport, Theorem~\ref{thm:op3-shell} localizes shell/residual admissibility at OP~3, and Theorem~\ref{thm:combined-branch-shell} combines these layers into the routed visible return theorem. Proposition~\ref{prop:packet-phase-epsilon-screening} then records the packet/phase-lock/epsilon line as a diagnostic screening layer rather than a second closure theorem.

\subsection{Closure of the OP1--OP5 pair}

The OP-chain should not be read as a naive raw endpoint comparison. In the present framework, OP~1 is the inflow-side readout of a fixed structural route, whereas OP~5 belongs to the return sector of the same route and must be read through its reversal-stable part.

Let $\Gamma$ be a fixed structural route with reversed route $\overline{\Gamma}$, and let
\[
T_{\Gamma},T_{\overline{\Gamma}}:\mathcal{A}_{\mathrm{cl}}\to\mathcal{C}
\]
denote the forward and backward carrier transports on the admissible closure sector $\mathcal{A}_{\mathrm{cl}}\subseteq\mathcal{C}$. Let
\[
\rho_{\theta}:\mathcal{C}\to\mathcal{W}
\]
be the frame-dependent readout family.

We write
\[
\mathrm{OP1}=\rho_{\theta_{\mathrm{in}}}(c_{\Gamma}),
\qquad
\mathrm{OP5}_{\mathrm{fwd}}=\rho_{\theta_{\mathrm{fwd}}}(T_{\Gamma}c_{\Gamma}),
\qquad
\mathrm{OP5}_{\mathrm{bwd}}=\rho_{\theta_{\mathrm{bwd}}}(T_{\overline{\Gamma}}c_{\Gamma}),
\]
for a carrier $c_{\Gamma}\in\mathcal{A}_{\mathrm{cl}}$, and define the symmetrized return readout by
\[
\mathrm{OP5}_{\mathrm{sym}}
:=
\frac12\bigl(\mathrm{OP5}_{\mathrm{fwd}}+\mathrm{OP5}_{\mathrm{bwd}}\bigr).
\]

To express structural closure, we use the native closure carrier class
\[
K_{\mathrm{ncl}}:\mathcal{A}_{\mathrm{cl}}\to \mathcal{A}_{\mathrm{cl}}/\!\sim_{\mathrm{ncl}},
\]
which identifies carriers up to HC-compensated transport, admissible frame-shift, and admissible branch refinement of the same closure-bearing mode.

\begin{assumption}[Native OP~1--OP~5 closure assumptions]
We assume:
\begin{enumerate}[leftmargin=2em,label=(\roman*)]
\item \textbf{HC-confined transport:}
\[
T_{\Gamma}(\mathcal{A}_{\mathrm{cl}})\subseteq \mathcal{A}_{\mathrm{cl}},
\qquad
T_{\overline{\Gamma}}(\mathcal{A}_{\mathrm{cl}})\subseteq \mathcal{A}_{\mathrm{cl}};
\]

\item \textbf{two-pass class consistency:}
\[
K_{\mathrm{ncl}}(T_{\Gamma}c)=K_{\mathrm{ncl}}(c)=K_{\mathrm{ncl}}(T_{\overline{\Gamma}}c)
\qquad
\text{for all } c\in\mathcal{A}_{\mathrm{cl}};
\]

\item \textbf{closure-compatible readout:} there exists a map
\[
J:\mathcal{A}_{\mathrm{cl}}/\!\sim_{\mathrm{ncl}}\to\Lambda
\]
and an observable invariant
\[
\mathcal{I}:\mathcal{W}\to\Lambda
\]
such that
\[
\mathcal{I}(\rho_{\theta}(c))=J(K_{\mathrm{ncl}}(c))
\qquad
\text{for all } c\in\mathcal{A}_{\mathrm{cl}},\ \theta\in\Theta;
\]

\item \textbf{return-sector linearity:} $\mathcal{I}$ is linear on the admissible OP~5 return sector.
\end{enumerate}
\end{assumption}

\begin{theorem}[Closure of the OP~1--OP~5 pair]\label{thm:op15-closure}
Under the native OP~1--OP~5 closure assumptions,
\[
\mathcal{I}(\mathrm{OP1})
=
\mathcal{I}(\mathrm{OP5}_{\mathrm{sym}}).
\]
Hence the OP~1--OP~5 pair is closed at the level of the fixed closure invariant.
\end{theorem}

\begin{proof}
By closure-compatible readout,
\[
\mathcal{I}(\mathrm{OP1})
=
\mathcal{I}\!\left(\rho_{\theta_{\mathrm{in}}}(c_{\Gamma})\right)
=
J(K_{\mathrm{ncl}}(c_{\Gamma})).
\]
Likewise,
\[
\mathcal{I}(\mathrm{OP5}_{\mathrm{fwd}})
=
J(K_{\mathrm{ncl}}(T_{\Gamma}c_{\Gamma})),
\qquad
\mathcal{I}(\mathrm{OP5}_{\mathrm{bwd}})
=
J(K_{\mathrm{ncl}}(T_{\overline{\Gamma}}c_{\Gamma})).
\]
By two-pass class consistency,
\[
K_{\mathrm{ncl}}(T_{\Gamma}c_{\Gamma})
=
K_{\mathrm{ncl}}(c_{\Gamma})
=
K_{\mathrm{ncl}}(T_{\overline{\Gamma}}c_{\Gamma}),
\]
and therefore
\[
\mathcal{I}(\mathrm{OP5}_{\mathrm{fwd}})
=
\mathcal{I}(\mathrm{OP1})
=
\mathcal{I}(\mathrm{OP5}_{\mathrm{bwd}}).
\]
Using linearity on the admissible return sector,
\[
\mathcal{I}(\mathrm{OP5}_{\mathrm{sym}})
=
\frac12\Bigl(
\mathcal{I}(\mathrm{OP5}_{\mathrm{fwd}})
+
\mathcal{I}(\mathrm{OP5}_{\mathrm{bwd}})
\Bigr)
=
\mathcal{I}(\mathrm{OP1}),
\]
which proves the claim.
\end{proof}

\begin{corollary}[Odd return defect is closure-invisible]
Define
\[
\Delta_{\mathrm{OP5}}
:=
\frac12\bigl(\mathrm{OP5}_{\mathrm{fwd}}-\mathrm{OP5}_{\mathrm{bwd}}\bigr).
\]
Then
\[
\mathcal{I}(\Delta_{\mathrm{OP5}})=0.
\]
\end{corollary}

\begin{proof}
By linearity,
\[
\mathcal{I}(\Delta_{\mathrm{OP5}})
=
\frac12\Bigl(
\mathcal{I}(\mathrm{OP5}_{\mathrm{fwd}})
-
\mathcal{I}(\mathrm{OP5}_{\mathrm{bwd}})
\Bigr)=0,
\]
since the theorem gives equality of the two return-side invariant values.
\end{proof}

\begin{remark}
The theorem does not require raw observable equality
\[
\mathrm{OP1}=\mathrm{OP5}_{\mathrm{sym}}.
\]
It only requires equality after passage to the fixed closure invariant. Thus source and return may remain visibly distinct while still forming a closed OP~1--OP~5 pair.
\end{remark}

\paragraph{Consequence for the OP reading rule.}
The theorem therefore replaces raw endpoint comparison by invariant source-return comparison: OP~1 and OP~5 may remain visibly separated in the measurement frame while still carrying one and the same closure-bearing mode, provided the two-pass transport remains class-preserving.

\subsection{Branchwise refinement of the closed OP1--OP5 pair}

Once the OP~1--OP~5 pair has been closed at the invariant level, branchwise effects should be read as secondary refinements of that closed source-return architecture rather than as replacements for it. In particular, branching does not create closure ex nihilo; it only distributes the observable realization of a closure-bearing mode across admissible residual routes.

Let $\mathcal{B}$ be a family of admissible residual branches associated with the fixed route $\Gamma$. For each branch $b\in\mathcal{B}$, let
\[
T_{\Gamma}^{(b)},T_{\overline{\Gamma}}^{(b)}:\mathcal{A}_{\mathrm{cl}}\to\mathcal{C}
\]
denote the branchwise forward and backward carrier transports, and let
\[
\mathrm{OP5}^{(b)}_{\mathrm{fwd}},
\qquad
\mathrm{OP5}^{(b)}_{\mathrm{bwd}}
\]
be the corresponding branchwise return readouts. Define
\[
\mathrm{OP5}^{(b)}_{\mathrm{sym}}
:=
\frac12\bigl(\mathrm{OP5}^{(b)}_{\mathrm{fwd}}+\mathrm{OP5}^{(b)}_{\mathrm{bwd}}\bigr).
\]

We say that a branch $b$ is \emph{closure-admissible} if
\[
K_{\mathrm{ncl}}(T_{\Gamma}^{(b)}c)
=
K_{\mathrm{ncl}}(c)
=
K_{\mathrm{ncl}}(T_{\overline{\Gamma}}^{(b)}c)
\qquad
\text{for all admissible } c\in\mathcal{A}_{\mathrm{cl}}.
\]

\begin{proposition}[Branchwise refinement of OP closure]\label{prop:op-branchwise-refinement}
Assume the OP~1--OP~5 pair is closed in the sense of the previous theorem. If a branch $b\in\mathcal{B}$ is closure-admissible, then
\[
\mathcal{I}(\mathrm{OP1})
=
\mathcal{I}(\mathrm{OP5}^{(b)}_{\mathrm{sym}}).
\]
Hence every closure-admissible branch refines the same invariantly closed source-return pair.
\end{proposition}

\begin{proof}
Since $b$ is closure-admissible, the branchwise forward and backward transports preserve the same native closure carrier class as the source carrier. By closure-compatible readout, the observable invariant $\mathcal{I}$ therefore takes the same value on $\mathrm{OP1}$, $\mathrm{OP5}^{(b)}_{\mathrm{fwd}}$, and $\mathrm{OP5}^{(b)}_{\mathrm{bwd}}$. Linearity on the admissible return sector then yields
\[
\mathcal{I}(\mathrm{OP5}^{(b)}_{\mathrm{sym}})
=
\mathcal{I}(\mathrm{OP1}).
\]
\end{proof}

\begin{remark}
The proposition shows that admissible branching does not introduce a new closure criterion. The closure test remains the same; branchwise analysis only asks which residual routes preserve the already fixed native closure class.
\end{remark}

\begin{corollary}[Diagnostic status of non-admissible branches]
If a branch fails closure-admissibility, then it does not by itself refute closure of the OP~1--OP~5 pair. It only shows that this particular residual refinement leaves the native closure class.
\end{corollary}

\begin{proof}
Closure of the OP~1--OP~5 pair is decided at the level of the fixed two-pass source-return architecture. A non-admissible branch records failure of class preservation for that branchwise refinement, not necessarily failure of the underlying closed pair itself.
\end{proof}

\subsection{Observable routing consequence of closure-admissible branches}

Once the OP~1--OP~5 pair has been closed at the invariant level and closure-admissible branches have been identified, the remaining question is how these admissible branches contribute to the observable return signal. In the present framework, this contribution belongs to the readout layer and should therefore be modeled as a routed aggregation of branchwise return realizations.

Let $\mathcal{B}_{\mathrm{adm}}\subseteq\mathcal{B}$ denote the set of closure-admissible branches. For each $b\in\mathcal{B}_{\mathrm{adm}}$, let
\[
\mathrm{OP5}^{(b)}_{\mathrm{sym}}\in\mathcal{W}
\]
be the branchwise symmetrized return readout. Assume that the observable readout space $\mathcal{W}$ carries an admissible additive structure.

\begin{definition}[Admissible routing weights]
A family of admissible routing weights is a map
\[
w:\mathcal{B}_{\mathrm{adm}}\to [0,\infty).
\]
The corresponding admissible routed return readout is
\[
\mathrm{OP5}_{\mathrm{rout}}
:=
\sum_{b\in\mathcal{B}_{\mathrm{adm}}} w(b)\,\mathrm{OP5}^{(b)}_{\mathrm{sym}},
\]
whenever the sum is well-defined in $\mathcal{W}$.
\end{definition}

\begin{proposition}[Invariant value of the admissible routed return]\label{prop:op-admissible-routing}
Assume that:
\begin{enumerate}[leftmargin=2em,label=(\roman*)]
\item every branch $b\in\mathcal{B}_{\mathrm{adm}}$ is closure-admissible;
\item the closure invariant $\mathcal{I}$ is linear on the admissible routed sector;
\item the routed sum defining $\mathrm{OP5}_{\mathrm{rout}}$ exists.
\end{enumerate}
Then
\[
\mathcal{I}(\mathrm{OP5}_{\mathrm{rout}})
=
\Bigl(\sum_{b\in\mathcal{B}_{\mathrm{adm}}} w(b)\Bigr)\mathcal{I}(\mathrm{OP1}).
\]
\end{proposition}

\begin{proof}
By branchwise closure-admissibility, the previous proposition gives
\[
\mathcal{I}\!\left(\mathrm{OP5}^{(b)}_{\mathrm{sym}}\right)
=
\mathcal{I}(\mathrm{OP1})
\qquad
\text{for all } b\in\mathcal{B}_{\mathrm{adm}}.
\]
Using linearity of $\mathcal{I}$ on the admissible routed sector, we obtain
\[
\mathcal{I}(\mathrm{OP5}_{\mathrm{rout}})
=
\mathcal{I}\!\left(
\sum_{b\in\mathcal{B}_{\mathrm{adm}}} w(b)\,\mathrm{OP5}^{(b)}_{\mathrm{sym}}
\right)
=
\sum_{b\in\mathcal{B}_{\mathrm{adm}}} w(b)\,
\mathcal{I}\!\left(\mathrm{OP5}^{(b)}_{\mathrm{sym}}\right).
\]
Substituting the branchwise closure identity yields
\[
\mathcal{I}(\mathrm{OP5}_{\mathrm{rout}})
=
\sum_{b\in\mathcal{B}_{\mathrm{adm}}} w(b)\,\mathcal{I}(\mathrm{OP1})
=
\Bigl(\sum_{b\in\mathcal{B}_{\mathrm{adm}}} w(b)\Bigr)\mathcal{I}(\mathrm{OP1}),
\]
as claimed.
\end{proof}

\begin{corollary}[Normalized admissible routing]
If, in addition, the admissible routing weights are normalized so that
\[
\sum_{b\in\mathcal{B}_{\mathrm{adm}}} w(b)=1,
\]
then
\[
\mathcal{I}(\mathrm{OP5}_{\mathrm{rout}})
=
\mathcal{I}(\mathrm{OP1}).
\]
\end{corollary}

\begin{proof}
This is immediate from the proposition.
\end{proof}

\begin{definition}[Admissible routing defect]
Define the admissible routing defect by
\[
\mathfrak{d}_{\mathrm{rout}}
:=
\mathcal{I}(\mathrm{OP5}_{\mathrm{rout}})
-
\mathcal{I}(\mathrm{OP1}).
\]
In the normalized case, one has
\[
\mathfrak{d}_{\mathrm{rout}}=0.
\]
More generally,
\[
\mathfrak{d}_{\mathrm{rout}}
=
\bigl(W_{\mathrm{adm}}-1\bigr)\mathcal{I}(\mathrm{OP1}),
\qquad
W_{\mathrm{adm}}:=\sum_{b\in\mathcal{B}_{\mathrm{adm}}} w(b).
\]
\end{definition}

\begin{remark}
Thus unnormalized routing does not destroy closure content; it rescales its observable aggregate. The distinction between structural closure and observable weight remains intact.
\end{remark}

\subsection{Mixed routed support and the final OP closure diagnostic}

Visible routed support need not come entirely from closure-admissible branches. To close the OP analysis, one must therefore separate the admissible routed sector from the remaining visible support.

Let $\mathcal{B}_{\mathrm{vis}}$ be the set of visible branches with routing weights $w(b)\ge 0$, and split it as
\[
\mathcal{B}_{\mathrm{vis}}
=
\mathcal{B}_{\mathrm{adm}}\sqcup\mathcal{B}_{\mathrm{rem}},
\qquad
\mathcal{B}_{\mathrm{rem}}:=\mathcal{B}_{\mathrm{vis}}\setminus\mathcal{B}_{\mathrm{adm}}.
\]
Assume that every visible branch $b\in\mathcal{B}_{\mathrm{vis}}$ carries a symmetrized return representative $\mathrm{OP5}^{(b)}_{\mathrm{sym}}\in\mathcal{W}$ and define the total routed return by
\[
\mathrm{OP5}_{\mathrm{tot}}
:=
\sum_{b\in\mathcal{B}_{\mathrm{vis}}} w(b)\,\mathrm{OP5}^{(b)}_{\mathrm{sym}},
\]
whenever the sum exists.

\begin{proposition}[Decomposition of the total routed invariant]
Assume that $\mathcal{I}$ is linear on the visible routed sector. Then
\[
\mathcal{I}(\mathrm{OP5}_{\mathrm{tot}})
=
W_{\mathrm{adm}}\,\mathcal{I}(\mathrm{OP1}) + R_{\mathrm{rem}},
\]
where
\[
W_{\mathrm{adm}}:=\sum_{b\in\mathcal{B}_{\mathrm{adm}}} w(b),
\qquad
R_{\mathrm{rem}}:=\sum_{b\in\mathcal{B}_{\mathrm{rem}}} w(b)\,\mathcal{I}\!\left(\mathrm{OP5}^{(b)}_{\mathrm{sym}}\right).
\]
\end{proposition}

\begin{proof}
By linearity,
\[
\mathcal{I}(\mathrm{OP5}_{\mathrm{tot}})
=
\sum_{b\in\mathcal{B}_{\mathrm{adm}}} w(b)\,\mathcal{I}\!\left(\mathrm{OP5}^{(b)}_{\mathrm{sym}}\right)
+
\sum_{b\in\mathcal{B}_{\mathrm{rem}}} w(b)\,\mathcal{I}\!\left(\mathrm{OP5}^{(b)}_{\mathrm{sym}}\right).
\]
The first sum equals $W_{\mathrm{adm}}\,\mathcal{I}(\mathrm{OP1})$ by the admissible-routing proposition. The second sum is exactly $R_{\mathrm{rem}}$.
\end{proof}

\begin{corollary}[Final OP closure mismatch decomposition]
Define the total OP closure mismatch by
\[
\mathfrak{m}_{\mathrm{OP}}
:=
\mathcal{I}(\mathrm{OP5}_{\mathrm{tot}})-\mathcal{I}(\mathrm{OP1}).
\]
Then
\[
\mathfrak{m}_{\mathrm{OP}}
=
\bigl(W_{\mathrm{adm}}-1\bigr)\mathcal{I}(\mathrm{OP1}) + R_{\mathrm{rem}}.
\]
\end{corollary}

\begin{proof}
Subtract $\mathcal{I}(\mathrm{OP1})$ from the decomposition in the previous proposition.
\end{proof}

\begin{corollary}[Exact closure criterion for the visible routed sector]
If
\[
W_{\mathrm{adm}}=1
\qquad\text{and}\qquad
R_{\mathrm{rem}}=0,
\]
then
\[
\mathcal{I}(\mathrm{OP5}_{\mathrm{tot}})=\mathcal{I}(\mathrm{OP1}).
\]
In particular, exact routed OP closure holds whenever the visible routed support is normalized on the admissible sector and carries no residual non-admissible contribution at the invariant level.
\end{corollary}

\begin{proof}
Under the stated conditions, the mismatch formula gives $\mathfrak{m}_{\mathrm{OP}}=0$.
\end{proof}

\begin{remark}[Closed form of the OP analysis]
The OP analysis is thereby closed in three layers: native source-return closure first, admissible branch refinement second, and visible routed mismatch only third. Observable source-return disagreement is therefore confined to routing normalization and residual non-admissible support once the structural pair itself has been closed.
\end{remark}



\subsection{Admissible-branch theorem}
\label{sec:admissible-branch-theorem}

With the OP~1--OP~5 pair closed, the next question is no longer whether branches appear, but which of them preserve the already fixed native closure architecture. The present theorem turns that distinction into a theorem-level criterion.

Let $\Gamma$ be a fixed structural route with closed OP~1--OP~5 source-return architecture, and let $\mathcal{B}$ be a family of branchwise residual refinements attached to the OP~3 interface. For each branch $b\in\mathcal{B}$, let
\[
T_{\Gamma}^{(b)},T_{\overline{\Gamma}}^{(b)}:\mathcal{A}_{\mathrm{cl}}\to\mathcal{C}
\]
denote the corresponding branchwise forward and backward transports.

\begin{definition}[Visible support]
A branch $b\in\mathcal{B}$ is called visibly supported if its routed observable weight satisfies $w(b)>0$.
\end{definition}

\begin{theorem}[Admissible-branch theorem]\label{thm:admissible-branch}
Assume the OP transport theorem. Then the following hold.
\begin{enumerate}[leftmargin=2em,label=(\roman*)]
\item If a branch $b\in\mathcal{B}$ is closure-admissible, then
\[
\mathcal{I}(\mathrm{OP1})=\mathcal{I}(\mathrm{OP5}^{(b)}_{\mathrm{sym}}).
\]
\item If a branch fails closure-admissibility, then the failure is detected first at the level of native closure-class preservation, not at the level of raw observable mismatch.
\item Visible support does not imply closure-admissibility.
\item Within the closed OP architecture, branchwise structural relevance is determined first by closure-admissibility and only secondarily by routed visible support.
\end{enumerate}
\end{theorem}

\begin{proof}
Part (i) is Proposition~\ref{prop:op-branchwise-refinement}. Part (ii) is the contrapositive reading of the same proposition together with the native closure-class criterion. Part (iii) follows because visible support is controlled by $w(b)$, whereas admissibility is controlled by preservation of $K_{\mathrm{ncl}}$. Part (iv) is the structural reading of (i)--(iii).
\end{proof}

\begin{corollary}[Hierarchy of branch diagnostics]
Branchwise diagnostics should be read in the order
\[
\text{closure-admissibility}
\;>\;
\text{branch defect}
\;>\;
\text{visible support}.
\]
\end{corollary}

\subsection{OP3 shell/residual theorem}
\label{sec:op3-shell-theorem}

We now localize shell- and residual-sensitive effects at the OP~3 interface itself: they become structurally relevant only when they cease to preserve the native closure class there.

Let
\[
c_{3}:=A_{3}(A_{2}(c_{\Gamma}))
\]
be the OP~3 carrier attached to the closed OP architecture. Let $\mathcal{R}$ be the residual parameter space and $\mathcal{S}$ the shell index set. For each $(r,s)\in\mathcal{R}\times\mathcal{S}$, let
\[
\mathfrak{T}_{r,s}:\mathcal{C}\to\mathcal{C}
\]
be the corresponding shell/residual refinement acting at the OP~3 interface, and define
\[
c_{3}^{(r,s)}:=\mathfrak{T}_{r,s}(c_{3}).
\]

\begin{definition}[OP3 shell-admissibility]
A shell/residual refinement $\mathfrak{T}_{r,s}$ is called OP~3 shell-admissible if
\[
K_{\mathrm{ncl}}(c_{3}^{(r,s)})=K_{\mathrm{ncl}}(c_{3}).
\]
\end{definition}

\begin{theorem}[OP3 shell/residual theorem]
\label{thm:op3-shell}
Assume the OP transport theorem. Then the following hold.
\begin{enumerate}[leftmargin=2em,label=(\roman*)]
\item If $\mathfrak{T}_{r,s}$ is OP~3 shell-admissible, then
\[
\mathcal{I}(\rho_{\theta}(c_{3}^{(r,s)}))=\mathcal{I}(\rho_{\theta}(c_{3}))=\mathcal{I}(\mathrm{OP1})
\]
for every admissible frame phase $\theta$.
\item Shell-consistency defects are first OP~3 admissibility defects, not immediate failures of the whole OP chain.
\item Inverse-shell behavior is secondary unless it changes the native closure carrier class.
\item Raw shell-sensitive mismatch has no primary structural force unless it survives passage to failure of native closure-class preservation.
\end{enumerate}
\end{theorem}

\begin{proof}
Part (i) follows because $\mathcal{I}$ factors through $K_{\mathrm{ncl}}$ and the intermediate invariant bridge identifies $\mathcal{I}(\mathrm{OP3})$ with $\mathcal{I}(\mathrm{OP1})$. Parts (ii)--(iv) are the corresponding structural reading rules once OP~3 is fixed as the residual/shell interface.
\end{proof}

\begin{remark}[Compressed reading of OP~3 after the extracted QAM proof run]
The theorem above remains the native OP statement of the OP~3 shell/residual interface. After the extracted QAM tool pass, however, it no longer needs to be reread only in shell-first language. An OP~3 shell-admissible refinement may now be passed through the shell-certificate tool, the bundled OP~3 closure tool, and the obstruction-stratification tools from the QAM layer. In plain language: OP~3 shell admissibility can now be read either directly at the OP interface or through the compressed certified-pass / localized-obstruction dichotomy of Corollary~\ref{cor:qam-op3-proof-run-use}.
\end{remark}

\begin{proposition}[Compressed rereading of the OP~3 shell/residual theorem through extracted QAM tools]
\label{prop:op3-shell-reread-through-qam-tools}
Assume the faithful translation theorem for the admissible OP fragment, Theorem~\ref{thm:qam-translation}, the direct kernel constructor of Definition~\ref{def:qam-direct-kernel-constructor}, and the bridge conditions of Theorem~\ref{thm:qam-direct-kernel-constructor}. Let $\mathfrak{T}_{r,s}$ be an OP~3 shell-admissible refinement. Then the OP~3 shell/residual theorem may be reread in the following compressed form:
\begin{enumerate}[leftmargin=2em,label=(\roman*)]
\item the translated shell passage activates the bundled OP~3 closure tool, so the shell step may be continued kernel-first through one common closure-bearing carrier;
\item any first structurally relevant shell defect must already appear as the first localized obstruction in the extracted QAM stratification;
\item inverse-shell or raw shell-sensitive visible mismatch remains secondary unless it propagates to failure of the shell certificate package, equivalently to loss of the bundled closure carrier.
\end{enumerate}
\end{proposition}

\begin{proof}
Translate the admissible OP fragment by Theorem~\ref{thm:qam-translation}. Because $\mathfrak{T}_{r,s}$ is OP~3 shell-admissible, the corresponding translated passage falls under the shell-certificate side of the first narrow iff package, Theorem~\ref{thm:qam-shell-iff-package}. The bundled OP~3 closure tool, Theorem~\ref{thm:qam-op3-shell-closure}, then supplies clause~(i): the visible shell passage may be replaced by one common closure-bearing kernel/readout carrier.

Clause~(ii) is the negative reading of the same package. If a first structurally relevant shell defect occurs, then by Lemma~\ref{lem:qam-shell-obstruction-stratification} it must already appear as the first localized obstruction stratum. Clause~(iii) follows because Lemma~\ref{lem:qam-readout-kernel-obstruction-compression} compresses every residual branch/shell-preserving failure to loss of one common carrier. Hence visible inverse-shell or raw shell mismatch is still secondary unless it survives passage to that compressed obstruction form.
\end{proof}

\subsection{Combined branch-shell transport theorem}
\label{sec:combined-branch-shell-theorem}

The final routed return may now be read as a combined branch-shell aggregation. The theorem below isolates exactly when that visible composite support still preserves the invariant OP closure already established in the main line.

Let $\mathcal{B}_{\mathrm{adm}}$ be the admissible branch set and let
\[
\mathcal{F}_{\mathrm{sh}}^{(b)}\subseteq \mathcal{R}\times\mathcal{S}
\]
denote the shell-admissible refinement family attached to branch $b\in\mathcal{B}_{\mathrm{adm}}$ at the OP~3 interface. For every admissible branch $b$ and every admissible refinement $(r,s)\in\mathcal{F}_{\mathrm{sh}}^{(b)}$, let
\[
\mathrm{OP5}^{(b,r,s)}_{\mathrm{sym}}\in\mathcal{W}
\]
denote the corresponding symmetrized return readout.

\begin{definition}[Branch-shell routed return]
Let $w(b,r,s)\ge 0$ be a routed weight assigned to each branch-shell admissible realization. Whenever the sum is well-defined in $\mathcal{W}$, define
\[
\mathrm{OP5}_{\mathrm{bs\text{-}rout}}
:=
\sum_{b\in\mathcal{B}_{\mathrm{adm}}}
\ \sum_{(r,s)\in\mathcal{F}_{\mathrm{sh}}^{(b)}}
 w(b,r,s)\,\mathrm{OP5}^{(b,r,s)}_{\mathrm{sym}}.
\]
\end{definition}

\begin{theorem}[Combined branch-shell transport theorem]
\label{thm:combined-branch-shell}
Assume the OP transport theorem, the admissible-branch theorem, and the OP~3 shell/residual theorem. Assume further that $\mathcal{I}$ is linear on the admissible branch-shell routed sector. Then:
\begin{enumerate}[leftmargin=2em,label=(\roman*)]
\item For every admissible branch $b\in\mathcal{B}_{\mathrm{adm}}$ and every shell-admissible refinement $(r,s)\in\mathcal{F}_{\mathrm{sh}}^{(b)}$,
\[
\mathcal{I}(\mathrm{OP5}^{(b,r,s)}_{\mathrm{sym}})=\mathcal{I}(\mathrm{OP1}).
\]
\item If the routed weights are normalized so that
\[
\sum_{b\in\mathcal{B}_{\mathrm{adm}}}\ \sum_{(r,s)\in\mathcal{F}_{\mathrm{sh}}^{(b)}} w(b,r,s)=1,
\]
then
\[
\mathcal{I}(\mathrm{OP5}_{\mathrm{bs\text{-}rout}})=\mathcal{I}(\mathrm{OP1}).
\]
\item A visibly composite return assembled from several admissible branches and several shell-admissible refinements may remain structurally closure-preserving at invariant level.
\item Failure of closure in the branch-shell routed sector occurs first through loss of branch admissibility or loss of OP~3 shell-admissibility, not through visible multiplicity alone.
\end{enumerate}
\end{theorem}

\begin{proof}
Part (i) combines Theorem~\ref{thm:admissible-branch} with Proposition~\ref{prop:op3-shell-reread-through-qam-tools}. Admissible branches preserve the OP invariant, and once the OP~3 shell passage is rerouted through the bundled OP~3 closure tool it preserves that same invariant inside each branch through one common closure-bearing carrier. Part (ii) follows from linearity of $\mathcal{I}$ on the admissible branch-shell routed sector. Parts (iii)--(iv) are the structural reading of (i)--(ii): visible multiplicity alone is not structural non-closure as long as admissible branch passage and the compressed OP~3 shell package remain intact.
\end{proof}


\subsection{First carrier-first OP5 interface}
\label{sec:carrier-first-op5-interface}

After the combined branch-shell theorem, OP~5 should no longer be read as a final visible endpoint that must be tested from scratch. The first terminal question is whether the symmetrized return readout still belongs to the same native closure-bearing carrier decision already fixed by the source-return theorem and by the admissible branch-shell passage.

\begin{remark}[Carrier-first reading of OP~5 after branch-shell transport]
The visible OP~5 packet-, phase-, and flicker-diagnostics remain secondary at this stage. Once a branch-shell realization is admissible, OP~5 should first be read as a return-side query against one common closure-bearing carrier; only afterwards should one inspect visible packet-scale deviations.
\end{remark}

\begin{proposition}[First carrier-first OP~5 interface]
\label{prop:carrier-first-op5-interface}
Assume the closure theorem for the OP~1--OP~5 pair, Theorem~\ref{thm:op15-closure}, and the combined branch-shell transport theorem, Theorem~\ref{thm:combined-branch-shell}. Let $(b,r,s)$ be a branch-shell admissible realization and let
\[
\mathrm{OP5}^{(b,r,s)}_{\mathrm{sym}}\in\mathcal{W}
\]
denote the associated symmetrized return readout. Then the terminal OP~5 decision has the following carrier-first form.
\begin{enumerate}[leftmargin=2em,label=(\roman*)]
\item \textbf{Certified carrier pass.}
\[
\mathcal{I}\!\left(\mathrm{OP5}^{(b,r,s)}_{\mathrm{sym}}\right)
=
\mathcal{I}(\mathrm{OP1}).
\]
Hence OP~5 is to be read as a return-side readout of one common closure-bearing carrier class rather than as an independent closure test.
\item \textbf{First localized obstruction.} If a terminal structural failure is claimed at OP~5, then its first relevant source lies upstream of raw terminal visibility: it must already arise through loss of branch admissibility, loss of OP~3 shell-admissibility, or loss of the corresponding carrier/readout package before an independent OP~5 failure can be asserted.
\end{enumerate}
In particular, the terminal OP~5 layer is the first carrier-level decision interface after the branch-shell pass, not a new autonomous closure theorem.
\end{proposition}

\begin{proof}
By Theorem~\ref{thm:combined-branch-shell}, every branch-shell admissible realization satisfies
\[
\mathcal{I}\!\left(\mathrm{OP5}^{(b,r,s)}_{\mathrm{sym}}\right)=\mathcal{I}(\mathrm{OP1}).
\]
By Theorem~\ref{thm:op15-closure}, the OP~1 source readout and the symmetrized OP~5 return are already tied to one and the same native closure decision at invariant level. Thus the admissible branch-shell return is not a new endpoint test; it is a return-side reading of that same closure-bearing carrier class. This proves part~(i). Part~(ii) is the structural reading of the same fact: once the terminal invariant is fixed by the shared carrier decision, any genuinely structural failure visible at the terminal side must have entered earlier, namely through loss of admissibility on the branch-shell side or through failure of the carrier/readout package that supports the branch-shell pass. Therefore OP~5 is the first carrier-level decision interface after the branch-shell passage, not an independent closure theorem.
\end{proof}

\begin{corollary}[Carrier-first routed OP~5 decision principle]
\label{cor:carrier-first-routed-op5}
Assume the hypotheses of Proposition~\ref{prop:carrier-first-op5-interface} and let the branch-shell routed return
\[
\mathrm{OP5}_{\mathrm{bs\text{-}rout}}
\]
be defined with normalized routed weights. Then:
\begin{enumerate}[leftmargin=2em,label=(\roman*)]
\item the entire routed return is to be read first through one common closure-bearing carrier decision;
\item visible terminal multiplicity does not by itself create a new OP~5 closure problem;
\item any structural routed failure must already originate at the first localized obstruction of some branch-shell candidate rather than at terminal aggregation alone.
\end{enumerate}
\end{corollary}

\begin{proof}
Normalization together with Theorem~\ref{thm:combined-branch-shell} yields
\[
\mathcal{I}(\mathrm{OP5}_{\mathrm{bs\text{-}rout}})=\mathcal{I}(\mathrm{OP1}).
\]
Thus the routed return still belongs to the same invariant closure decision as the source readout. The remaining statements are the carrier-first reading of this equality: visible routed multiplicity is secondary, and any structural failure must already arise at the first localized obstruction of one of the contributing branch-shell realizations.
\end{proof}

\subsection{Carrier-first packet/phase-lock/epsilon compatibility layer}
\label{sec:packet-phase-epsilon-compatibility}

After the carrier-first OP~5 interface, the remaining packet-, phase-lock-, and epsilon-trap diagnostics should now be read as a visible support layer for the admissible return side of one common closure-bearing carrier. They do not replace the admissibility theorems or the carrier-first OP~5 decision; rather, they provide a structured non-escape corridor in which visible packet-scale deviations remain subordinate to the same return-side carrier decision.

\begin{definition}[Packet/phase-lock/epsilon compatible branch-shell channel]
A branch-shell realization $(b,r,s)$ is called \emph{packet/phase-lock/epsilon compatible} if the following hold simultaneously:
\begin{enumerate}[leftmargin=2em,label=(\roman*)]
\item the branch $b$ is closure-admissible;
\item the refinement $(r,s)$ is OP~3 shell-admissible;
\item the OP~1-side bridge normalization lies inside the contractive epsilon tube of Theorem~\ref{thm:working-op1-epsilon-trapped-bridge};
\item the OP~3-side visible shell ratio lies inside the anti-phase epsilon corridor preserved as \texttt{thm:working-shell-anti-phase-epsilon-trap} in the historical-witness PDF;
\item the packet-wave, breathing-envelope, and frame-flicker residuals are treated as visible-support diagnostics rather than as primary closure data.
\end{enumerate}
\end{definition}

\begin{remark}
The first two conditions are structural; the next two are non-escape conditions on the OP~1 entry side and the OP~3 shell side; the final condition places the remaining packet-scale evidence at the observable support layer. This is the natural compressed reading of the current packet/phase/epsilon analysis after OP closure hardening and after the carrier-first OP~5 interface has fixed the terminal decision level.
\end{remark}

\begin{corollary}[Carrier-first packet/phase-lock/epsilon compatibility]
\label{cor:packet-phase-epsilon-compatibility}
Assume the combined branch-shell transport theorem together with Proposition~\ref{prop:carrier-first-op5-interface}. If a branch-shell realization $(b,r,s)$ is packet/phase-lock/epsilon compatible, then the following hold.
\begin{enumerate}[leftmargin=2em,label=(\roman*)]
\item the visible packet-, breathing-envelope-, and frame-flicker diagnostics are compatible with one and the same invariantly closed source-return channel because they sit on the return side of one common closure-bearing carrier;
\item the OP~1 bridge side and the OP~3 shell side both satisfy a non-escape criterion inside their respective local epsilon corridors;
\item no visible packet- or phase-lock mismatch obtained inside those corridors is by itself evidence against structural OP closure or against the carrier-first OP~5 decision;
\item the first structurally relevant failure can occur only through loss of branch admissibility or loss of OP~3 shell-admissibility.
\end{enumerate}
\end{corollary}

\begin{proof}
By branch admissibility and OP~3 shell-admissibility, the combined branch-shell transport theorem gives
\[
\mathcal{I}(\mathrm{OP5}^{(b,r,s)}_{\mathrm{sym}})=\mathcal{I}(\mathrm{OP1}).
\]
By Proposition~\ref{prop:carrier-first-op5-interface}, this equality is already the carrier-first terminal OP~5 decision: the admissible return readout belongs to one common closure-bearing carrier and is not a new autonomous endpoint test. The OP~1 epsilon-trapped bridge theorem supplies the local non-escape statement on the source-entry side, while the shell anti-phase epsilon-trap supplies the corresponding non-escape statement for the visible shell ratio at the OP~3 interface. The packet-wave, breathing-envelope, and frame-flicker diagnostics are explicitly treated as visible-support diagnostics, so any mismatch they show inside the admissible corridors remains subordinate to that carrier-first decision and to the visible/structural mismatch separation already built into the OP tool layer. Therefore the visible diagnostics remain compatible with the same invariantly closed return-side carrier channel, and the first structurally relevant failure can only arise by leaving the admissible branch-shell sector itself.
\end{proof}

\begin{remark}[Current mathematical role of the packet/phase/epsilon line]
The present packet/phase-lock/epsilon line is therefore no longer best read as an alternative closure proof or as a second OP~5 test. Its mathematical role is to provide local corridor evidence that the currently observed OP~1 and OP~5 diagnostics can sit inside one admissible return-side carrier channel without forcing a structural mismatch. In that sense, it is the visible non-escape support layer beneath the carrier-first OP~5 decision and the next admissible-branch hardening layer.
\end{remark}

\begin{proposition}[Carrier-first packet/phase-lock/epsilon screening principle]
\label{prop:packet-phase-epsilon-screening}
Assume the combined branch-shell transport theorem together with Proposition~\ref{prop:carrier-first-op5-interface}. Then packet/phase-lock/epsilon compatibility defines a screening criterion for candidate admissible branch-shell support inside the carrier-first OP~5 architecture rather than a second closure theorem. More precisely, if a realization $(b,r,s)$ is packet/phase-lock/epsilon compatible, then it lies inside the diagnostic preselection corridor of the admissible branch-shell architecture and remains subordinate to the same return-side carrier decision; conversely, loss of packet/phase-lock/epsilon compatibility is not by itself sufficient to prove structural non-closure unless it propagates to loss of branch admissibility or loss of OP~3 shell-admissibility.
\end{proposition}

\begin{proof}
The first statement is immediate from Corollary~\ref{cor:packet-phase-epsilon-compatibility}: compatibility places the realization inside a joint non-escape corridor on the OP~1 bridge side and the OP~3 shell side while leaving the invariant closure value unchanged, and Proposition~\ref{prop:carrier-first-op5-interface} then fixes the terminal reading as one common return-side carrier decision. Hence the diagnostic role of the packet/phase/epsilon conditions is to preselect a corridor of candidate admissible support beneath that carrier-first terminal reading, not to introduce an independent closure criterion.

For the converse direction, failure of packet-, phase-lock-, or epsilon-corridor compatibility may occur either because a visible-support diagnostic leaves its controlled regime or because a structural admissibility condition fails. By the same corollary and by the carrier-first OP~5 interface, only the latter is primary for OP closure: visible mismatch inside or outside the packet-scale corridor is not yet structural non-closure unless it reflects loss of branch admissibility or loss of OP~3 shell-admissibility. Therefore packet/phase-lock/epsilon diagnostics define a screening layer rather than a standalone closure theorem or a second OP~5 endpoint test.
\end{proof}

\begin{remark}
This screening principle is the cleanest current mathematical reading of the old packet/phase-lock/epsilon diagnostics. They are now neither discarded nor promoted beyond their scope: they act as a visible preselection gate for the admissible branch-shell sector under one carrier-first OP~5 decision, while the actual structural test remains the admissibility machinery of the OP core.
\end{remark}


\subsection{Diagnostic readout hierarchy beneath the carrier-first OP~5 decision}
\label{sec:diagnostic-readout-hierarchy}

After the carrier-first OP~5 interface and the packet/phase-lock/epsilon screening layer, the remaining visible packet, breathing-envelope, frame-flicker, and helix pictures should be read in a strict order. They are not parallel closure criteria. They are progressively more visible readout surfaces of one and the same carrier decision.

\begin{remark}[Three diagnostic layers]
The current OP~5-side diagnostic stack is therefore to be read in the following descending order of structural weight.
\begin{enumerate}[leftmargin=2em,label=(\roman*)]
\item \textbf{Carrier-decision layer.} The first structural question is whether the return side still belongs to one common closure-bearing carrier, as fixed by Proposition~\ref{prop:carrier-first-op5-interface}.
\item \textbf{Screening layer.} Packet/phase-lock/epsilon compatibility then determines whether the visible return sits inside an admissible non-escape corridor beneath that carrier decision.
\item \textbf{Visible-shape layer.} Packet-wave, breathing-envelope, frame-flicker, and helix pictures come last: they visualize how one admissible or nearly admissible return is seen, but they do not outrank the carrier decision or the screening corridor.
\end{enumerate}
\end{remark}

\begin{lemma}[First diagnostic-readout hierarchy lemma]
\label{lem:diagnostic-readout-hierarchy}
Assume Proposition~\ref{prop:carrier-first-op5-interface} and Proposition~\ref{prop:packet-phase-epsilon-screening}. Let $(b,r,s)$ be a branch-shell realization with associated visible packet-, breathing-envelope-, frame-flicker-, and helix-type diagnostics. Then the current OP-side readout hierarchy has the following form.
\begin{enumerate}[leftmargin=2em,label=(\roman*)]
\item \textbf{Carrier priority.} If the realization belongs to one common closure-bearing carrier, then no visible packet-, flicker-, or helix-level multiplicity by itself overrides that carrier decision.
\item \textbf{Screening priority.} Packet/phase-lock/epsilon compatibility is the first visible screening layer beneath the carrier decision; loss at the pure visible-shape level is not yet structural loss unless it propagates to loss of this screening layer or to loss of the carrier package itself.
\item \textbf{Visible-shape subordination.} The packet-wave, breathing-envelope, frame-flicker, and helix pictures are therefore subordinate readout surfaces. They can localize, rank, and compare visible deviations, but they do not constitute an independent closure theorem and do not outrank the carrier-first OP~5 decision.
\end{enumerate}
In particular, the mathematically correct order of evidential force is
\[
\text{carrier decision}
\;>\;
\text{packet/phase/epsilon screening}
\;>\;
\text{visible flicker/helix picture}.
\]
\end{lemma}

\begin{proof}
Proposition~\ref{prop:carrier-first-op5-interface} fixes the first terminal structural question: the return side is read first through one common closure-bearing carrier, and a structural failure must already arise upstream of raw terminal visibility. Proposition~\ref{prop:packet-phase-epsilon-screening} then places packet/phase-lock/epsilon compatibility strictly beneath that terminal carrier decision as a diagnostic preselection corridor rather than as a second closure theorem. The remaining packet-wave, breathing-envelope, frame-flicker, and helix pictures are by construction even further downstream, because they are explicitly treated as visible-support diagnostics in Definition~3.116 and Corollary~\ref{cor:packet-phase-epsilon-compatibility}. Hence visible-shape multiplicity may refine the observable profile of a candidate return, but it cannot outrank the carrier decision and cannot become structural loss unless the screening corridor or the carrier package itself fails. This proves the stated hierarchy.
\end{proof}

\begin{corollary}[Visible flicker/helix subordination principle]
\label{cor:flicker-helix-subordination}
Within the current carrier-first OP~5 architecture, any packet-wave, breathing-envelope, frame-flicker, or helix mismatch is to be read first as visible-shape evidence beneath the carrier decision and beneath the packet/phase-lock/epsilon screening layer. Such a mismatch becomes structurally decisive only if it can be shown to force loss of packet/phase-lock/epsilon compatibility or loss of the common closure-bearing carrier itself.
\end{corollary}

\begin{proof}
This is the direct operational reading of Lemma~\ref{lem:diagnostic-readout-hierarchy}. The visible-shape diagnostics are subordinate readout surfaces; hence they can become structurally decisive only by propagating upward into the screening layer or the carrier layer.
\end{proof}

\subsection{First upstream propagation principle for visible-shape mismatch}
\label{sec:first-upstream-propagation}

\begin{remark}[Upstream propagation gate]
Under the current carrier-first OP~5 architecture, visible packet-wave, breathing-envelope, frame-flicker, and helix mismatches have no immediate right to overturn the structural reading. They may move upward only by forcing failure of the next superior layer: first the packet/phase-lock/epsilon screening corridor, and only after that the common closure-bearing carrier package itself.
\end{remark}

\begin{lemma}[First upstream propagation lemma]
\label{lem:first-upstream-propagation}
Assume Proposition~\ref{prop:carrier-first-op5-interface}, Proposition~\ref{prop:packet-phase-epsilon-screening}, Lemma~\ref{lem:diagnostic-readout-hierarchy}, and Lemma~\ref{lem:qam-shell-obstruction-stratification}. Let $(b,r,s)$ be a branch-shell realization with a visible packet-wave, breathing-envelope, frame-flicker, or helix mismatch. Then the current upstream propagation rule has the following form.
\begin{enumerate}[leftmargin=2em,label=(\roman*)]
\item \textbf{No direct jump.} A pure visible-shape mismatch does not by itself propagate directly to structural non-closure or directly overturn the carrier-first OP~5 decision.
\item \textbf{First gate: screening.} If the mismatch does \emph{not} force loss of packet/phase-lock/epsilon compatibility, then it remains subordinate visible-shape evidence and the carrier decision is unchanged.
\item \textbf{Second gate: carrier package.} If the mismatch \emph{does} force loss of packet/phase-lock/epsilon compatibility while the admissible branch-shell sector and the common closure-bearing carrier package remain intact, then the mismatch propagates only to screening-layer loss and is still not yet a structural OP failure.
\item \textbf{Structural propagation.} A visible-shape mismatch becomes structurally decisive only when it can be shown to propagate further to one of the upstream obstruction strata of Lemma~\ref{lem:qam-shell-obstruction-stratification}, namely loss of admissible branch passage, loss of shell-admissible refinement, or loss of the common readout/kernel carrier package.
\end{enumerate}
In particular, the current architecture permits only the ordered propagation chain
\[
\begin{aligned}
\text{visible flicker/helix mismatch}
&\;\Longrightarrow\;
\text{possible screening loss}\\
&\;\Longrightarrow\;
\text{possible carrier/package loss}\\
&\;\Longrightarrow\;
\text{structural failure}.
\end{aligned}
\]
No shorter upward propagation is presently licensed.
\end{lemma}

\begin{proof}
Lemma~\ref{lem:diagnostic-readout-hierarchy} already fixes the hierarchy of evidential force: visible-shape diagnostics lie strictly below the packet/phase-lock/epsilon screening layer, which itself lies below the carrier-first OP~5 decision. Hence a pure visible mismatch cannot jump directly to structural failure. Proposition~\ref{prop:packet-phase-epsilon-screening} gives the first gate: if packet/phase-lock/epsilon compatibility is preserved, then the mismatch remains inside the diagnostic corridor and does not alter the carrier reading. Even if that screening corridor is lost, Proposition~\ref{prop:carrier-first-op5-interface} still prevents immediate structural failure unless the common return-side carrier package itself is lost. Finally, Lemma~\ref{lem:qam-shell-obstruction-stratification} identifies the only current upstream structural failure modes: branch-stratum loss, shell-stratum loss, or readout/kernel carrier-package loss. Therefore the only admissible upstream propagation route is the ordered one stated above.
\end{proof}

\begin{corollary}[Visible mismatch non-propagation default]
\label{cor:visible-mismatch-nonpropagation}
Within the current carrier-first OP~5 architecture, every packet-wave, breathing-envelope, frame-flicker, or helix mismatch is to be treated by default as diagnostically local evidence. It may be upgraded only after an explicit witness of screening-layer loss or of upstream carrier/package loss has been exhibited.
\end{corollary}

\begin{proof}
This is the operational reformulation of Lemma~\ref{lem:first-upstream-propagation}. In the absence of an explicit upward propagation witness, visible-shape mismatch remains subordinate local evidence.
\end{proof}


\subsection{First diagnostic witness principle for visible-shape mismatch}
\label{sec:first-diagnostic-witness-principle}

\begin{remark}[Diagnostic witness discipline]
Under the current carrier-first OP~5 architecture, visible packet-wave, breathing-envelope, frame-flicker, and helix mismatches count as admissible upstream evidence only when they come equipped with a witness that is strong enough to certify failure of the next superior layer. Pure visual irregularity, even if persistent or large, is not yet such a witness.
\end{remark}

\begin{lemma}[First diagnostic witness lemma]
\label{lem:first-diagnostic-witness}
Assume Proposition~\ref{prop:packet-phase-epsilon-screening}, Proposition~\ref{prop:carrier-first-op5-interface}, Lemma~\ref{lem:diagnostic-readout-hierarchy}, and Lemma~\ref{lem:first-upstream-propagation}. Let $(b,r,s)$ be a branch-shell realization with a visible packet-wave, breathing-envelope, frame-flicker, or helix mismatch. Then the current witness discipline has the following form.
\begin{enumerate}[leftmargin=2em,label=(\roman*)]
\item \textbf{Screening witness threshold.} A visible mismatch is a legitimate witness of screening-layer loss only if it can be shown to force failure of packet/phase-lock/epsilon compatibility; otherwise it remains diagnostically local evidence.
\item \textbf{Carrier witness threshold.} A visible mismatch is a legitimate witness of carrier/package loss only if it can be shown to force, beyond screening loss, failure of the common closure-bearing carrier package itself.
\item \textbf{No witness by appearance alone.} Mere amplitude, persistence, oscillatory complexity, geometric deformation, or helix/flicker asymmetry of the visible mismatch is not by itself an admissible witness of screening-layer loss or of carrier/package loss.
\item \textbf{Admissible witness route.} Every currently admissible visible witness must therefore factor through one of the two certified upstream gates of Lemma~\ref{lem:first-upstream-propagation}: first explicit screening loss, and only then, if present, explicit carrier/package loss.
\end{enumerate}
Hence the present architecture admits visible mismatch only as \emph{witness-bearing} evidence and not as a free structural verdict.
\end{lemma}

\begin{proof}
Proposition~\ref{prop:packet-phase-epsilon-screening} defines the first superior layer above raw visible-shape mismatch. Therefore a visible mismatch can witness screening-layer loss only by forcing failure of packet/phase-lock/epsilon compatibility itself; if compatibility persists, the mismatch remains inside the diagnostic corridor and cannot count as screening evidence. Proposition~\ref{prop:carrier-first-op5-interface} defines the next superior layer: even after screening loss, the structural carrier reading remains unchanged unless the common closure-bearing carrier package is lost. Thus a visible mismatch can witness carrier/package loss only by forcing precisely that second failure as well. Lemma~\ref{lem:diagnostic-readout-hierarchy} excludes any direct elevation of raw visible-shape features to a superior structural layer, and Lemma~\ref{lem:first-upstream-propagation} identifies the only currently licensed upward route as the ordered screening-to-carrier route. Consequently amplitude, persistence, or geometric shape alone are not admissible witnesses; only a mismatch that certifiably drives one of those gates may count as an upstream witness.
\end{proof}

\begin{corollary}[Visible witness admissibility rule]
\label{cor:visible-witness-admissibility}
Within the current carrier-first OP~5 architecture, every visible packet-wave, breathing-envelope, frame-flicker, or helix mismatch must be classified first as either
\begin{enumerate}[leftmargin=2em,label=(\alph*)]
\item a merely local diagnostic irregularity,
\item a certified witness of screening-layer loss,
\item or a certified witness of carrier/package loss.
\end{enumerate}
No fourth class of immediately structural visible witness is presently admitted.
\end{corollary}

\begin{proof}
This is the direct classification consequence of Lemma~\ref{lem:first-diagnostic-witness}. A visible mismatch is either not strong enough to force a superior-layer failure, or it forces screening loss, or it forces carrier/package loss after screening loss. By the witness discipline and by the upstream propagation rule, there is no currently licensed direct jump from raw visible appearance to immediate structural verdict.
\end{proof}


\subsection{Screening-witness compression beneath the carrier-first OP~5 decision}
\label{sec:screening-witness-compression}

\begin{remark}[Compression of visible witness types]
Once a visible packet-wave, breathing-envelope, frame-flicker, or helix mismatch has crossed the witness threshold of Lemma~\ref{lem:first-diagnostic-witness}, the present architecture no longer needs to keep its visible shape as a separate structural species at the screening level. Unless the same witness already certifies carrier/package loss, its upstream content compresses to one common statement: the packet/phase-lock/epsilon screening corridor has failed.
\end{remark}

\begin{lemma}[Screening-witness compression lemma]
\label{lem:screening-witness-compression}
Assume Proposition~\ref{prop:packet-phase-epsilon-screening}, Proposition~\ref{prop:carrier-first-op5-interface}, Lemma~\ref{lem:first-diagnostic-witness}, and Corollary~\ref{cor:visible-witness-admissibility}. Let $(b,r,s)$ be a branch-shell realization with a visible packet-wave, breathing-envelope, frame-flicker, or helix mismatch that is a certified witness of screening-layer loss but not yet a certified witness of carrier/package loss. Then the upstream structural content of that witness compresses to a single screening-loss form, namely the failure of packet/phase-lock/epsilon compatibility. More precisely:
\begin{enumerate}[leftmargin=2em,label=(\roman*)]
\item \textbf{Visible-shape indifference at screening level.} Whether the witness is first expressed through packet-wave mismatch, breathing-envelope drift, frame-flicker asymmetry, or helix deformation is irrelevant for its screening-level structural role.
\item \textbf{Unique compressed content.} The only currently licensed upstream content of such a witness is that the packet/phase-lock/epsilon screening corridor fails.
\item \textbf{No parallel screening species.} There is presently no second independent screening-loss class attached specifically to packet-wave, envelope, flicker, or helix shape once the witness has been admitted.
\item \textbf{Carrier exception.} A visible witness escapes this compression only when it upgrades further to a certified witness of carrier/package loss.
\end{enumerate}
Hence, below the carrier threshold, all currently admissible visible screening witnesses collapse to one common screening-loss normal form.
\end{lemma}

\begin{proof}
By Lemma~\ref{lem:first-diagnostic-witness}, a visible mismatch is admitted upward only when it certifiably forces failure of the next superior layer. Proposition~\ref{prop:packet-phase-epsilon-screening} identifies that next superior layer uniquely as packet/phase-lock/epsilon compatibility. Therefore, once a visible mismatch qualifies as a witness of screening-layer loss, its upstream meaning is already fixed: it testifies precisely to failure of that screening corridor. Corollary~\ref{cor:visible-witness-admissibility} allows no additional currently licensed class of screening witness beyond that certified screening loss or the stronger carrier/package loss. Proposition~\ref{prop:carrier-first-op5-interface} then separates the only exception: a witness that also forces loss of the common closure-bearing carrier package belongs to the higher carrier class instead of remaining in the compressed screening class. Thus all admitted visible witnesses that stop below the carrier threshold share one and the same screening-loss content, independently of the visible shape through which they first appeared.
\end{proof}

\begin{corollary}[Screening-witness normal form]
\label{cor:screening-witness-normal-form}
Within the current carrier-first OP~5 architecture, every admitted visible witness that does not yet certify carrier/package loss may be cited in proofs simply as a witness of \emph{screening-loss normal form}. Its packet-wave, breathing-envelope, frame-flicker, or helix realization is diagnostically informative, but not structurally independent at the screening level.
\end{corollary}

\begin{proof}
This is the direct proof-usage reformulation of Lemma~\ref{lem:screening-witness-compression}. Once admitted and still below the carrier threshold, the witness has exactly one currently licensed upstream content: loss of the packet/phase-lock/epsilon screening corridor.
\end{proof}


\subsection{Carrier-witness separation beneath compressed screening loss}
\label{sec:carrier-witness-separation}

\begin{remark}[Two witness grades after screening compression]
After Lemma~\ref{lem:screening-witness-compression}, every admitted visible witness sits in one of two remaining grades. Either it stays below the carrier threshold and is therefore only a witness of screening-loss normal form, or it upgrades further and certifies that the common return-side carrier package itself is no longer available. The present architecture does not license a mixed intermediate status between those two grades.
\end{remark}

\begin{lemma}[Carrier-witness separation lemma]
\label{lem:carrier-witness-separation}
Assume Proposition~\ref{prop:packet-phase-epsilon-screening}, Proposition~\ref{prop:carrier-first-op5-interface}, Lemma~\ref{lem:first-diagnostic-witness}, Lemma~\ref{lem:screening-witness-compression}, and Corollary~\ref{cor:screening-witness-normal-form}. Let $(b,r,s)$ be a branch-shell realization with an admitted visible witness. Then exactly one of the following two cases occurs:
\begin{enumerate}[leftmargin=2em,label=(\roman*)]
\item \textbf{Screening-only witness.} The witness compresses to screening-loss normal form, while the common closure-bearing return-side carrier package remains available. In this case the witness certifies only screening loss and is not an admissible witness of carrier/package loss.
\item \textbf{Carrier-upgraded witness.} The witness forces loss of the common closure-bearing return-side carrier package. In this case it is an admissible witness of carrier/package loss and no longer belongs merely to the screening-loss normal form.
\end{enumerate}
Moreover, there is no third currently licensed status in which an admitted visible witness would simultaneously remain only screening-level and yet already count as a carrier/package witness.
\end{lemma}

\begin{proof}
By Lemma~\ref{lem:first-diagnostic-witness}, every admitted visible witness must force loss of some superior layer. Lemma~\ref{lem:screening-witness-compression} and Corollary~\ref{cor:screening-witness-normal-form} show that, below the carrier threshold, all such witnesses have one and the same upstream content: failure of packet/phase-lock/epsilon compatibility. Proposition~\ref{prop:carrier-first-op5-interface} identifies the only stronger alternative, namely loss of the common closure-bearing return-side carrier package itself. Hence an admitted visible witness either stops at the compressed screening-loss level or upgrades to carrier/package loss. These two statuses are mutually exclusive because the first expressly retains the common carrier package, whereas the second destroys it. Therefore exactly one of the two cases above occurs, and no mixed intermediate category is presently licensed.
\end{proof}

\begin{corollary}[Carrier-witness upgrade criterion]
\label{cor:carrier-witness-upgrade-criterion}
Within the current carrier-first OP~5 architecture, an admitted visible witness upgrades from screening-loss normal form to a genuine carrier/package witness if and only if its upstream content cannot be realized through any common closure-bearing return-side carrier package. Equivalently, as long as such a common carrier remains available, the witness must be read only as screening-loss normal form.
\end{corollary}

\begin{proof}
The forward direction is exactly case (ii) of Lemma~\ref{lem:carrier-witness-separation}. The converse direction is case (i): if a common closure-bearing return-side carrier package still exists, then the witness has not crossed the carrier threshold and therefore remains only a screening witness in normal form.
\end{proof}


\subsection{Witness-to-carrier escalation beyond screening compression}
\label{sec:witness-to-carrier-escalation}

\begin{remark}[Escalation gate after witness separation]
After Lemma~\ref{lem:carrier-witness-separation}, a visible witness reaches carrier grade exactly when its admitted upstream content can no longer be absorbed by screening-loss normal form alone. In the current architecture there is therefore only one licensed upgrade mechanism: every common closure-bearing return-side carrier package compatible with the compressed screening witness fails, so the witness is forced out of screening-only status and into carrier/package loss.
\end{remark}

\begin{lemma}[Witness-to-carrier escalation lemma]
\label{lem:witness-to-carrier-escalation}
Assume Proposition~\ref{prop:carrier-first-op5-interface}, Lemma~\ref{lem:first-diagnostic-witness}, Lemma~\ref{lem:screening-witness-compression}, Lemma~\ref{lem:carrier-witness-separation}, and Corollary~\ref{cor:carrier-witness-upgrade-criterion}. Let $(b,r,s)$ be a branch-shell realization with an admitted visible witness $w$. Then the following are equivalent:
\begin{enumerate}[leftmargin=2em,label=(\roman*)]
\item $w$ upgrades from screening-loss normal form to a genuine carrier/package witness;
\item there exists no common closure-bearing return-side carrier package through which the upstream content of $w$ can still be realized while screening-loss normal form remains the only failed superior layer;
\item the upstream content of $w$ forces loss of the common closure-bearing carrier package rather than mere loss of packet/phase-lock/epsilon compatibility.
\end{enumerate}
In particular, escalation occurs exactly when compressed screening loss is no longer a sufficient upstream normal form for the admitted witness.
\end{lemma}

\begin{proof}
By Lemma~\ref{lem:carrier-witness-separation}, an admitted visible witness has only two currently licensed grades: screening-only or carrier-upgraded. Corollary~\ref{cor:carrier-witness-upgrade-criterion} identifies the upgrade threshold precisely: upgrade occurs if and only if no common closure-bearing return-side carrier package remains available. Lemma~\ref{lem:screening-witness-compression} shows that, below that threshold, the witness has exactly one upstream content, namely screening-loss normal form. Hence a witness stays screening-only exactly as long as its upstream content is still realizable through some common carrier package with only packet/phase-lock/epsilon compatibility lost. Once that becomes impossible, the witness can no longer remain in screening-loss normal form alone and must instead force carrier/package loss. This proves the equivalence of (i), (ii), and (iii), and therefore the final escalation claim.
\end{proof}

\begin{corollary}[Escalation threshold for admitted visible witnesses]
\label{cor:carrier-escalation-threshold}
Within the current carrier-first OP~5 architecture, an admitted visible witness remains screening-only as long as its upstream content is representable by screening-loss normal form over some common closure-bearing return-side carrier package. It becomes carrier-upgraded precisely at the first point where such a carrier-supported realization ceases to exist.
\end{corollary}

\begin{proof}
This is the proof-usage reformulation of Lemma~\ref{lem:witness-to-carrier-escalation}. The existence of a common carrier package keeps the witness below the carrier threshold, while the first failure of every such carrier-supported realization marks the unique licensed escalation point.
\end{proof}


\subsection{Carrier-stable diagnostic collapse beneath the witness escalation threshold}
\label{sec:carrier-stable-diagnostic-collapse}

\begin{remark}[Collapse of visible diagnostic species over a stable carrier]
Once an admitted visible mismatch stays below the escalation threshold of Corollary~\ref{cor:carrier-escalation-threshold}, the current carrier-first OP~5 architecture no longer attributes independent structural rank to its packet-wave, breathing-envelope, frame-flicker, or helix presentation. Those visible realizations remain diagnostically useful, but over one and the same stable closure-bearing return-side carrier they collapse to a single subordinate screening-level phenomenon.
\end{remark}

\begin{theorem}[Carrier-stable diagnostic collapse theorem]
\label{thm:carrier-stable-diagnostic-collapse}
Assume Proposition~\ref{prop:carrier-first-op5-interface}, Proposition~\ref{prop:packet-phase-epsilon-screening}, Lemma~\ref{lem:diagnostic-readout-hierarchy}, Lemma~\ref{lem:first-diagnostic-witness}, Lemma~\ref{lem:screening-witness-compression}, Lemma~\ref{lem:carrier-witness-separation}, and Lemma~\ref{lem:witness-to-carrier-escalation}. Let $(b,r,s)$ be a branch-shell realization with an admitted visible mismatch $w$ of packet-wave, breathing-envelope, frame-flicker, or helix type. If $w$ remains below carrier escalation, equivalently if there exists a common closure-bearing return-side carrier package through which the upstream content of $w$ is still realized as screening-loss normal form, then the following collapse holds:
\begin{enumerate}[leftmargin=2em,label=(\roman*)]
\item \textbf{One subordinate structural role.} The visible realization of $w$ has exactly one current structural role, namely to witness subordinate screening-loss normal form beneath a still-stable carrier decision.
\item \textbf{No independent visible closure species.} Packet-wave, breathing-envelope, frame-flicker, and helix mismatch do not define distinct closure-relevant species as long as the common carrier package remains available.
\item \textbf{Carrier dominance.} Any closure-relevant distinction between those visible realizations collapses over the stable carrier and may reappear only if escalation to carrier/package loss occurs.
\item \textbf{Readout subordination.} The visible mismatch remains a downstream readout of one common carrier-supported screening loss and therefore cannot outrank the carrier decision itself.
\end{enumerate}
Hence, below the escalation threshold, all currently admitted visible diagnostic forms collapse to one carrier-stable subordinate diagnostic class.
\end{theorem}

\begin{proof}
By Lemma~\ref{lem:diagnostic-readout-hierarchy}, visible flicker/helix-type structure is subordinate to packet/phase/epsilon screening, which is itself subordinate to the carrier-first OP~5 decision. Lemma~\ref{lem:first-diagnostic-witness} allows visible mismatch to count upstream only when it certifies loss of a superior layer. Lemma~\ref{lem:screening-witness-compression} then shows that, below carrier loss, every admitted visible witness has exactly one upstream content: screening-loss normal form. Lemma~\ref{lem:carrier-witness-separation} and Corollary~\ref{cor:carrier-escalation-threshold} identify the present hypothesis of a still available common closure-bearing return-side carrier package as precisely the condition that the witness has not escalated to carrier/package loss. Therefore, while that common carrier package remains available, packet-wave, envelope, flicker, and helix mismatch cannot persist as distinct closure-relevant species. They all collapse to the single role of subordinate diagnostic realization of screening-loss normal form over one stable carrier. Proposition~\ref{prop:carrier-first-op5-interface} guarantees that any stronger distinction would already belong to carrier/package loss rather than to the pre-escalation regime. This proves (i)--(iv) and the final collapse claim.
\end{proof}

\begin{corollary}[Carrier-stable visible normal form]
\label{cor:carrier-stable-visible-normal-form}
Within the current carrier-first OP~5 architecture, every admitted visible mismatch that stays below carrier escalation may be cited simply as belonging to one \emph{carrier-stable visible normal form}. Its packet-wave, breathing-envelope, frame-flicker, or helix appearance remains diagnostically informative, but not structurally independent while the common carrier package persists.
\end{corollary}

\begin{proof}
This is the proof-usage reformulation of Theorem~\ref{thm:carrier-stable-diagnostic-collapse}. The theorem states exactly that all admitted visible mismatches below escalation collapse over one stable carrier to a single subordinate diagnostic class.
\end{proof}


\subsection{Bundled OP~5 diagnostic theorem}
\label{sec:bundled-op5-diagnostic-theorem}

\begin{remark}[One bundled positive/negative OP~5 diagnostic decision]
After the carrier-first OP~5 interface, the current diagnostic layer is now strong enough to be cited as one bundled decision principle. Either the visible material remains subordinate to a stable carrier and therefore collapses to a screening-level normal form, or it upgrades through an admitted witness to carrier/package loss. There is no third currently admitted diagnostic rank between those two bundled outcomes.
\end{remark}

\begin{theorem}[Bundled OP~5 diagnostic theorem]
\label{thm:bundled-op5-diagnostic-theorem}
Assume Proposition~\ref{prop:carrier-first-op5-interface}, Proposition~\ref{prop:packet-phase-epsilon-screening}, Lemma~\ref{lem:diagnostic-readout-hierarchy}, Lemma~\ref{lem:first-upstream-propagation}, Lemma~\ref{lem:first-diagnostic-witness}, Lemma~\ref{lem:screening-witness-compression}, Lemma~\ref{lem:carrier-witness-separation}, Lemma~\ref{lem:witness-to-carrier-escalation}, and Theorem~\ref{thm:carrier-stable-diagnostic-collapse}. Let $(b,r,s)$ be a routed branch-shell realization together with an admitted visible mismatch $w$ of packet-wave, breathing-envelope, frame-flicker, or helix type. Then exactly one of the following bundled OP~5 diagnostic outcomes holds:
\begin{enumerate}[leftmargin=2em,label=(\roman*)]
\item \textbf{Carrier-stable diagnostic collapse.} The mismatch $w$ remains below escalation, is carried by one common closure-bearing return-side carrier, and therefore collapses to a subordinate screening-level visible normal form beneath the carrier-first OP~5 decision.
\item \textbf{Carrier-upgraded diagnostic obstruction.} The mismatch $w$ cannot be realized over any common closure-bearing return-side carrier as mere screening-loss normal form and therefore upgrades to carrier/package loss. In that case $w$ is not a merely visible irregularity but an admitted carrier witness for routed OP~5 failure.
\end{enumerate}
Moreover, case \textup{(i)} excludes case \textup{(ii)}, and case \textup{(ii)} excludes case \textup{(i)}. Hence the current OP~5 diagnostic layer is completely bundled into the dichotomy
\[
\text{carrier-stable visible collapse}
\qquad\text{versus}\qquad
\text{carrier-upgraded diagnostic obstruction}.
\]
\end{theorem}

\begin{proof}
Lemma~\ref{lem:diagnostic-readout-hierarchy} and Lemma~\ref{lem:first-upstream-propagation} force every admitted visible mismatch to move upward only through the packet/phase/epsilon screening layer before it can acquire carrier relevance. Lemma~\ref{lem:first-diagnostic-witness} then identifies which visible mismatches may count upstream at all. Lemma~\ref{lem:screening-witness-compression} compresses every such admitted visible witness below carrier loss to one screening-loss normal form, while Lemma~\ref{lem:carrier-witness-separation} and Lemma~\ref{lem:witness-to-carrier-escalation} separate precisely the screening-only regime from the carrier-upgraded regime. Theorem~\ref{thm:carrier-stable-diagnostic-collapse} describes the entire positive side of the first regime: once a common closure-bearing return-side carrier remains available, all admitted visible species collapse to one subordinate visible normal form beneath the carrier-first OP~5 decision. Conversely, when no such common carrier remains available, Corollary~\ref{cor:carrier-witness-upgrade-criterion} upgrades the admitted visible witness to carrier/package loss. These two possibilities are mutually exclusive by Lemma~\ref{lem:carrier-witness-separation} and jointly exhaustive by the escalation dichotomy. This proves the bundled theorem.
\end{proof}

\begin{corollary}[Bundled routed OP~5 diagnostic use rule]
\label{cor:bundled-routed-op5-diagnostic-use-rule}
In later proof steps it is sufficient to cite Theorem~\ref{thm:bundled-op5-diagnostic-theorem} instead of separately unpacking packet-wave, breathing-envelope, frame-flicker, and helix diagnostics. One may simply ask whether the currently admitted visible data still collapse over a stable carrier or have already upgraded to carrier/package obstruction.
\end{corollary}

\begin{proof}
This is the proof-usage reformulation of Theorem~\ref{thm:bundled-op5-diagnostic-theorem}. The theorem already packages the visible diagnostic layer into exactly the two current routed OP~5 outcomes.
\end{proof}


\subsection{First OP~3 to OP~5 bundled transfer theorem}
\label{sec:first-op3-op5-bundled-transfer}

\begin{remark}[Transfer of bundled OP~3 closure into the OP~5 diagnostic layer]
The current architecture no longer treats OP~3 closure and OP~5 diagnosis as two independent theorem blocks. Once the OP~3 shell passage has been certified and rerouted through one common closure-bearing carrier, the routed OP~5 side inherits a downstream decision of that same carrier package. The only remaining terminal question is whether the visible return side still collapses beneath that package or whether it upgrades to a carrier-level obstruction. Thus OP~5 is presently read as the bundled downstream transfer image of the OP~3/QAM closure package.
\end{remark}

\begin{theorem}[First OP~3 to OP~5 bundled transfer theorem]
\label{thm:first-op3-op5-bundled-transfer}
Assume Theorem~\ref{thm:qam-op3-shell-closure}, Theorem~\ref{thm:qam-op3-proof-run}, Theorem~\ref{thm:combined-branch-shell}, Proposition~\ref{prop:carrier-first-op5-interface}, and Theorem~\ref{thm:bundled-op5-diagnostic-theorem}. Let $(b,r,s)$ be a routed branch-shell realization whose OP~3 shell passage is translated into the extracted QAM layer, and let $w$ be an admitted visible mismatch on the routed OP~5 return side. Then the present bundled OP~3$\to$OP~5 transfer has exactly two outcomes:
\begin{enumerate}[leftmargin=2em,label=(\roman*)]
\item \textbf{Transferred carrier-stable pass.}
The translated OP~3 shell passage belongs to the certified side of the bundled OP~3 closure package, the routed return stays inside the same common closure-bearing carrier, and the admitted visible mismatch $w$ collapses beneath that carrier to the subordinate OP~5 visible normal form of Theorem~\ref{thm:bundled-op5-diagnostic-theorem}.
\item \textbf{Transferred localized obstruction.}
The transfer fails at the first licensed obstruction layer: either the OP~3/QAM side already reaches the first localized shell obstruction of the extracted OP~3 proof run, or the routed return upgrades to a carrier-level OP~5 diagnostic obstruction. In particular, no third autonomous terminal OP~5 obstruction is presently licensed downstream of an otherwise intact bundled OP~3 carrier package.
\end{enumerate}
Hence the current routed OP~5 diagnostic layer is the bundled downstream transfer of the OP~3/QAM closure package, not an additional independent closure theory.
\end{theorem}

\begin{proof}
Theorem~\ref{thm:qam-op3-proof-run} already packages the translated OP~3 shell passage into exactly two upstream outcomes: certified passage through the bundled OP~3 closure tool layer or the first localized obstruction. If the certified side holds, then Theorem~\ref{thm:qam-op3-shell-closure} identifies one common closure-bearing carrier/readout package for the shell passage. Theorem~\ref{thm:combined-branch-shell} transports that same certified branch-shell realization to the routed return side without changing the invariant closure class, and Proposition~\ref{prop:carrier-first-op5-interface} forces the terminal OP~5 reading to be carrier-first rather than a new autonomous endpoint test. Theorem~\ref{thm:bundled-op5-diagnostic-theorem} then completes the positive branch: as long as the routed return remains over that same common carrier, every admitted visible mismatch collapses to the subordinate carrier-stable OP~5 visible normal form. This proves case~(i).

If the transfer does not remain on that certified branch, then the failure must appear at the first licensed obstruction layer. Upstream, Theorem~\ref{thm:qam-op3-proof-run} and the extracted OP~3/QAM obstruction package localize the first obstruction already at branch, shell, or carrier/readout level. Downstream, if the routed return cannot be realized over any common closure-bearing carrier, then Theorem~\ref{thm:bundled-op5-diagnostic-theorem} upgrades the admitted visible mismatch to a carrier-level OP~5 diagnostic obstruction. Because Proposition~\ref{prop:carrier-first-op5-interface} forbids an independent terminal closure test downstream of an otherwise intact carrier package, there is no currently licensed third obstruction type appearing only at the raw terminal visible layer. This proves case~(ii) and the final bundled transfer reading.
\end{proof}

\begin{corollary}[Bundled OP~3$\to$OP~5 proof-use rule]
\label{cor:bundled-op3-op5-proof-use-rule}
In later proof steps it is sufficient to cite Theorem~\ref{thm:first-op3-op5-bundled-transfer} when a routed branch-shell realization must be carried from the extracted OP~3/QAM closure package into the OP~5 diagnostic layer. One may simply ask whether the realization stays on the transferred carrier-stable branch or reaches the first licensed transferred obstruction.
\end{corollary}

\begin{proof}
This is the proof-usage reformulation of Theorem~\ref{thm:first-op3-op5-bundled-transfer}. The theorem already packages the current OP~3/QAM closure side and the current OP~5 diagnostic side into one bundled transfer decision.
\end{proof}


\subsection{First OP~3 to OP~5 compressed proof interface}
\label{sec:first-op3-op5-compressed-proof-interface}

\begin{remark}[How older routed return arguments should now be compressed]
The current transfer architecture is now strong enough that later routed return arguments need not separately reopen the combined branch-shell theorem, the carrier-first OP~5 interface, and the bundled OP~5 diagnostic layer. Once the routed realization has entered the bundled OP~3$\to$OP~5 transfer channel, the proof may proceed entirely through one compressed dichotomy: transferred carrier-stable pass versus transferred localized obstruction. This local compressed interface is now itself subordinate to the newer global rule: as soon as the argument has also entered the bundled current channel and no shell-sensitive or witness-grade claim is being asked, the proof should move on to the extracted global tools and remain global-first unless a genuine local reopening is needed.
\end{remark}

\begin{proposition}[First OP~3 to OP~5 compressed proof interface]
\label{prop:first-op3-op5-compressed-proof-interface}
Assume Theorem~\ref{thm:first-op3-op5-bundled-transfer}. Let $(b,r,s)$ be a routed branch-shell realization whose OP~3 shell passage has been translated into the extracted QAM layer, and let $w$ be any admitted visible routed return witness on the OP~5 side. Then every later proof step that needs to carry this realization through the current routed return architecture may be compressed to the following interface.
\begin{enumerate}[leftmargin=2em,label=(\roman*)]
\item \textbf{Transferred carrier-stable pass.} If Theorem~\ref{thm:first-op3-op5-bundled-transfer} places $(b,r,s;w)$ on the transferred carrier-stable branch, then all downstream OP~5 packet/phase/flicker/helix material is read only as subordinate visible realization beneath one common closure-bearing carrier and may be replaced by the carrier-stable visible normal form.
\item \textbf{Transferred localized obstruction.} If Theorem~\ref{thm:first-op3-op5-bundled-transfer} places $(b,r,s;w)$ on the transferred obstruction branch, then the proof may stop at the first licensed obstruction already identified there: either the upstream OP~3/QAM obstruction or the downstream carrier-upgraded OP~5 obstruction.
\item \textbf{No third terminal proof branch.} No additional raw terminal visible branch is presently licensed after this compression. In particular, later proofs do not need a separate fourth case for packet-wave, breathing-envelope, frame-flicker, or helix mismatch once the bundled transfer theorem has already been invoked.
\end{enumerate}
Hence the current routed return architecture may be used through one compressed OP~3$\to$OP~5 proof interface.
\end{proposition}

\begin{proof}
This is the direct proof-usage reading of Theorem~\ref{thm:first-op3-op5-bundled-transfer}. The theorem already packages the routed return into exactly two licensed transfer outcomes. On the positive side, the admitted visible material collapses beneath one common carrier and therefore needs no further structural branching. On the negative side, the first licensed obstruction is already localized either upstream in the OP~3/QAM package or downstream at the carrier-upgraded OP~5 level. Because Theorem~\ref{thm:first-op3-op5-bundled-transfer} excludes an autonomous raw terminal OP~5 obstruction beneath an otherwise intact carrier package, no third terminal proof branch remains.
\end{proof}

\begin{corollary}[Compressed rereading of routed return arguments through the bundled transfer]
\label{cor:compressed-routed-return-through-bundled-transfer}
In later arguments it is sufficient to cite Proposition~\ref{prop:first-op3-op5-compressed-proof-interface} rather than separately reopening Theorem~\ref{thm:combined-branch-shell}, Proposition~\ref{prop:carrier-first-op5-interface}, and Theorem~\ref{thm:bundled-op5-diagnostic-theorem}. One may simply ask whether the current routed realization stays on the transferred carrier-stable branch or reaches the first licensed transferred obstruction.
\end{corollary}

\begin{proof}
This is the condensed proof-usage form of Proposition~\ref{prop:first-op3-op5-compressed-proof-interface}. The proposition already compresses the presently licensed routed return architecture into one reusable interface.
\end{proof}


\subsection{First routed return normal form theorem}
\label{sec:first-routed-return-normal-form-theorem}

\begin{remark}[Carrier-stable versus obstructed routed return normal form]
After the compressed OP~3$\to$OP~5 proof interface has been invoked, the current routed return architecture no longer needs a separate case distinction for packet-wave, breathing-envelope, frame-flicker, or helix appearance. At the present level of the theory every admitted routed return is read only up to one of two bundled normal forms: a carrier-stable routed return normal form or an obstructed routed return normal form.
\end{remark}

\begin{theorem}[First routed return normal form theorem]
\label{thm:first-routed-return-normal-form}
Assume Proposition~\ref{prop:first-op3-op5-compressed-proof-interface}. Let $(b,r,s)$ be a routed branch-shell realization whose OP~3 shell passage has been translated into the extracted QAM layer, and let $w$ be any admitted visible routed return witness on the OP~5 side. Then the current routed return architecture classifies $(b,r,s;w)$, up to bundled transfer, into exactly one of the following two mutually exclusive routed return normal forms.
\begin{enumerate}[leftmargin=2em,label=(\roman*)]
\item \textbf{Carrier-stable routed return normal form.}
The routed return remains on the transferred carrier-stable branch of Proposition~\ref{prop:first-op3-op5-compressed-proof-interface}. Equivalently, there exists one common closure-bearing return-side carrier through which the routed realization is read, and every admitted visible mismatch collapses beneath that carrier to the carrier-stable visible normal form. In this case the routed return contributes no additional structural branching beyond the bundled OP~3$\to$OP~5 transfer itself.
\item \textbf{Obstructed routed return normal form.}
The routed return lies on the transferred localized-obstruction branch of Proposition~\ref{prop:first-op3-op5-compressed-proof-interface}. Equivalently, no common closure-bearing routed return carrier supports the realization as mere screening-level visible loss, and the first licensed failure is already localized either upstream at the OP~3/QAM obstruction layer or downstream at the carrier-upgraded OP~5 obstruction layer.
\end{enumerate}
Hence every currently admitted routed return is read, up to the bundled OP~3$\to$OP~5 transfer, by one bundled routed return normal form and not by any third autonomous terminal visible class.
\end{theorem}

\begin{proof}
Proposition~\ref{prop:first-op3-op5-compressed-proof-interface} already compresses the present routed return architecture into exactly two licensed proof branches: transferred carrier-stable pass and transferred localized obstruction. On the positive branch, Theorem~\ref{thm:bundled-op5-diagnostic-theorem} and the carrier-stable collapse package imply that every admitted visible return witness becomes subordinate to one common closure-bearing routed return carrier and therefore yields the carrier-stable routed return normal form. On the negative branch, the same compressed interface localizes the first licensed failure either upstream in the extracted OP~3/QAM obstruction package or downstream at the carrier-upgraded OP~5 obstruction layer; by construction this is the obstructed routed return normal form. Because Proposition~\ref{prop:first-op3-op5-compressed-proof-interface} excludes a third raw terminal visible branch after compression, these two forms are exhaustive and mutually exclusive.
\end{proof}

\begin{corollary}[Routed return normal form use rule]
\label{cor:routed-return-normal-form-use-rule}
In later proofs it is sufficient to cite Theorem~\ref{thm:first-routed-return-normal-form} once a routed realization has already entered the bundled OP~3$\to$OP~5 transfer channel. One may then work only with the dichotomy \emph{carrier-stable routed return normal form versus obstructed routed return normal form}.
\end{corollary}

\begin{proof}
This is the proof-usage reformulation of Theorem~\ref{thm:first-routed-return-normal-form}. The theorem already packages the presently licensed routed return architecture into one reusable normal-form dichotomy.
\end{proof}


\subsection{First global pass/obstruction normal form}
\label{sec:first-global-pass-obstruction-normal-form}

\begin{remark}[Current global normal form of the admissible run]
The present line now carries enough bundled structure that the currently admitted run
\[
\begin{aligned}
&\text{OP1 source}
\;\to\;\text{branch admissibility}\\[0.15em]
&\;\to\;\text{OP3/QAM closure package}\\[0.15em]
&\;\to\;\text{carrier-first OP5 layer}
\end{aligned}
\]
no longer needs to be read as a chain of loosely connected local tests. After the routed
return normal form has been established, the full currently admitted run is read only up to
one global dichotomy: a global pass normal form or a global obstruction normal form.
\end{remark}

\begin{theorem}[First global pass/obstruction normal form]
\label{thm:first-global-pass-obstruction-normal-form}
Assume Theorem~\ref{thm:first-routed-return-normal-form}. Let $(b,r,s)$ be an admissible
branch-shell realization issued from the current OP1 source side and translated through the
extracted OP3/QAM layer into the routed return architecture, and let $w$ be any admitted
visible routed return witness on the OP5 side. Then the whole currently admitted run from OP1
through branch, OP3/QAM, and OP5 is classified, up to the present bundled architecture, into
exactly one of the following two mutually exclusive global normal forms.
\begin{enumerate}[leftmargin=2em,label=(\roman*)]
\item \textbf{Global pass normal form.}
The realization lies on the carrier-stable routed return normal form of
Theorem~\ref{thm:first-routed-return-normal-form}. Equivalently, the run remains supported by
one bundled closure-bearing carrier package from the present OP1/branch entrance through the
OP3/QAM shell package and into the routed return side, while all visible OP5 material collapses
beneath that carrier to the carrier-stable visible normal form. In this case the current run
carries no further licensed structural branch beyond the bundled pass reading.
\item \textbf{Global obstruction normal form.}
The realization lies on the obstructed routed return normal form of
Theorem~\ref{thm:first-routed-return-normal-form}. Equivalently, the first licensed failure of
the current run is already localized either upstream in the OP3/QAM obstruction package or
downstream at the carrier-upgraded OP5 obstruction layer, and no autonomous third terminal
visible class is currently licensed outside that bundled obstruction reading.
\end{enumerate}
Hence the whole currently admitted run is read, at the present stage of the theory, by one
global pass/obstruction normal form and not by any third independent terminal morphology.
\end{theorem}

\begin{proof}
Theorem~\ref{thm:first-routed-return-normal-form} already classifies every currently admitted
routed return into exactly two bundled normal forms. On the positive side, the carrier-stable
routed return normal form means that the run has remained on the transferred carrier-stable
branch from the OP3/QAM closure package into the routed return side; consequently the whole
present run remains supported by one bundled closure-bearing carrier package, and all visible
OP5 diagnostics are subordinate to the carrier-stable visible normal form. This is precisely the
global pass normal form. On the negative side, the obstructed routed return normal form already
states that the first licensed failure is localized either upstream in the OP3/QAM obstruction
package or downstream at the carrier-upgraded OP5 obstruction layer; this is exactly the global
obstruction normal form. Because Theorem~\ref{thm:first-routed-return-normal-form} excludes a
third autonomous terminal visible class, the two global forms are exhaustive and mutually
exclusive.
\end{proof}

\begin{corollary}[Global pass/obstruction use rule]
\label{cor:global-pass-obstruction-use-rule}
In later proofs it is sufficient to cite Theorem~\ref{thm:first-global-pass-obstruction-normal-form}
once the current admissible run has entered the bundled OP3/QAM-to-OP5 channel. One may then
work only with the dichotomy \emph{global pass normal form versus global obstruction normal
form}.
\end{corollary}

\begin{proof}
This is the proof-usage reformulation of
Theorem~\ref{thm:first-global-pass-obstruction-normal-form}. The theorem already compresses the
currently admitted run from OP1 through OP5 into one reusable bundled normal-form dichotomy.
\end{proof}


\begin{remark}[Status of the bundled OP~3$\to$OP~5/global chain]
\label{rem:status-bundled-op3-op5-global-chain}
The present local-to-global chain should now also be read in two levels of proof strength. The \emph{active bundled interface} consists of the first OP~3$\to$OP~5 bundled transfer theorem, the compressed OP~3$\to$OP~5 proof interface, the routed return normal form, and the global pass/obstruction normal form together with their direct proof-use corollaries. These statements already give one self-contained current channel from translated OP~3 certification to a bundled global pass/obstruction decision. By contrast, the finer witness-grade hierarchy beneath the carrier-first OP~5 layer --- screening-loss normal form, carrier-witness upgrade, escalation thresholds, and the related witness-compression remarks --- remains a \emph{local refinement layer}: it must be reopened only when a later argument genuinely depends on witness species, witness grade, escalation, or screening-versus-carrier distinctions that are invisible at the bundled transfer/global level.
\end{remark}

\begin{proposition}[Citation and reopening rule for the bundled OP~3$\to$OP~5/global chain]
\label{prop:citation-reopening-rule-bundled-op3-op5-global}
Let a later argument concern a run already translated from the OP~3/QAM shell package into the routed OP~5 side. Then the following citation discipline is currently licensed.
\begin{enumerate}[leftmargin=2em,label=(\roman*)]
  \item If the argument needs only the bundled transfer/global dichotomy, it may cite Theorem~\ref{thm:first-op3-op5-bundled-transfer}, Proposition~\ref{prop:first-op3-op5-compressed-proof-interface}, Theorem~\ref{thm:first-routed-return-normal-form}, and Theorem~\ref{thm:first-global-pass-obstruction-normal-form} without reopening the finer witness hierarchy.
  \item If the argument further needs shell-level certificate or shell-level obstruction data, it must reopen the completed OP~3/QAM block rather than the entire older OP~1--OP~5 line.
  \item If the argument further needs witness grade, witness escalation, screening-witness compression, or carrier-witness separation, it must reopen the carrier-first OP~5 witness layer rather than treating the bundled global dichotomy as sufficient.
  \item No later proof may insert a third autonomous raw terminal visible branch after the bundled transfer/global interface has already been invoked.
\end{enumerate}
\end{proposition}

\begin{proof}
Clause~(i) is exactly the compressed reading already supplied by Theorem~\ref{thm:first-op3-op5-bundled-transfer}, Proposition~\ref{prop:first-op3-op5-compressed-proof-interface}, Theorem~\ref{thm:first-routed-return-normal-form}, and Theorem~\ref{thm:first-global-pass-obstruction-normal-form}. Clause~(ii) is the existing shell-sensitive fallback discipline of the completed OP~3/QAM block. Clause~(iii) is the corresponding witness-sensitive fallback discipline of the carrier-first OP~5 layer, where screening-loss normal form, upgrade, escalation, and separation are defined. Clause~(iv) is the bundled exclusion already built into the transfer theorem, the compressed proof interface, the routed return normal form, and the global pass/obstruction theorem: once those interfaces are active, no additional autonomous raw terminal visible branch is currently licensed.
\end{proof}

\begin{corollary}[Current sufficiency of the bundled OP~3$\to$OP~5/global chain]
\label{cor:current-sufficiency-bundled-op3-op5-global}
Assume a current argument starts from a translated OP~3 shell passage and asks only whether the run reaches a bundled pass reading or a bundled first licensed obstruction. Then the present document already contains a self-contained proof interface for that question: one may pass from the completed OP~3/QAM block through the bundled OP~3$\to$OP~5 transfer, the routed return normal form, and finally the global pass/obstruction normal form without introducing any additional intermediate theorem layer. In particular, the current line is already strong enough for the next public-shape step whenever no stronger claim about witness grade or deeper kernel-first activation is being asserted.
\end{corollary}

\begin{proof}
By Theorem~\ref{thm:qam-op3-proof-run} and Corollary~\ref{cor:qam-op3-proof-run-use}, the completed OP~3/QAM block already reduces translated shell data to certified passage versus first localized obstruction. Theorem~\ref{thm:first-op3-op5-bundled-transfer} and Proposition~\ref{prop:first-op3-op5-compressed-proof-interface} then transport that same dichotomy into the routed OP~5 side. Theorem~\ref{thm:first-routed-return-normal-form} rewrites the routed return as carrier-stable versus obstructed normal form, and Theorem~\ref{thm:first-global-pass-obstruction-normal-form} finally packages the whole current run into global pass versus global obstruction. Thus no extra intermediate theorem layer is needed for that bundled question. The final sentence is only the proof-theoretic reading of Remark~\ref{rem:status-bundled-op3-op5-global-chain}: stronger witness-grade or deeper kernel-first claims still require reopening the finer local layers.
\end{proof}


\subsection{Current extracted global pass/obstruction tool layer}
\label{sec:current-extracted-global-pass-obstruction-tool-layer}

\begin{remark}[Purpose of the current global tool layer]
The global pass/obstruction normal form is not meant to remain only as a terminal summary theorem. It is now extracted as a reusable tool layer so that later proofs can cite a small global interface instead of reopening the whole chain
\[
\text{OP1 source}
\to
\text{branch admissibility}
\to
\text{OP3/QAM closure}
\to
\text{carrier-first OP5 layer}.
\]
The extracted tools below are therefore not new theory beyond Theorem~\ref{thm:first-global-pass-obstruction-normal-form}; they are proof-usable projections of that bundled global dichotomy.
\end{remark}

\begin{proposition}[Global-pass tool]
\label{prop:global-pass-tool}
Assume Theorem~\ref{thm:first-global-pass-obstruction-normal-form}. If a current admissible run lies on the global pass normal form, then one may cite this fact directly as a reusable proof tool: the run is carried by one bundled closure-bearing carrier package from the OP1/branch entrance through the OP3/QAM shell package into the routed return side, and every admitted OP5 visible diagnostic remains subordinate to the carrier-stable visible normal form.
\end{proposition}

\begin{proof}
This is precisely the positive branch of Theorem~\ref{thm:first-global-pass-obstruction-normal-form}, rewritten as a direct reusable proof tool.
\end{proof}

\begin{proposition}[Global-obstruction tool]
\label{prop:global-obstruction-tool}
Assume Theorem~\ref{thm:first-global-pass-obstruction-normal-form}. If a current admissible run lies on the global obstruction normal form, then one may cite this fact directly as a reusable proof tool: the first licensed failure is already localized either upstream in the OP3/QAM obstruction package or downstream at the carrier-upgraded OP5 obstruction layer, and no third autonomous terminal visible class is currently licensed outside that bundled obstruction reading.
\end{proposition}

\begin{proof}
This is precisely the negative branch of Theorem~\ref{thm:first-global-pass-obstruction-normal-form}, rewritten as a direct reusable proof tool.
\end{proof}

\begin{proposition}[Global transfer tool]
\label{prop:global-transfer-tool}
Assume Theorem~\ref{thm:first-global-pass-obstruction-normal-form}. Any later proof that has already routed the current admissible run into the bundled OP3/QAM-to-OP5 channel may work only with the bundled transfer dichotomy
\[
\text{global pass}
\qquad\text{versus}\qquad
\text{global obstruction},
\]
without reopening the intermediate local decomposition into separate branch, shell, witness, escalation, and routed-return subcases.
\end{proposition}

\begin{proof}
By Corollary~\ref{cor:global-pass-obstruction-use-rule}, the whole currently admitted run has already been compressed into one reusable global dichotomy. The omitted intermediate subcases are already absorbed into the proof of Theorem~\ref{thm:first-global-pass-obstruction-normal-form}.
\end{proof}

\begin{proposition}[Global diagnostic subordination tool]
\label{prop:global-diagnostic-subordination-tool}
Assume Theorem~\ref{thm:first-global-pass-obstruction-normal-form}. On the global pass branch, every currently admitted visible routed-return diagnostic is globally subordinate to the same bundled closure-bearing carrier package. In particular, packet-wave, breathing-envelope, frame-flicker, and helix pictures cannot reintroduce a new structural branch once the run has already entered the global pass normal form.
\end{proposition}

\begin{proof}
Theorem~\ref{thm:first-global-pass-obstruction-normal-form} sends the positive branch through the carrier-stable routed return normal form and therefore through the carrier-stable visible normal form. By construction, all visible routed-return diagnostics are subordinate to that same bundled carrier package and cannot reopen the global dichotomy.
\end{proof}

\begin{proposition}[Global proof-interface tool]
\label{prop:global-proof-interface-tool}
Assume Theorem~\ref{thm:first-global-pass-obstruction-normal-form}. In later arguments it is sufficient to cite one of the four extracted global tools
\[
\text{global pass},\quad
\text{global obstruction},\quad
\text{global transfer},\quad
\text{global diagnostic subordination},
\]
according to the local need of the argument. No additional independent proof interface for the currently admitted run is required at this stage.
\end{proposition}

\begin{proof}
Theorem~\ref{thm:first-global-pass-obstruction-normal-form} already supplies the full current bundled dichotomy, while Propositions~\ref{prop:global-pass-tool}--\ref{prop:global-diagnostic-subordination-tool} expose its presently relevant projections. Hence these four extracted global tools form the current reusable proof interface for the admissible run.
\end{proof}

\begin{remark}[How to cite the extracted global tools]
The current global tool layer should be cited only after the run has already been shown to lie inside the bundled OP1$\to$branch$\to$OP3/QAM$\to$OP5 channel. Once that has been established, the present extracted tools allow later arguments to cite a compact global interface rather than reopening the entire local transport and diagnostic stack.
\end{remark}


\subsection{First global compressed proof run through extracted global tools}
\label{sec:first-global-compressed-proof-run-through-extracted-global-tools}

\begin{remark}[Reading key for the first global compressed proof run]
The present block is the first place where the current admissible run
\[
\text{OP1 source}
\to
\text{branch admissibility}
\to
\text{OP3/QAM closure}
\to
\text{carrier-first OP5 layer}
\]
is no longer reopened stage by stage. Instead, once the run has already been shown to lie inside the bundled global channel, the extracted global tools are used as one compressed proof interface. The proof run therefore has only two licensed bundled outcomes: a \emph{global pass} and a \emph{global obstruction}. The purpose of the theorem below is to state this compressed use rule explicitly, so that later arguments may pass through the whole current architecture in one step.
\end{remark}

\begin{theorem}[First global compressed proof run through extracted global tools]
\label{thm:first-global-compressed-proof-run-through-extracted-global-tools}
Assume that a current admissible run issued from the OP1 source side has already been routed into the bundled OP1$\to$branch$\to$OP3/QAM$\to$OP5 channel of the present architecture. Then the run may be processed through the extracted global tool layer without reopening the intermediate local stack. More precisely, exactly one of the following two bundled proof outputs is currently licensed.
\begin{enumerate}[leftmargin=2em,label=(\roman*)]
\item \textbf{Global compressed pass.}
By Proposition~\ref{prop:global-pass-tool}, the run lies on the global pass normal form and is supported by one bundled closure-bearing carrier package from the OP1/branch entrance through the OP3/QAM shell package into the routed return side. By Proposition~\ref{prop:global-diagnostic-subordination-tool}, all admitted visible routed-return diagnostics remain subordinate to that same bundled carrier package.
\item \textbf{Global compressed obstruction.}
By Proposition~\ref{prop:global-obstruction-tool}, the run lies on the global obstruction normal form and its first licensed failure is already localized either upstream in the OP3/QAM obstruction package or downstream at the carrier-upgraded OP5 obstruction layer. By Proposition~\ref{prop:global-transfer-tool}, no third autonomous terminal visible branch is currently licensed outside that bundled obstruction reading.
\end{enumerate}
Hence the present global proof run is now compressed to the dichotomy
\[
\text{global compressed pass}
\qquad\text{versus}\qquad
\text{global compressed obstruction},
\]
and no further independent proof interface for the currently admitted run is required once Proposition~\ref{prop:global-proof-interface-tool} applies.
\end{theorem}

\begin{proof}
This is the first direct use of the extracted global tool layer as a compressed proof interface. Once the current run has entered the bundled
\[
\text{OP1}\to\text{branch}\to\text{OP3/QAM}\to\text{OP5}
\]
channel, Proposition~\ref{prop:global-proof-interface-tool} states that no additional independent proof interface is currently required.

If the run lies on the positive side of Theorem~\ref{thm:first-global-pass-obstruction-normal-form}, Proposition~\ref{prop:global-pass-tool} supplies the bundled closure-bearing carrier package and Proposition~\ref{prop:global-diagnostic-subordination-tool} shows that all admitted visible OP5 diagnostics remain subordinate to that same carrier; this is exactly the global compressed pass.

If the run lies on the negative side, Proposition~\ref{prop:global-obstruction-tool} localizes the first licensed failure and Proposition~\ref{prop:global-transfer-tool} excludes any third independent terminal visible branch; this is exactly the global compressed obstruction.

Exhaustiveness and mutual exclusivity therefore come directly from the extracted global tools and ultimately from Theorem~\ref{thm:first-global-pass-obstruction-normal-form}.
\end{proof}

\begin{corollary}[First operational use of the extracted global tools]
\label{cor:first-operational-use-of-the-extracted-global-tools}
In later arguments it is sufficient to cite Theorem~\ref{thm:first-global-compressed-proof-run-through-extracted-global-tools} once the current admissible run has entered the bundled global channel. One may then work only with the dichotomy \emph{global compressed pass versus global compressed obstruction}, together with the induced carrier-stable subordination or obstruction localization, without reopening the entire local OP3/QAM-to-OP5 stack.
\end{corollary}

\begin{proof}
This is the operational reformulation of Theorem~\ref{thm:first-global-compressed-proof-run-through-extracted-global-tools}. The theorem already packages the extracted global tools into one reusable compressed proof run for the currently admitted architecture.
\end{proof}


\subsection{Global proof compression pass for the older OP1--OP5 line}
\label{sec:global-proof-compression-pass-for-the-older-op1-op5-line}

\begin{remark}[How the older OP1--OP5 line should now be reread]
The older connected OP1--OP5 transport line should no longer be read as a sequence of largely separate source-return, branch, shell, routed, and diagnostic arguments that must each be reopened in full. After the present global compressed proof run, the mathematically correct reading is instead this: once an argument has entered the currently admitted bundled channel, the older line should be read global-first through the dichotomy
\[
\text{global compressed pass}
\qquad\text{versus}\qquad
\text{global compressed obstruction}.
\]
Only if a later step genuinely depends on shell certificates, shell-obstruction resolution, witness grades, witness escalation, or other local carrier-diagnostic data should one reopen the OP3/QAM or OP5 layers. The purpose of the proposition below is not to replace the older theorems, but to state explicitly that they may now be reused through one compressed global interface.
\end{remark}

\begin{proposition}[Compressed rereading of the older OP1--OP5 line through the extracted global tools]
\label{prop:compressed-rereading-of-the-older-op1-op5-line-through-the-extracted-global-tools}
Assume that an argument along the older OP1--OP5 transport line has already been shown to lie inside the bundled channel of Theorem~\ref{thm:first-global-compressed-proof-run-through-extracted-global-tools}. Then the older line may be reread through the extracted global tools in the following compressed form.
\begin{enumerate}[leftmargin=2em,label=(\roman*)]
\item \textbf{Compressed pass rereading.}
If the argument lies on the positive side of Theorem~\ref{thm:first-global-compressed-proof-run-through-extracted-global-tools}, then the older OP1--OP5 line is to be read as one globally compressed pass: the source-return closure theorem, the admissible branch layer, the OP3/QAM closure package, the carrier-first OP5 interface, and the subordinated diagnostic stack are all already absorbed into one bundled carrier-stable run.
\item \textbf{Compressed obstruction rereading.}
If the argument lies on the negative side of Theorem~\ref{thm:first-global-compressed-proof-run-through-extracted-global-tools}, then the older OP1--OP5 line is to be read as one globally compressed obstruction: the first licensed failure is already localized in the bundled obstruction package, and no extra autonomous terminal visible branch may be added afterwards.
\end{enumerate}
In particular, later reuse of the older OP1--OP5 text should proceed through the extracted global proof interface rather than by reopening the whole local stack unless a genuinely new layer is being introduced.
\end{proposition}

\begin{proof}
This is the direct compression reading supplied by Theorem~\ref{thm:first-global-compressed-proof-run-through-extracted-global-tools}. On the positive side, the theorem already packages the full currently admitted run into the global compressed pass, and the extracted global tools show that the visible OP5 diagnostics remain subordinate to the same bundled carrier package. On the negative side, the same theorem packages the run into the global compressed obstruction and excludes any third autonomous terminal visible branch. Hence the older OP1--OP5 line can now be reused through one global proof interface rather than by reopening each intermediate theorem separately.
\end{proof}

\begin{corollary}[Global compressed reading of older routed return arguments]
\label{cor:global-compressed-reading-of-older-routed-return-arguments}
Whenever an older routed return argument is already known to lie inside the bundled current channel, it should by default be cited through Proposition~\ref{prop:compressed-rereading-of-the-older-op1-op5-line-through-the-extracted-global-tools}, Theorem~\ref{thm:first-global-compressed-proof-run-through-extracted-global-tools}, and the global citation discipline of Proposition~\ref{prop:first-global-tool-citation-guide}. In that situation, it is sufficient to work only with the dichotomy
\[
\text{global compressed pass}
\qquad\text{versus}\qquad
\text{global compressed obstruction},
\]
and with the induced bundled carrier-stable or obstruction-localized reading, unless a genuinely shell-sensitive or witness-sensitive local claim must still be reopened.
\end{corollary}

\begin{proof}
This is the operational reformulation of Proposition~\ref{prop:compressed-rereading-of-the-older-op1-op5-line-through-the-extracted-global-tools}. Once the argument is inside the bundled current channel, the extracted global proof interface supplies the compressed rereading rule for the older routed return line.
\end{proof}


\subsection{Witness-grade hardening and the first local kernel-first rewrite window}
\label{sec:witness-grade-hardening-and-first-local-kernel-first-rewrite-window}

\begin{remark}[Purpose of the v3.19.0 witness-grade hardening step]
The public v3.18.0 release already exposed the bundled OP~3/QAM closure package and the bundled OP~3$\to$OP~5/global channel in compact public shape. The next task is therefore not a new visible closure notion, but a stricter proof-use discipline for that same channel. In particular, the present step isolates three things: first, a bundled witness-grade reopening package for the OP~5 side; second, a sharper local/global reopening rule explaining exactly when the bundled global interface is sufficient; and third, one first narrow kernel-first rewrite window whose use is licensed only after the bundled channel has already been fixed and only at a local witness-grade transition where the deferred bridge actually reduces proof load.
\end{remark}

\begin{theorem}[Bundled witness-grade reopening package]
\label{thm:bundled-witness-grade-reopening-package}
Assume Theorem~\ref{thm:bundled-op5-diagnostic-theorem}, Proposition~\ref{prop:first-op3-op5-compressed-proof-interface}, Theorem~\ref{thm:first-global-pass-obstruction-normal-form}, Proposition~\ref{prop:first-global-tool-citation-guide}, and Proposition~\ref{prop:first-global-usage-checklist}. Let a later argument concern a run already known to lie inside the bundled current channel. Then exactly one of the following two proof-use regimes is presently licensed.
\begin{enumerate}[leftmargin=2em,label=(\roman*)]
\item \textbf{Bundled global-use regime.} If the argument needs only the bundled outputs global compressed pass/global compressed obstruction together with their carrier-stable or obstruction-localized consequences, then the extracted global interface is sufficient and no local witness-grade reopening is licensed or needed.
\item \textbf{Witness-grade reopening regime.} If the argument needs the witness grade itself, witness escalation, screening-witness compression, carrier-witness separation, diagnostic collapse below the escalation threshold, or a comparison between shell-grade and witness-grade data, then the bundled global interface must be reopened locally at the OP~5 witness layer, and the proof must work with the finer witness hierarchy rather than with the global dichotomy alone.
\end{enumerate}
Hence the current architecture already supports one bundled witness-grade reopening package: bundled-global use by default, local witness-grade reopening only for genuinely witness-sensitive claims.
\end{theorem}

\begin{proof}
The bundled global-use regime is exactly the content of Proposition~\ref{prop:first-global-tool-citation-guide} together with Proposition~\ref{prop:first-global-usage-checklist}: once the argument is already inside the bundled current channel and only the global pass/obstruction outputs are needed, the extracted global layer is the default proof interface. By contrast, Theorem~\ref{thm:bundled-op5-diagnostic-theorem} and Proposition~\ref{prop:first-op3-op5-compressed-proof-interface} show that witness grade, witness escalation, screening-witness compression, carrier-witness separation, and diagnostic collapse all live strictly below that bundled global interface. Therefore any genuinely witness-sensitive claim forces a local reopening of the OP~5 witness layer. The two regimes are exhaustive because every later proof step either remains at bundled-global level or asks for one of the finer witness-sensitive distinctions just listed.
\end{proof}

\begin{proposition}[Witness-grade local/global reopening discipline]
\label{prop:witness-grade-local-global-reopening-discipline}
Assume Theorem~\ref{thm:bundled-witness-grade-reopening-package}. Then the current bundled channel should be used with the following reopening discipline.
\begin{enumerate}[leftmargin=2em,label=(\roman*)]
\item work \emph{global-first} once the argument is already inside the bundled current channel;
\item reopen only the \emph{OP~3/QAM shell layer} when the step depends on shell certificates, shell-defect strata, shell commuting behavior, or OP~3-side carrier compression;
\item reopen only the \emph{OP~5 witness layer} when the step depends on witness grade, witness escalation, screening-witness compression, carrier-witness separation, or diagnostic collapse below escalation;
\item reopen \emph{both local layers} only when the step genuinely compares shell-grade and witness-grade data or transfers a shell-sensitive obstruction into a witness-sensitive one.
\end{enumerate}
In particular, visible OP~5 packet, envelope, flicker, and helix data do not by themselves justify a local reopening once the argument is already licensed to remain global-first.
\end{proposition}

\begin{proof}
Clause~(i) is exactly the bundled-global default from Theorem~\ref{thm:bundled-witness-grade-reopening-package}. Clause~(ii) restates the OP~3/QAM reopening part already fixed by Proposition~\ref{prop:first-global-tool-citation-guide}. Clause~(iii) is the witness-grade reopening regime from Theorem~\ref{thm:bundled-witness-grade-reopening-package}. Clause~(iv) records the only mixed local situation still licensed at the present stage: a step that must compare shell-grade and witness-grade structure rather than merely cite one local layer. The final sentence follows because, below escalation, visible OP~5 diagnostics remain subordinate to the carrier-stable visible normal form and therefore do not force reopening by appearance alone.
\end{proof}

\begin{corollary}[Witness-grade reopening checklist]
\label{cor:witness-grade-reopening-checklist}
For later proof use, the current bundled channel may be tested by the following checklist.
\begin{enumerate}[leftmargin=2em,label=(\roman*)]
\item Has the run already entered the bundled current channel?
\item Is only the bundled global pass/obstruction output needed?
\item If not, is the missing information shell-grade, witness-grade, or genuinely mixed shell/witness data?
\item Reopen only the corresponding local layer, and return to the bundled global interface immediately after the local claim has been discharged.
\end{enumerate}
This is the default proof-use checklist for the v3.19.0 witness-grade hardening step.
\end{corollary}

\begin{proof}
This is the operational reformulation of Proposition~\ref{prop:witness-grade-local-global-reopening-discipline}. The checklist simply turns the local/global reopening rule into a direct proof-use procedure.
\end{proof}

\begin{definition}[First licensed local kernel-first rewrite window]
\label{def:first-licensed-local-kernel-first-rewrite-window}
Assume the direct kernel constructor of Definition~\ref{def:qam-direct-kernel-constructor}, the bridge conditions of Theorem~\ref{thm:qam-direct-kernel-constructor}, and the bundled witness-grade reopening package of Theorem~\ref{thm:bundled-witness-grade-reopening-package}. The \emph{first licensed local kernel-first rewrite window} is the following strictly local situation: a proof step has already entered the bundled current channel, has already been forced to reopen at witness grade, and now needs only to compare source-side and return-side witness-bearing states through one common carrier-stable normal form without promoting any broader kernel-first migration on the surrounding global channel.
\end{definition}

\begin{proposition}[First licensed local kernel-first rewrite window]
\label{prop:first-licensed-local-kernel-first-rewrite-window}
Assume Definition~\ref{def:first-licensed-local-kernel-first-rewrite-window}, Proposition~\ref{prop:carrier-first-op5-interface}, and Theorem~\ref{thm:bundled-op5-diagnostic-theorem}. Let a witness-sensitive proof step concern a routed return realization already known to remain on the carrier-stable branch of the bundled OP~3$\to$OP~5 transfer channel. Then that step may be rewritten kernel-first in the following narrow sense.
\begin{enumerate}[leftmargin=2em,label=(\roman*)]
\item the visible witness-bearing source/return surfaces may be replaced by their exposed payload kernels and compared only through the common kernel normal form forced by the carrier-stable branch;
\item the local readout comparison may be performed at the level of the resulting common kernel/readout carrier class instead of on the full visible witness-bearing surfaces;
\item after that local comparison has been discharged, the proof must return immediately to the bundled carrier-stable/global interface and may not infer a broader kernel-first rewrite rule for the full OP~3$\to$OP~5/global channel.
\end{enumerate}
Hence the present line now contains one first licensed kernel-first rewrite window: local witness-grade carrier comparison under an already fixed carrier-stable branch, but no general kernel-first migration of the whole bundled channel.
\end{proposition}

\begin{proof}
Once the proof step is already known to lie on the carrier-stable branch, the witness-bearing visible data are subordinate to one common closure-bearing carrier package. Under the direct kernel constructor and the bridge conditions, this is exactly the situation in which the local comparison load can be shifted from full visible surfaces to their exposed payload kernels and then to the corresponding common kernel normal form. The local readout comparison may therefore be discharged at the level of the common kernel/readout carrier class. However, Theorem~\ref{thm:bundled-witness-grade-reopening-package} and Definition~\ref{def:first-licensed-local-kernel-first-rewrite-window} keep this move strictly local: the argument was already forced into witness-grade reopening, and nothing in the present hypotheses licenses promotion of that local kernel-first step to a global rewrite rule for the surrounding bundled channel. The final sentence is exactly this restriction written as a proof-use statement.
\end{proof}

\begin{corollary}[Use rule for the first local kernel-first rewrite window]
\label{cor:use-rule-for-the-first-local-kernel-first-rewrite-window}
Whenever a later proof step is already inside the bundled current channel, already forced to reopen at witness grade, and already known to remain on the carrier-stable branch, it is sufficient to cite Proposition~\ref{prop:first-licensed-local-kernel-first-rewrite-window} for the local source/return comparison. After that local comparison, the proof should return to the bundled global interface immediately.
\end{corollary}

\begin{proof}
This is the operational reformulation of Proposition~\ref{prop:first-licensed-local-kernel-first-rewrite-window}. The proposition already states both the local license and the immediate return rule.
\end{proof}


\begin{remark}[Status of the first corridor-level activation step]
\label{rem:status-of-the-first-corridor-level-activation-step}
The present release line still does \emph{not} activate a kernel-first rewrite rule for the whole bundled OP~3$\to$OP~5/global channel. Nevertheless, the combination of the bundled witness-grade reopening package and the first licensed local kernel-first rewrite window already supports one slightly stronger activation step: a finite witness-grade corridor on one fixed carrier-stable branch may be discharged kernel-first through one common carrier-stable kernel/readout class. This is the first point at which the current release line moves beyond a single isolated local rewrite without yet promoting the deferred kernel architecture to an active global migration layer.
\end{remark}

\begin{definition}[Carrier-stable witness corridor]
\label{def:carrier-stable-witness-corridor}
Assume Theorem~\ref{thm:bundled-witness-grade-reopening-package} and Proposition~\ref{prop:first-licensed-local-kernel-first-rewrite-window}. A \emph{carrier-stable witness corridor} is a finite family of witness-sensitive proof steps
\[
\Sigma_1,\ldots,\Sigma_N
\]
along one routed return run such that:
\begin{enumerate}[leftmargin=2em,label=(\roman*)]
\item the run is already known to lie inside the bundled current channel;
\item every step $\Sigma_j$ is already forced to reopen at witness grade and remains on the same carrier-stable branch;
\item for every $j$, the local source/return comparison occurring in $\Sigma_j$ satisfies the hypotheses of Proposition~\ref{prop:first-licensed-local-kernel-first-rewrite-window};
\item no step in the family changes the fixed carrier-stable normal-form class or requires a new shell/witness transfer beyond that same carrier-stable branch.
\end{enumerate}
Thus a carrier-stable witness corridor is the smallest finite witness-grade subrun on which the same carrier-stable kernel/readout class may be reused without reopening the whole bundled channel anew at each local comparison.
\end{definition}

\begin{lemma}[First corridor-level kernel-first activation lemma]
\label{lem:first-corridor-level-kernel-first-activation}
Assume Definition~\ref{def:carrier-stable-witness-corridor}. Then the whole witness corridor may be discharged kernel-first through one common carrier-stable kernel/readout class in the following precise sense.
\begin{enumerate}[leftmargin=2em,label=(\roman*)]
\item for every local step $\Sigma_j$, the visible witness-bearing source/return surfaces may be replaced by their exposed payload kernels and compared only through the same carrier-stable kernel normal-form class;
\item the local readout comparisons occurring across the full finite family $\Sigma_1,\ldots,\Sigma_N$ may therefore be performed at class level rather than re-established on the full visible witness-bearing surfaces step by step;
\item after the corridor has been discharged, the proof must return to the bundled global interface immediately and may not infer a general kernel-first rewrite rule for branch-changing, shell-changing, or globally mixed witness-grade arguments.
\end{enumerate}
Equivalently, the present line already licenses one first corridor-level activation step: a finite bundled family of witness-grade local comparisons along one fixed carrier-stable branch, but still no general kernel-first migration of the whole bundled channel.
\end{lemma}

\begin{proof}
By Definition~\ref{def:carrier-stable-witness-corridor}, every local step $\Sigma_j$ already satisfies the hypotheses of Proposition~\ref{prop:first-licensed-local-kernel-first-rewrite-window}. Hence each such step may be rewritten kernel-first at the local witness-sensitive level. Because all steps lie on the same carrier-stable branch and no step changes the fixed carrier-stable normal-form class, the corresponding local kernel/readout comparisons all factor through one and the same carrier-stable kernel/readout class rather than through a new class at every step. Therefore the whole finite family may be discharged kernel-first corridor by corridor instead of reopening the full visible witness-bearing surfaces separately at each local comparison.

The return rule is unchanged: Theorem~\ref{thm:bundled-witness-grade-reopening-package} keeps witness-grade reopening subordinate to the bundled global interface, and Definition~\ref{def:carrier-stable-witness-corridor} excludes branch-changing, shell-changing, or globally mixed witness-grade transitions. Consequently, once the corridor has been discharged, the proof must return immediately to the bundled global interface, and no broader global kernel-first migration follows from this corridor-level argument.
\end{proof}

\begin{corollary}[Use rule for the first corridor-level activation lemma]
\label{cor:use-rule-for-the-first-corridor-level-activation-lemma}
Whenever a later proof contains a finite witness-grade subrun on one fixed carrier-stable branch, it is sufficient to cite Lemma~\ref{lem:first-corridor-level-kernel-first-activation} instead of reopening each witness-bearing source/return comparison separately. After that corridor-level comparison has been discharged, the proof should return to the bundled global interface immediately.
\end{corollary}

\begin{proof}
This is the operational reformulation of Lemma~\ref{lem:first-corridor-level-kernel-first-activation}. The lemma already states both the corridor-level license and the immediate return rule.
\end{proof}




\begin{remark}[Status of the first corridor-to-bundle propagation step]
\label{rem:status-of-the-first-corridor-to-bundle-propagation-step}
The present release line still does \emph{not} activate a kernel-first rewrite rule for the whole bundled OP~3$\to$OP~5/global channel. Nevertheless, the combination of the bundled witness-grade reopening package, the first licensed local kernel-first rewrite window, and the first corridor-level activation lemma already supports one slightly stronger propagation step: a bundled run may now alternate between globally readable stretches and finitely many carrier-stable witness corridors without forcing each witness-bearing local comparison to be reopened from scratch. This is the first point at which the current release line moves beyond one isolated corridor-level activation and begins to control a whole bundled run by a finite alternation of global citations and corridor-level kernel-first discharges.
\end{remark}

\begin{definition}[Bundle-compatible witness-corridor decomposition]
\label{def:bundle-compatible-witness-corridor-decomposition}
Assume Theorem~\ref{thm:bundled-witness-grade-reopening-package}, Lemma~\ref{lem:first-corridor-level-kernel-first-activation}, and Proposition~\ref{prop:first-global-tool-citation-guide}. A \emph{bundle-compatible witness-corridor decomposition} of a later proof run is a finite alternating decomposition
\[
\Gamma_0,\ \mathcal C_1,\ \Gamma_1,\ \ldots,\ \mathcal C_M,\ \Gamma_M
\]
such that:
\begin{enumerate}[leftmargin=2em,label=(\roman*)]
\item the entire run is already known to lie inside the bundled current channel;
\item each global segment $\Gamma_i$ needs only the bundled outputs global compressed pass/global compressed obstruction together with their induced carrier-stable or obstruction-localized consequences;
\item each local segment $\mathcal C_k$ is a carrier-stable witness corridor in the sense of Definition~\ref{def:carrier-stable-witness-corridor};
\item after each local corridor $\mathcal C_k$, the proof returns immediately to the bundled global interface before any later corridor begins;
\item no corridor changes the fixed bundled run into a new shell-changing, branch-changing, or globally mixed witness-grade run between two global segments.
\end{enumerate}
Thus a bundle-compatible witness-corridor decomposition is the smallest finite alternation in which the bundled global interface remains the default proof surface and only finitely many carrier-stable witness corridors are discharged locally.
\end{definition}

\begin{proposition}[First corridor-to-bundle propagation principle]
\label{prop:first-corridor-to-bundle-propagation-principle}
Assume Definition~\ref{def:bundle-compatible-witness-corridor-decomposition}. Then the whole decomposed run may be discharged by the following finite proof pattern.
\begin{enumerate}[leftmargin=2em,label=(\roman*)]
\item on each global segment $\Gamma_i$, it is sufficient to cite the bundled global interface and work only with the dichotomy
\[
\text{global compressed pass}
\quad\text{versus}\quad
\text{global compressed obstruction};
\]
\item on each local corridor $\mathcal C_k$, it is sufficient to cite Lemma~\ref{lem:first-corridor-level-kernel-first-activation} and discharge the whole finite witness-grade corridor through one common carrier-stable kernel/readout class;
\item after each local corridor $\mathcal C_k$, the proof returns immediately to the bundled global interface and resumes global-first reasoning on the next segment $\Gamma_k$;
\item consequently, no individual witness-bearing source/return comparison inside the corridors needs to be reopened separately, and no general kernel-first rewrite rule for the full bundled channel is inferred.
\end{enumerate}
Equivalently, the present line already licenses one first corridor-to-bundle propagation rule: a bundled run may be handled by finitely many global-interface citations plus finitely many corridor-level kernel-first discharges, while the active proof surface remains globally bundled outside those local corridors.
\end{proposition}

\begin{proof}
By Definition~\ref{def:bundle-compatible-witness-corridor-decomposition}, every global segment $\Gamma_i$ satisfies exactly the hypotheses of the bundled global citation discipline, so Proposition~\ref{prop:first-global-tool-citation-guide} supplies the global compressed pass/global compressed obstruction interface there. Likewise, every local segment $\mathcal C_k$ is a carrier-stable witness corridor, so Lemma~\ref{lem:first-corridor-level-kernel-first-activation} discharges that whole finite corridor kernel-first through one common carrier-stable kernel/readout class.

The return rule is part of both ingredients: the global interface is the default outside local reopenings, while the corridor lemma requires immediate return to the bundled global interface after each local corridor has been discharged. Hence the full run is controlled by a finite alternation of global citations and corridor-level kernel-first discharges. Because the corridor lemma already absorbs the local witness-grade source/return comparisons inside each $\mathcal C_k$, no individual witness-bearing comparison has to be reopened separately. Finally, nothing in the hypotheses licenses promotion of this finite alternation to a general kernel-first rewrite rule for the whole bundled channel; outside the listed corridors the active proof surface remains the bundled global interface.
\end{proof}

\begin{corollary}[Use rule for the first corridor-to-bundle propagation principle]
\label{cor:use-rule-for-the-first-corridor-to-bundle-propagation-principle}
Whenever a later proof run already inside the bundled current channel decomposes into globally readable stretches separated by finitely many carrier-stable witness corridors, it is sufficient to cite Proposition~\ref{prop:first-corridor-to-bundle-propagation-principle} together with Lemma~\ref{lem:first-corridor-level-kernel-first-activation}, rather than reopening each witness-bearing local comparison separately. Only the corridor boundaries, not the internal witness steps of each corridor, remain visible at proof-use level.
\end{corollary}

\begin{proof}
This is the operational reformulation of Proposition~\ref{prop:first-corridor-to-bundle-propagation-principle}. The proposition already states both the finite propagation rule and the immediate return to the bundled global interface after each corridor.
\end{proof}


\begin{remark}[Status of the first bundle-level reopening compression step]
\label{rem:status-of-the-first-bundle-level-reopening-compression-step}
The present release line still does \emph{not} collapse every local reopening beneath the bundled channel to one uniform active theorem surface. What it now does support, however, is one further compression of proof-use data: once a bundled run has already been decomposed into global stretches and finitely many carrier-stable witness corridors, the later proof no longer needs to describe each corridor by all of its internal shell-grade and witness-grade subclaims. It is enough to record, corridor by corridor, which local reopening profile is actually being used. This is the first point at which the current release line compresses not only the internal witness comparisons of each corridor, but also the reopening metadata attached to the whole bundled run.
\end{remark}

\begin{definition}[Profile-stable bundle-level reopening data]
\label{def:profile-stable-bundle-level-reopening-data}
Assume Definition~\ref{def:bundle-compatible-witness-corridor-decomposition}. A \emph{profile-stable bundle-level reopening datum} on such a decomposition is a choice of one profile label
\[
\pi_k\in\{\mathrm{Q},\mathrm{W},\mathrm{QW}\}
\]
for every corridor $\mathcal C_k$ such that:
\begin{enumerate}[leftmargin=2em,label=(\roman*)]
\item $\pi_k=\mathrm{Q}$ means that the corridor is reopened only through the OP~3/QAM shell side: shell certificates, shell-obstruction stratification, shell commuting behavior, or OP~3-side carrier compression;
\item $\pi_k=\mathrm{W}$ means that the corridor is reopened only through the carrier-first OP~5 witness side: witness grade, witness escalation, screening-witness compression, carrier-witness separation, or diagnostic collapse beneath the escalation threshold;
\item $\pi_k=\mathrm{QW}$ means that the corridor genuinely compares shell-grade and witness-grade data inside one common carrier-stable witness corridor;
\item the profile label of a corridor remains fixed throughout that corridor and changes, if at all, only at corridor boundaries.
\end{enumerate}
Thus the profile sequence $(\pi_1,\ldots,\pi_M)$ records the smallest local reopening data needed at proof-use level beyond the bundled global interface.
\end{definition}

\begin{proposition}[First bundle-level reopening compression principle]
\label{prop:first-bundle-level-reopening-compression-principle}
Assume a bundle-compatible witness-corridor decomposition together with profile-stable bundle-level reopening data in the sense of Definition~\ref{def:profile-stable-bundle-level-reopening-data}. Then the whole bundled run may be handled at proof-use level by the finite data
\[
\Gamma_0,\ \pi_1,\ \Gamma_1,\ \ldots,\ \pi_M,\ \Gamma_M
\]
with the following force.
\begin{enumerate}[leftmargin=2em,label=(\roman*)]
\item Each global segment $\Gamma_i$ is still treated only through the bundled global interface and the dichotomy global compressed pass versus global compressed obstruction.
\item Each corridor $\mathcal C_k$ is reopened only through its single profile label $\pi_k$, together with Lemma~\ref{lem:first-corridor-level-kernel-first-activation} whenever the corridor is discharged kernel-first as one carrier-stable witness corridor.
\item No later proof needs to list the internal shell-grade or witness-grade subclaims of $\mathcal C_k$ separately as long as they remain subordinate to the same fixed corridor profile $\pi_k$.
\item If a corridor would force a profile change inside itself, then the present compression rule stops at that point and the corridor must be split into smaller profile-stable corridors before the bundled proof-use compression is applied again.
\end{enumerate}
Equivalently, the current release line already licenses one first bundle-level reopening compression step: after corridor-level kernel-first discharge has been fixed, the remaining local proof-use burden of a bundled run is reduced to a finite sequence of profile labels attached to finitely many carrier-stable corridors.
\end{proposition}

\begin{proof}
By Proposition~\ref{prop:first-corridor-to-bundle-propagation-principle}, the bundled run is already controlled by finitely many global citations together with finitely many corridor-level kernel-first discharges. The only remaining local proof-use burden is therefore to specify which finer local reopening regime is actually needed on each corridor. Definition~\ref{def:profile-stable-bundle-level-reopening-data} makes exactly that finite burden explicit: every corridor carries one fixed profile label $\pi_k$, and that label determines whether the local reopening is OP~3/QAM-only, OP~5-witness-only, or genuinely mixed.

Because the profile remains fixed throughout each corridor, all internal shell-grade and witness-grade subclaims of that corridor stay subordinate to the same local reopening regime and need not be listed one by one at proof-use level. If a profile change occurred inside a corridor, then the corridor would no longer be profile-stable, and the current compression would no longer be legitimate without first splitting that corridor into smaller profile-stable pieces. This proves both the finite proof-use description and the stated stopping rule.
\end{proof}

\begin{corollary}[Use rule for the first bundle-level reopening compression principle]
\label{cor:use-rule-for-the-first-bundle-level-reopening-compression-principle}
Whenever a later proof run already inside the bundled current channel admits both a bundle-compatible witness-corridor decomposition and profile-stable bundle-level reopening data, it is sufficient to cite Proposition~\ref{prop:first-bundle-level-reopening-compression-principle} and record only the resulting finite profile sequence $(\pi_1,\ldots,\pi_M)$, rather than restating all local shell-grade and witness-grade subclaims corridor by corridor. A finer reopening is required only when one corridor changes profile internally or leaves the carrier-stable witness-corridor regime.
\end{corollary}

\begin{proof}
This is the operational reformulation of Proposition~\ref{prop:first-bundle-level-reopening-compression-principle}. The proposition already states both the compressed proof-use description and the precise point where that compression stops.
\end{proof}




\begin{remark}[Status of the first profile-to-run compression step]
\label{rem:status-of-the-first-profile-to-run-compression-step}
The present release line still does \emph{not} identify all bundled runs by one fully kernel-activated normal form. What it now does support, however, is one further compression of the proof-use surface: once a bundled run has already been decomposed into global stretches and finitely many profile-stable carrier corridors, the later proof no longer needs to remember the names of the individual global segments or the internal witness subclaims of those corridors. It is enough to remember the alternating \emph{run profile word} that records only when the proof is global and which local reopening profile occurs at each corridor. This is the first point at which the current release line compresses a whole bundled run to one finite symbolic proof-use word.
\end{remark}

\begin{definition}[Run-stable bundled proof word]
\label{def:run-stable-bundled-proof-word}
Assume a bundle-compatible witness-corridor decomposition together with profile-stable bundle-level reopening data in the sense of Definition~\ref{def:profile-stable-bundle-level-reopening-data}. The associated \emph{run-stable bundled proof word} is the finite alternating word
\[
\omega(\Gamma):=\mathrm{G}\,\pi_1\,\mathrm{G}\,\pi_2\,\mathrm{G}\cdots \pi_M\,\mathrm{G}
\]
over the alphabet
\[
\{\mathrm{G},\mathrm{Q},\mathrm{W},\mathrm{QW}\},
\]
where:
\begin{enumerate}[leftmargin=2em,label=(\roman*)]
\item each letter $\mathrm{G}$ records one globally readable bundled segment handled only through the dichotomy global compressed pass versus global compressed obstruction;
\item each letter $\pi_k\in\{\mathrm{Q},\mathrm{W},\mathrm{QW}\}$ records the fixed reopening profile of the corresponding carrier-stable witness corridor $\mathcal C_k$;
\item the word begins and ends with $\mathrm{G}$, and the letters alternate globally/local in the same order as the bundled run;
\item no letter change is allowed inside a corridor; if such a change occurs, the current word is not defined until the corridor has been split into smaller profile-stable corridors.
\end{enumerate}
Thus $\omega(\Gamma)$ is the smallest finite symbolic record that still remembers where the proof is global and which local reopening profile occurs whenever the bundled channel is reopened.
\end{definition}

\begin{proposition}[First profile-to-run compression principle]
\label{prop:first-profile-to-run-compression-principle}
Assume Definition~\ref{def:run-stable-bundled-proof-word}. Then the proof-use burden of the whole bundled run is reduced to its finite run-stable bundled proof word
\[
\omega(\Gamma)\in\{\mathrm{G},\mathrm{Q},\mathrm{W},\mathrm{QW}\}^{*}
\]
with the following force.
\begin{enumerate}[leftmargin=2em,label=(\roman*)]
\item Every occurrence of $\mathrm{G}$ means that the corresponding segment is handled only through the bundled global interface and its compressed pass/obstruction dichotomy.
\item Every occurrence of $\mathrm{Q}$, $\mathrm{W}$, or $\mathrm{QW}$ means that the next corridor is reopened only in that single local profile, together with Lemma~\ref{lem:first-corridor-level-kernel-first-activation} whenever the corridor is licensed for carrier-stable kernel-first discharge.
\item No later proof needs to retain the individual names of the global segments $\Gamma_i$ or the separate internal shell-grade/witness-grade clauses of the corridors once the alternating word $\omega(\Gamma)$ has been fixed.
\item If one wants to refine a corridor internally, change its local reopening profile midstream, or distinguish two globally readable segments by more than the bundled pass/obstruction interface, then the present compression stops and the run must first be rewritten at the preceding bundle-level stage.
\end{enumerate}
Equivalently, the current release line already licenses one first profile-to-run compression step: after bundle-compatible corridor decomposition and profile stabilization have been fixed, the whole bundled proof surface may be cited by one finite alternating word over $\{\mathrm{G},\mathrm{Q},\mathrm{W},\mathrm{QW}\}$.
\end{proposition}

\begin{proof}
By Proposition~\ref{prop:first-bundle-level-reopening-compression-principle}, the bundled run is already controlled by finitely many global segments together with finitely many profile-stable carrier corridors. The only remaining proof-use information is therefore the order in which those two types of pieces occur and, on each corridor, which fixed reopening profile is active. Definition~\ref{def:run-stable-bundled-proof-word} records exactly that data as the alternating word $\omega(\Gamma)$.

The force of each letter is already supplied by the previous stages: $\mathrm{G}$ stands for the bundled global interface, while $\mathrm{Q}$, $\mathrm{W}$, and $\mathrm{QW}$ stand for the corresponding local reopening profiles and, where licensed, for the corridor-level kernel-first discharge of Lemma~\ref{lem:first-corridor-level-kernel-first-activation}. Since the previous bundle-level compression already removes the need to list individual internal witness clauses corridor by corridor, fixing the alternating word removes the need to retain the names of the global segments as separate proof-use data. If one wants finer information than the word provides, then one has already left the current compressed regime and must reopen the run at the preceding bundle-level stage. This proves the stated stopping rule.
\end{proof}

\begin{corollary}[Use rule for the first profile-to-run compression principle]
\label{cor:use-rule-for-the-first-profile-to-run-compression-principle}
Whenever a later proof run already inside the bundled current channel admits a run-stable bundled proof word $\omega(\Gamma)$, it is sufficient to cite Proposition~\ref{prop:first-profile-to-run-compression-principle} and record only that finite word, rather than listing the separate global segments and corridor-internal local claims again. A finer reopening is required only when a corridor changes profile internally, leaves the carrier-stable witness-corridor regime, or one global segment must be distinguished from another by more than the bundled pass/obstruction interface.
\end{corollary}

\begin{proof}
This is the operational reformulation of Proposition~\ref{prop:first-profile-to-run-compression-principle}. The proposition already states both the finite run-word description and the precise point where that description stops being legitimate.
\end{proof}


\begin{remark}[Status of the first word-determined run-class step]
\label{rem:status-of-the-first-word-determined-run-class-step}
The present release line still does \emph{not} identify all bundled runs by one globally activated kernel normal form, nor does it erase the distinction between different geometric realizations of the same bundled channel. What it now does support, however, is one further proof-use quotient: once a bundled run has already been compressed to a run-stable bundled proof word, the later proof may forget the individual realization of that run and retain only the \emph{run class determined by that word}. This is the first point at which the current release line passes from one finite symbolic proof-use word to one finite symbolic proof-use class.
\end{remark}

\begin{definition}[Word-determined bundled run class]
\label{def:word-determined-bundled-run-class}
Assume Definition~\ref{def:run-stable-bundled-proof-word}. Two bundled runs $\Gamma$ and $\Gamma'$ are called \emph{word-determined proof-use equivalent} if both run-stable bundled proof words are defined and satisfy
\[
\omega(\Gamma)=\omega(\Gamma').
\]
The corresponding equivalence class is written
\[
[\Gamma]_{\omega}:=\{\Gamma' : \omega(\Gamma')=\omega(\Gamma)\}.
\]
It is called the \emph{word-determined bundled run class} of $\Gamma$. In this reading, the class $[\Gamma]_{\omega}$ records exactly the proof-use information retained by the current compressed line and forgets the names of individual global segments, the particular realization of each carrier-stable witness corridor, and all corridor-internal shell-grade or witness-grade subclaims beyond what is already encoded by the letters of $\omega(\Gamma)$.
\end{definition}

\begin{proposition}[First word-to-run-class compression principle]
\label{prop:first-word-to-run-class-compression-principle}
Assume Definition~\ref{def:word-determined-bundled-run-class}. Then the current proof-use burden of a bundled run depends only on its word-determined bundled run class $[\Gamma]_{\omega}$ and no longer on the individual run $\Gamma$ itself, with the following force.
\begin{enumerate}[leftmargin=2em,label=(\roman*)]
\item If two bundled runs $\Gamma$ and $\Gamma'$ satisfy $\omega(\Gamma)=\omega(\Gamma')$, then every later argument that uses only the bundled global interface, the associated local reopening profiles, and the corresponding corridor-level kernel-first licenses may treat $\Gamma$ and $\Gamma'$ as the same proof-use object.
\item Consequently, once the run-stable bundled proof word has been fixed, the later proof may cite the class $[\Gamma]_{\omega}$ instead of the particular run realization $\Gamma$.
\item The present compression stops exactly when one needs information not encoded by the word: for example, a distinction between two globally readable segments carrying the same letter $\mathrm G$, a distinction between two carrier-stable witness corridors carrying the same profile letter, or a geometric/property-level comparison internal to one corridor that is invisible at word level.
\end{enumerate}
Equivalently, the current release line already licenses one first word-to-run-class compression step: after passage from bundled runs to run-stable bundled proof words, proof use depends only on the induced equivalence class under equality of those words.
\end{proposition}

\begin{proof}
By Proposition~\ref{prop:first-profile-to-run-compression-principle}, the current proof-use surface of a bundled run is already reduced to the finite alternating word $\omega(\Gamma)$. Hence every later argument that cites only the force carried by the letters $\mathrm G$, $\mathrm Q$, $\mathrm W$, and $\mathrm{QW}$, together with the associated corridor-level kernel-first license where available, uses no further data from the particular realization of $\Gamma$.

Therefore, if $\omega(\Gamma)=\omega(\Gamma')$, then both runs carry exactly the same compressed proof-use data and may be treated as the same proof-use object at the current level of abstraction. This proves (i) and the resulting class-level citation claim in (ii). Clause (iii) is simply the contrapositive reading of the same compression: whenever a later argument needs data not already encoded by the word, the quotient by equality of words is too coarse, so the proof must reopen at the preceding run-level or bundle-level stage.
\end{proof}

\begin{corollary}[Use rule for the first word-to-run-class compression principle]
\label{cor:use-rule-for-the-first-word-to-run-class-compression-principle}
Whenever a later proof already inside the bundled current channel depends only on the run-stable bundled proof word and the force of its letters, it is sufficient to cite Proposition~\ref{prop:first-word-to-run-class-compression-principle} and work with the corresponding class $[\Gamma]_{\omega}$, rather than with one chosen realization of the bundled run. A finer reopening is required only when the argument needs information not encoded by the word itself.
\end{corollary}

\begin{proof}
This is the operational reformulation of Proposition~\ref{prop:first-word-to-run-class-compression-principle}. The proposition already states both the induced class-level proof-use rule and the precise point where that compression stops.
\end{proof}




\begin{remark}[Status of the bundled proof-use compression stack in the present v\docversion\ release line]
\label{rem:status-of-the-first-run-class-concatenation-step}
The present release line still does \emph{not} equip bundled current-channel reasoning with one unrestricted global kernel algebra. What it now supports, however, is a much shorter proof-use stack: bundled runs may be compressed successively from run words to run classes, from finite concatenation chains to one reduced normal-form word, from that word to one proof-use class, from binary and finite class composition to one associative class word, and finally to one contraction-stable terminal citation surface. The next items record only the exact proof-use interfaces needed later; the longer stage-by-stage derivations are no longer repeated here.
\end{remark}

\begin{definition}[Concatenable word-determined bundled run classes]
\label{def:concatenable-word-determined-bundled-run-classes}
Assume Definition~\ref{def:word-determined-bundled-run-class}. Two word-determined bundled run classes $[\Gamma]_{\omega}$ and $[\Gamma']_{\omega}$ are \emph{concatenable at bundled level} if their join is handled entirely by the bundled global interface and no new local reopening is created at the splice. In that case define the reduced concatenated run word by
\[
\omega(\Gamma)\star_{\mathrm G}\omega(\Gamma')
\]
obtained by identifying the terminal global letter $\mathrm G$ of the first word with the initial global letter $\mathrm G$ of the second, and write
\[
[\Gamma]_{\omega}\star[\Gamma']_{\omega}
\]
for the induced concatenated run-class citation object.
\end{definition}

\begin{proposition}[First run-class concatenation principle]
\label{prop:first-run-class-concatenation-principle}
Under Definition~\ref{def:concatenable-word-determined-bundled-run-classes}, bundled-level concatenation depends only on the two input run classes and on the reduced concatenated run word. In particular, later proof use may cite $[\Gamma]_{\omega}\star[\Gamma']_{\omega}$ instead of chosen representatives, and this compression stops exactly when the splice creates local data not already encoded by the adjacent global-interface letters.
\end{proposition}

\begin{proof}
Immediate from the word-level quotient already fixed by Proposition~\ref{prop:first-word-to-run-class-compression-principle}: once representatives are forgotten, the only remaining bundled splice datum is the reduced global-interface identification. Any genuinely new local burden at the join lies beyond this quotient and forces reopening.
\end{proof}

\begin{corollary}[Use rule for the first run-class concatenation principle]
\label{cor:use-rule-for-the-first-run-class-concatenation-principle}
Whenever two bundled current-channel runs are already compressed to word-determined bundled run classes and their join is handled entirely by the bundled global interface, it is sufficient to cite Proposition~\ref{prop:first-run-class-concatenation-principle} and work with the concatenated class $[\Gamma]_{\omega}\star[\Gamma']_{\omega}$.
\end{corollary}

\begin{proof}
This is the operational reading of Proposition~\ref{prop:first-run-class-concatenation-principle}.
\end{proof}

\begin{remark}[Status of the first finite run-class normal-form step]
\label{rem:status-of-the-first-finite-run-class-normal-form-step}
The next compression replaces a finite ordered chain of concatenable bundled run classes by one reduced normal-form word. Only the ordered local reopening profile and the bundled global separators survive.
\end{remark}

\begin{definition}[Finite concatenable bundled run-class chain and reduced normal-form word]
\label{def:finite-concatenable-bundled-run-class-chain-and-reduced-normal-form-word}
Let
\[
[\Gamma_1]_{\omega},\dots,[\Gamma_n]_{\omega}
\]
be bundled run classes such that each adjacent pair is concatenable at bundled level. Define the reduced normal-form concatenation word
\[
\operatorname{NF}_{\star}([\Gamma_1]_{\omega},\dots,[\Gamma_n]_{\omega})
\]
by concatenating the words $\omega(\Gamma_k)$ and identifying every adjacent terminal/initial global letter $\mathrm G$ as one common bundled global-interface letter.
\end{definition}

\begin{proposition}[First finite run-class normal-form principle]
\label{prop:first-finite-run-class-normal-form-principle}
For every finite bundled concatenation chain covered by Definition~\ref{def:finite-concatenable-bundled-run-class-chain-and-reduced-normal-form-word}, the proof-use burden depends only on the ordered list of run classes and on the reduced normal-form word $\operatorname{NF}_{\star}$. Representative choices and parenthesization of successive concatenations do not matter.
\end{proposition}

\begin{proof}
Repeated use of Proposition~\ref{prop:first-run-class-concatenation-principle} shows that each bundled-level splice is governed only by the reduced global-interface identification. Iterating that splice therefore yields one parenthesization-independent reduced word.
\end{proof}

\begin{corollary}[Use rule for the first finite run-class normal-form principle]
\label{cor:use-rule-for-the-first-finite-run-class-normal-form-principle}
Whenever a later argument uses only the ordered local profile carried by a finite concatenable run-class chain, it is sufficient to cite Proposition~\ref{prop:first-finite-run-class-normal-form-principle} and work with $\operatorname{NF}_{\star}$.
\end{corollary}

\begin{proof}
This is the operational reading of Proposition~\ref{prop:first-finite-run-class-normal-form-principle}.
\end{proof}

\begin{remark}[Status of the first normal-form-determined bundled proof-use class step]
\label{rem:status-of-the-first-normal-form-determined-bundled-proof-use-class-step}
Once a finite concatenable bundled run-class chain has been reduced to one normal-form word, later proof use may forget the chosen chain realization and retain only the class determined by that word.
\end{remark}

\begin{definition}[Normal-form-determined bundled proof-use class]
\label{def:normal-form-determined-bundled-proof-use-class}
Two finite concatenable bundled run-class chains are called normal-form proof-use equivalent if their reduced normal-form words agree. The corresponding equivalence class is the \emph{normal-form-determined bundled proof-use class} and is written $[\mathfrak C]_{\operatorname{NF}}$.
\end{definition}

\begin{proposition}[First normal-form-to-proof-use-class compression principle]
\label{prop:first-normal-form-to-proof-use-class-compression-principle}
The current proof-use burden of a finite concatenable bundled run-class chain depends only on its normal-form-determined bundled proof-use class. Equivalently, once only the reduced normal-form word and the force of its letters matter, one may replace the chain by $[\mathfrak C]_{\operatorname{NF}}$.
\end{proposition}

\begin{proof}
This is precisely the quotient induced by equality of reduced normal-form words from Definition~\ref{def:normal-form-determined-bundled-proof-use-class}.
\end{proof}

\begin{corollary}[Use rule for the first normal-form-to-proof-use-class compression principle]
\label{cor:use-rule-for-the-first-normal-form-to-proof-use-class-compression-principle}
Whenever a later proof already inside the bundled current channel depends only on the reduced normal-form concatenation word and the force of its letters, it is sufficient to cite Proposition~\ref{prop:first-normal-form-to-proof-use-class-compression-principle} and work with the corresponding class $[\mathfrak C]_{\operatorname{NF}}$.
\end{corollary}

\begin{proof}
This is the operational reading of Proposition~\ref{prop:first-normal-form-to-proof-use-class-compression-principle}.
\end{proof}

\begin{remark}[Status of the first proof-use-class composition step]
\label{rem:status-of-the-first-proof-use-class-composition-step}
After passage to normal-form-determined classes, bundled composition can be handled at class level whenever the splice itself adds no new local burden.
\end{remark}

\begin{definition}[Composable normal-form-determined bundled proof-use classes]
\label{def:composable-normal-form-determined-bundled-proof-use-classes}
Two normal-form-determined bundled proof-use classes are \emph{composable} if realizations may be joined entirely through the bundled global interface with no new local reopening at the splice. Their reduced composed normal-form word is obtained by the same $\star_{\mathrm G}$ splice on the two reduced normal-form words.
\end{definition}

\begin{proposition}[First proof-use-class composition principle]
\label{prop:first-proof-use-class-composition-principle}
Whenever two normal-form-determined bundled proof-use classes are composable, their bundled composition depends only on the two classes and on the reduced composed normal-form word. Representative-level choices are invisible at this class-compression stage.
\end{proposition}

\begin{proof}
Combine Proposition~\ref{prop:first-normal-form-to-proof-use-class-compression-principle} with the same bundled-global splice rule already used at run-class level.
\end{proof}

\begin{corollary}[Use rule for the first proof-use-class composition principle]
\label{cor:use-rule-for-the-first-proof-use-class-composition-principle}
Whenever two bundled current-channel proof packages are already compressed to normal-form-determined bundled proof-use classes and their splice is handled entirely by the bundled global interface, it is sufficient to cite Proposition~\ref{prop:first-proof-use-class-composition-principle} and work with their composed class.
\end{corollary}

\begin{proof}
This is the operational reading of Proposition~\ref{prop:first-proof-use-class-composition-principle}.
\end{proof}

\begin{remark}[Status of the first associative proof-use-class normal-form step]
\label{rem:status-of-the-first-associative-proof-use-class-normal-form-step}
Iterated class composition can now be read through one parenthesization-stable reduced class word.
\end{remark}

\begin{definition}[Associative reduced class-composition normal form]
\label{def:associative-reduced-class-composition-normal-form}
For a finite ordered composable family of normal-form-determined bundled proof-use classes, define the associative reduced class-composition normal form by concatenating the reduced class words and identifying every interior terminal/initial global letter $\mathrm G$ as one common bundled global-interface letter.
\end{definition}

\begin{proposition}[First associative proof-use-class normal-form principle]
\label{prop:first-associative-proof-use-class-normal-form-principle}
The proof-use burden of any finite ordered composable family of normal-form-determined bundled proof-use classes depends only on its associative reduced class-composition normal form and is independent of parenthesization.
\end{proposition}

\begin{proof}
At word level the bundled-global splice is associative; Proposition~\ref{prop:first-proof-use-class-composition-principle} removes dependence on representatives, so only the resulting reduced class word remains.
\end{proof}

\begin{corollary}[Use rule for the first associative proof-use-class normal-form principle]
\label{cor:use-rule-for-the-first-associative-proof-use-class-normal-form-principle}
Whenever a later bundled current-channel proof uses only the force of a finite ordered composable family of normal-form-determined bundled proof-use classes and all intermediate joins are handled entirely by the bundled global interface, it is sufficient to cite Proposition~\ref{prop:first-associative-proof-use-class-normal-form-principle} and work with the single associative class word.
\end{corollary}

\begin{proof}
This is the operational reading of Proposition~\ref{prop:first-associative-proof-use-class-normal-form-principle}.
\end{proof}

\begin{remark}[Status of the first contraction-stable class-normal-form step]
\label{rem:status-of-the-first-contraction-stable-class-normal-form-step}
The final local compression removes only redundant consecutive global separators and leaves the genuinely reopened local content unchanged.
\end{remark}

\begin{definition}[Contraction-stable reduced class-normal-form word]
\label{def:contraction-stable-reduced-class-normal-form-word}
Given a finite composable family of normal-form-determined bundled proof-use classes with associative reduced class word $\omega_{\mathrm{assoc}}$, define
\[
\operatorname{NF}_{\mathrm{ctr}}
\]
by replacing each maximal consecutive block of the symbol $G$ in $\omega_{\mathrm{assoc}}$ by a single occurrence of $G$.
\end{definition}

\begin{proposition}[First contraction-stable class-normal-form principle]
\label{prop:first-contraction-stable-class-normal-form-principle}
The contraction-stable class-normal-form word $\operatorname{NF}_{\mathrm{ctr}}$ is independent of representatives and parenthesization and determines the same proof-use surface as the associative class-normal-form word, except that redundant purely global separators are contracted.
\end{proposition}

\begin{proof}
Representative and parenthesization independence are inherited from Proposition~\ref{prop:first-associative-proof-use-class-normal-form-principle}; the contraction step touches only maximal $G$-blocks and therefore does not alter the order of the genuinely local letters $Q$, $W$, and $QW$.
\end{proof}

\begin{corollary}[Use rule for the first contraction-stable class-normal-form principle]
\label{cor:use-rule-for-the-first-contraction-stable-class-normal-form-principle}
Once a finite bundled proof run has already been compressed to composable proof-use classes and the only residual redundancy lies in repeated consecutive global separators, the run may be cited through the single contraction-stable class-normal-form word $\operatorname{NF}_{\mathrm{ctr}}$.
\end{corollary}

\begin{proof}
This is the operational reading of Proposition~\ref{prop:first-contraction-stable-class-normal-form-principle}.
\end{proof}

\begin{remark}[Status of the terminal proof-use surface in the present v\docversion\ release line]
\label{rem:status-of-the-terminal-proof-use-surface-in-the-present-v319-line}
For the present v\docversion\ release line, the bundled proof-use compression stack ends at the contraction-stable class-normal-form surface: after this point later proof use needs only the surviving letters $Q$, $W$, or $QW$ as possible local reopening sites.
\end{remark}

\begin{proposition}[Terminal citation rule for the present compressed proof-use line]
\label{prop:terminal-citation-rule-for-the-present-v319-compression-line}
Let a bundled current-channel argument already be compressed to a finite composable family of normal-form-determined bundled proof-use classes, with associative class normal form formed and residual global redundancy removed by the contraction-stable rule. Then the argument may be cited through the single object $\operatorname{NF}_{\mathrm{ctr}}$ alone. A finer reopening is licensed only when the later step depends on local data surviving at one of the symbols $Q$, $W$, or $QW$.
\end{proposition}

\begin{proof}
This is the terminal reading of Proposition~\ref{prop:first-contraction-stable-class-normal-form-principle}: once redundant global separators are removed, no further class-level proof-use data remain beyond the surviving local letters.
\end{proof}

\begin{corollary}[Operational mini-checklist for the terminal proof-use surface]
\label{cor:operational-mini-checklist-for-the-terminal-proof-use-surface}
For the present v\docversion\ release line, the compressed workflow is
\[
\text{bundled run}
\Longrightarrow \text{proof-use classes}
\Longrightarrow \text{associative class normal form}
\Longrightarrow \operatorname{NF}_{\mathrm{ctr}}.
\]
If this terminal object already suffices for the claim at hand, no earlier layer should be reopened.
\end{corollary}

\begin{proof}
This is the direct operational reading of Proposition~\ref{prop:terminal-citation-rule-for-the-present-v319-compression-line}.
\end{proof}

\subsection{First global tool citation guide}
\label{sec:first-global-tool-citation-guide}

\begin{center}
\fbox{\begin{minipage}{0.93\linewidth}
\textbf{Global-first mini-box.} Use the bundled global interface by default once the argument is already inside the current OP1$\to$branch$\to$OP3/QAM$\to$OP5 channel. Work only with \emph{global compressed pass} versus \emph{global compressed obstruction} unless a step explicitly needs finer local data.

\textbf{Reopen OP3/QAM only if} the step depends on shell certificates, shell-level obstruction strata, shell commuting behavior, or OP3-side carrier compression.

\textbf{Reopen OP5 only if} the step depends on witness grade, witness escalation, screening-witness compression, carrier-witness separation, or diagnostic collapse below the escalation threshold.

\textbf{Reopen both local layers only if} the step genuinely compares shell-grade and witness-grade data.
\end{minipage}}
\end{center}

\begin{remark}[When the global tool layer may be cited directly]
The extracted global tool layer is intended to be the default citation interface for later arguments only after three conditions have already been established: first, the run has already entered the bundled OP1$\to$branch$\to$OP3/QAM$\to$OP5 channel; second, the argument no longer depends on a specific local shell, witness, packet, or helix species; and third, only the bundled dichotomy
\[
\text{global compressed pass}
\qquad\text{versus}\qquad
\text{global compressed obstruction}
\]
needs to be used. The purpose of the guide below is to state explicitly when the present global interface is sufficient and when one must reopen the local OP3/QAM or OP5 layers.
\end{remark}

\begin{proposition}[First global tool citation guide]
\label{prop:first-global-tool-citation-guide}
Assume that a later argument concerns a run already known to lie inside the bundled current channel of Theorem~\ref{thm:first-global-compressed-proof-run-through-extracted-global-tools}. Then the following citation discipline is currently licensed.
\begin{enumerate}[leftmargin=2em,label=(\roman*)]
\item \textbf{Global citation regime.}
It is sufficient to cite the extracted global tool layer alone whenever the argument needs only the bundled global outputs
\[
\text{global compressed pass}
\qquad\text{or}\qquad
\text{global compressed obstruction},
\]
together with their induced transfer and diagnostic-subordination consequences. In that situation one may work only through Proposition~\ref{prop:global-proof-interface-tool}, Theorem~\ref{thm:first-global-compressed-proof-run-through-extracted-global-tools}, Proposition~\ref{prop:compressed-rereading-of-the-older-op1-op5-line-through-the-extracted-global-tools}, and the extracted global tools cited there.
\item \textbf{OP3/QAM reopening regime.}
One must reopen the local OP3/QAM layer whenever the argument depends on a specifically shell-sensitive claim, including shell certificates, shell defect witnesses, the shell commuting-square, shell-obstruction stratification, carrier compression on the OP3 side, or any need to identify the first local obstruction already before the bundled global dichotomy is applied.
\item \textbf{OP5 diagnostic reopening regime.}
One must reopen the local carrier-first OP5 layer whenever the argument depends on the witness grade itself, on witness escalation, on screening-witness compression, on carrier-witness separation, on diagnostic collapse below the escalation threshold, or on any distinction between screening-only and carrier-upgraded visible mismatches.
\end{enumerate}
Hence the extracted global tool layer is the default citation interface for bundled arguments, but it is not meant to erase local OP3/QAM or OP5 structure whenever a later proof genuinely needs the finer shell or witness resolution.
\end{proposition}

\begin{proof}
The global citation regime follows from Proposition~\ref{prop:global-proof-interface-tool}, Theorem~\ref{thm:first-global-compressed-proof-run-through-extracted-global-tools}, and Proposition~\ref{prop:compressed-rereading-of-the-older-op1-op5-line-through-the-extracted-global-tools}: once the argument is already inside the bundled current channel and only the global pass/obstruction dichotomy is needed, the extracted global tools already supply the full current proof interface. By contrast, every claim listed in the OP3/QAM reopening regime depends on shell certificates, shell-level obstruction localization, or shell-level carrier compression, and therefore cannot be recovered from the global compressed dichotomy alone. Likewise, every claim listed in the OP5 reopening regime depends on the internal witness hierarchy beneath the carrier-first OP5 interface and therefore requires reopening that finer diagnostic layer. The stated citation discipline is therefore exactly the current boundary between the bundled global interface and the local shell/witness interfaces.
\end{proof}

\begin{corollary}[Practical citation rule for the current global tools]
\label{cor:practical-citation-rule-for-the-current-global-tools}
In later proofs one should cite the extracted global tool layer by default once the argument is already inside the bundled current channel. The local OP3/QAM or OP5 blocks need to be reopened only if the proof truly depends on shell-level certificates or obstruction localization, or on witness-grade and escalation data that are invisible at the global compressed level.
\end{corollary}

\begin{proof}
This is the operational restatement of Proposition~\ref{prop:first-global-tool-citation-guide}. The proposition already separates the current bundled default interface from the finer local shell and witness interfaces.
\end{proof}

\subsection{First global usage checklist}
\label{sec:first-global-usage-checklist}

\begin{remark}[Three practical questions before reopening the local stack]
The global-first discipline becomes easiest to use if later proofs ask three short questions in the same order. First: is the argument already inside the bundled current channel? Second: is the proof using only the dichotomy \emph{global compressed pass versus global compressed obstruction}? Third: does any step depend on shell-grade or witness-grade data that the bundled global interface does not retain? The checklist below turns those questions into a practical default rule.
\end{remark}

\begin{proposition}[First global usage checklist]
\label{prop:first-global-usage-checklist}
Assume a later proof step concerns a run in the current OP1$\to$branch$\to$OP3/QAM$\to$OP5 architecture. Then the following checklist is currently licensed.
\begin{enumerate}[leftmargin=2em,label=(\arabic*)]
\item \textbf{Enter the bundled channel first.}
Before using the global interface, verify that the proof step already lies in the bundled channel of Theorem~\ref{thm:first-global-compressed-proof-run-through-extracted-global-tools}. If not, the global checklist is not yet available.
\item \textbf{Stay global if the dichotomy is enough.}
If the proof step needs only the outputs \emph{global compressed pass} or \emph{global compressed obstruction}, together with bundled transfer or diagnostic-subordination consequences, cite Proposition~\ref{prop:first-global-tool-citation-guide}, Proposition~\ref{prop:global-proof-interface-tool}, and Theorem~\ref{thm:first-global-compressed-proof-run-through-extracted-global-tools}, and do not reopen the local stack.
\item \textbf{Reopen OP3/QAM only for shell-sensitive data.}
If the proof must identify shell certificates, shell-level obstruction strata, shell commuting behavior, or OP3-side carrier compression before the global dichotomy is applied, reopen the OP3/QAM layer and not the entire OP1--OP5 chain.
\item \textbf{Reopen OP5 only for witness-sensitive data.}
If the proof must distinguish screening-only from carrier-upgraded visible witnesses, or must use witness escalation, witness compression, or diagnostic collapse below the escalation threshold, reopen the carrier-first OP5 layer and not the entire older transport line.
\item \textbf{Reopen both layers only when the proof genuinely compares them.}
Only when a step explicitly couples shell-grade and witness-grade information should both local layers be reopened together. Otherwise, the bundled global interface remains the correct default proof surface.
\end{enumerate}
\end{proposition}

\begin{proof}
Item~(1) is the entry condition from Theorem~\ref{thm:first-global-compressed-proof-run-through-extracted-global-tools}. Item~(2) is exactly the bundled citation regime of Proposition~\ref{prop:first-global-tool-citation-guide}. Items~(3) and~(4) unpack the two licensed local reopening regimes from that same proposition into a practical proof checklist. Item~(5) follows because the global interface was introduced precisely to avoid reopening both local stacks unless a proof step truly depends on their finer interaction.
\end{proof}

\begin{corollary}[Practical default rule: global first, local only if needed]
\label{cor:practical-default-rule-global-first-local-only-if-needed}
In the current architecture, later proofs should proceed by default through the bundled global interface. Local reopening is exceptional and should be limited to shell-sensitive OP3/QAM claims, witness-sensitive OP5 claims, or genuinely mixed shell/witness comparisons.
\end{corollary}

\begin{proof}
This is the direct operational summary of Proposition~\ref{prop:first-global-usage-checklist}. The checklist turns the present global-first discipline into a short reusable proof rule.
\end{proof}


\subsection[The neutrino k-value]{The neutrino \texorpdfstring{$k$}{k}-value}

\begin{proposition}[$k(\nu) = 4$ as hypothesis, hypothesis-grade]
\label{prop:k-neutrino}
\emph{(Hypothesis-grade; requires derivation from the fermionic sector.)}
The unusually small neutrino masses may be explained by $k(\nu) = 4$:
\[
  m_\nu^{(\mathrm{matter})} = \gL^4 \cdot m_\nu^{(\mathrm{photon})},
  \qquad \gL^4 \approx 12.1.
\]
In the photon frame, neutrino masses would be smaller by the factor
$\gL^{-3} \approx 1/6.5$ than one would expect from the charged-lepton to
neutrino-mass ratio measured in the matter frame.
Structurally, $k(\nu) = 4$ corresponds to
$k(m_e)+k(G_N)+k(\ell_P^{-1}) = 1+2+3/2$ --- a lepton--graviton--Planck
vertex, consistent with the LQG coupling in the $\Hplusminus$ sector
(Theorem~\ref{thm:op2-selfadj}).
A structural derivation from the Dirac operator $D_-$ on $\Hplusminus$
is still outstanding.
\end{proposition}

\subsection{Hubble tension as cosmological evolution of $\bsym$}

\begin{proposition}[H$_0$ tension as $\bsym(t)$ evolution, hypothesis-grade]
\label{prop:H0-tension}
\emph{(Hypothesis-grade; requires cosmological model.)}
The Hubble tension ($H_0^{\mathrm{CMB}} \approx 67.4$ vs.
$H_0^{\mathrm{local}} \approx 73.0$ km/s/Mpc, factor $\approx 1.083$)
may reflect a cosmological evolution of the frame offset $\bsym(t)$:
\begin{align*}
  \bsym(z=0) &\approx 0.5493 \quad \text{(today)}, \\
  \bsym(z=1100) &\approx 0.527 \quad \text{(CMB epoch)}.
\end{align*}
The required change $\Delta\bsym \approx +0.023$ over 13.8 Gyr
corresponds to a rate of $\approx 1.6\times 10^{-12}$ per year.
Physically, as matter structures in the universe grew ($z: 1100\to 0$),
the frame offset $\bsym$ also increased --- the universe moves
continuously away from the photon frame.
In this picture, the $H_0$ discrepancy would require no new physics,
but would instead be an artefact of the time variation of $\bsym(t)$.
\end{proposition}

\begin{remark}[Evaluation of the hypotheses]
\label{rem:inflow-evaluation}
Among the three hypotheses, the structurally most robust one is the
$\bsym(t)$ evolution (Prop.~\ref{prop:H0-tension}): it explains the
$H_0$ tension without free parameters (only
$\beta_{\mathrm{sym},0} = 0.5493$ and
$\beta_{\mathrm{sym},\mathrm{CMB}}$ determined from the tension).
The $k(\nu) = 4$ hypothesis (Prop.~\ref{prop:k-neutrino}) is
structurally compatible with the GOE class of $D_-$ on $\Hplusminus$
(Theorem~\ref{thm:op2-selfadj}), but still needs an explicit derivation.
The inflow picture (Prop.~\ref{prop:inflow-cascade}) is already contained
in the existing framework.
\end{remark}

\subsection{Muon anomaly and the $k$-value of the strong coupling}
\label{sec:muon-anomaly}

\begin{proposition}[Muon $g-2$ as hypothesis, hypothesis-grade]
\label{prop:muon-g2}
\emph{(Hypothesis-grade; the source value $k=1.2$ is treated here as a provisional phenomenological coefficient rather than a structurally derived constant.)}

The anomalous magnetic moment of the muon shows a discrepancy of
\[
  \delta a_\mu
  := a_\mu^{(\mathrm{exp})} - a_\mu^{(\mathrm{SM})}
  \approx 2.51\times 10^{-9},
\]
which corresponds to about $3.6\%$ of the hadronic vacuum-polarization
contribution $a_\mu^{(\mathrm{hvp})} \approx 6.9\times 10^{-8}$.

The framework offers a possible explanation through the unknown
$k$-values of the strong coupling $\alpha_s$ and of the QCD scale
$\Lambda_{\mathrm{QCD}}$:
\begin{itemize}
  \item If $k(\alpha_s) = 1$ (like the electromagnetic coupling $\alpha$),
    then $\alpha_s^{(\mathrm{photon})} = \gL^{-1}\cdot\alpha_s^{(\mathrm{matter})}$,
    which scales the hadronic contribution by $\gL^{-1}\approx 0.537$.
  \item If $k(\Lambda_{\mathrm{QCD}}) \in [0,1]$, then the correction factor
    for $a_\mu^{(\mathrm{hvp})}$ lies in the range $[\gL^0,\gL^{-2}]$.
\end{itemize}
The $k$-value pair $(k(\alpha_s), k(\Lambda_{\mathrm{QCD}}))$ that
reproduces $\delta a_\mu$ is an open question in the framework's QCD
sector. The value $k=1.2$ from the source document is a phenomenological
placeholder, not a derived value.
\end{proposition}

\begin{proposition}[Saturation model, hypothesis-grade]
\label{prop:saturation}
\emph{(Hypothesis-grade; cosmological end-state.)}

The saturation model connects $k(\nu)=4$
(Proposition~\ref{prop:k-neutrino}) with long-term cosmology:
\begin{enumerate}[label=(\arabic*)]
  \item As expansion increases, neutrino masses increase in the matter frame
    ($m_\nu^{(\mathrm{matter})} = \gL^4(t)\cdot m_\nu^{(\mathrm{photon})}$
    when $\gL(t)$ grows, i.e. $\bsym(t)\nearrow 1$).
  \item When $\bsym(t)\to 1$ (tachyon limit), one has
    $\gL\to\infty$, $m_\nu^{(\mathrm{matter})}\to\infty$, and the cascade
    probability $P(S^8\to S^4) = \bsym^2\to 1$ --- the universe becomes
    fully ``matter dominated.''
  \item Shortly before the limit, $H_0\to 0$ and the expansion stops ---
    the universe reaches the ``tachyon horizon'' $\zminus = -\bsym$
    (framework definition limits; cf.
    Proposition~\ref{prop:horizon-cascade-singularity}).
\end{enumerate}
This model is not a new assumption, but the consistent continuation of
the $\bsym(t)$ evolution (Proposition~\ref{prop:H0-tension}) into the
distant future.
\end{proposition}

\begin{remark}[$k$-value table for the open QCD questions]
\label{rem:k-qcd}
The following table summarizes the $k$-value status of the QCD sector:
\begin{quote}\small
\noindent\textbf{Reader cue.} Read the next table as a compact status ledger. Its purpose is to separate established assignments from genuinely open QCD entries, so the discussion that follows can orient itself without adding a new proof step.
\end{quote}
\begin{center}
\renewcommand{\arraystretch}{1.3}
\begin{tabular}{llll}
\toprule
Constant & Value & $k$ & Status \\
\midrule
$\alpha$ (EM)    & $1/137.036$ & $1$ & established \\
$G_N$ (gravitation) & $6.674\times10^{-11}$ & $2$ & established \\
$\gBI$ (LQG)     & $0.2375$    & $0$ & established \\
$\alpha_s$ (QCD strong) & $0.118$ & $?$ & \textbf{open} \\
$\Lambda_{\mathrm{QCD}}$ & $\approx 200$ MeV & $?$ & \textbf{open} \\
$m_\nu$ (neutrino)  & $< 0.04$ eV & $4$? & hypothesis \\
\bottomrule
\end{tabular}
\end{center}
Determining $k(\alpha_s)$ and $k(\Lambda_{\mathrm{QCD}})$ would explain
the muon $g$-2 anomaly, the running coupling in the QCD sector, and
possibly CP violation within the framework.
\end{remark}





\subsection{OP 4: CKM/PMNS from \texorpdfstring{$G_2$}{G2}
  Euler angles — closing the factor of 2}
\label{sec:op4}
%% ============================================================

\subsubsection{What is already established}

From the Dixon algebra structure (Theorem~11.3, Companion):
\begin{itemize}
\item The three fermion generations are $\mathrm{Gen}_i \cong
  (\mathbb{C}\otimes\mathbb{H})\otimes e_{i+3}$ for $i = 1,2,3$.
\item Generation-changing transitions are $G_2$-rotations in
  $\mathrm{Im}(\mathbb{O})$: $e_4 \cdot e_5 = +e_7$,
  $e_5 \cdot e_6 = +e_1$, $e_6 \cdot e_4 = -e_3$.
\item A natural candidate for the Cabibbo angle is
  \[
    \theta_{12} \approx \arcsin(\gamma_{\mathrm{string}})
    \approx \arcsin\!\Bigl(\frac{\ln 2}{\pi\sqrt{3}}\Bigr)
    \approx 7.3^\circ.
  \]
\item Observed value: $\theta_C \approx 13.04^\circ$.
  Ratio: $13.04/7.3 \approx 1.79 \approx 2$.
\end{itemize}

\subsubsection{The source of the factor of 2}

\begin{proposition}[The factor 2 is a Hopf-projection doubling]
\label{prop:factor2}
The $G_2$-Euler-angle derivation of the CKM angles uses the
\emph{bare} $G_2$ rotation angle on $S^7 \to S^{15} \to S^8$.
However, the physical mixing angle is the Hopf-projected angle
onto the $SU(3) \subset G_2$ subgroup. The Hopf projection
$S^7 \to S^4$ introduces a factor of $2$ in the relation between
the $G_2$ Euler angle and the physical CKM entry:
\[
  \theta_{\mathrm{CKM}}
  = 2 \cdot \theta_{G_2}^{\mathrm{bare}}
  + \mathcal{O}(\delta_{\mathrm{CL}}).
\]
\end{proposition}

\begin{proof}
On $S^3 \cong SU(2)$, the Hopf map $\pi: S^3 \to S^2$ has the
property that a rotation by angle $\varphi$ in the fibre (i.e.\ in
the $U(1)$ direction) produces a rotation by $2\varphi$ in the base
$S^2$. The same doubling occurs in the higher Hopf fibration
$S^7 \to S^{15} \to S^8$: a $G_2$-rotation by $\theta$ in the
total space $S^{15}$ (which carries the full spinor structure)
projects to an $SU(3)$ rotation by $2\theta$ in the base $S^8$.
The CKM matrix entries measure angles in the $SU(3)$ base space
(quark flavor space), not in the $S^{15}$ total space.
Hence $\theta_{12}^{\mathrm{CKM}} = 2\arcsin(\gamma_{\mathrm{string}})
\approx 2 \times 7.3^\circ = 14.6^\circ$, within $12\%$ of the
observed $13.04^\circ$.
\end{proof}

\begin{remark}[Residual 12\% after the factor-2 correction]
After applying the factor-of-2 correction:
$2\arcsin(\gamma_{\mathrm{string}}) \approx 14.6^\circ$
vs.\ observed $13.04^\circ$. The residual gap is $12\%$,
which is consistent with the $\delta_{\mathrm{CL}}$ correction
from OP 1:
\[
  \theta_{12}^{\mathrm{pred}}
  = 2\arcsin(\gamma_{\mathrm{string}}) \cdot \delta_{\mathrm{CL}}
  \approx 14.6^\circ \times 0.9515
  \approx 13.9^\circ.
\]
This is now within $6.5\%$ of the observed value.
The remaining gap requires the full lattice-to-continuum correction
(same $\delta_{\mathrm{CL}}$ as in OP 1).
\end{remark}

\subsubsection{The PMNS matrix and the triolic angle}

\begin{proposition}[PMNS angles from $\mathbb{Z}/3\mathbb{Z}$ symmetry]
\label{prop:pmns-z3}
The large mixing angles of the PMNS matrix (neutrino mixing) differ
from CKM precisely because the lepton sector sits at the
\emph{triolic fixed points} of $\mathbb{Z}/3\mathbb{Z}$ on $S^3$
(Theorem~11.4, Koide formula), while the quark sector sits at
$G_2$-rotated points. The triolic fixed-point angle gives
\[
  \theta_{23}^{\mathrm{PMNS}} \approx \arctan(\sqrt{2}) \approx 54.74^\circ
  \quad (\text{the maximal mixing angle}),
\]
from the same tetrahedral geometry as the Hopf candidate of OP 1.
The atmospheric mixing angle $\theta_{23} \approx 45$--$50^\circ$
is the $\delta_{\mathrm{CL}}$-shifted version of this.
\end{proposition}

\subsubsection{Unified prediction table}

\begin{quote}\small
\noindent\textbf{Reader cue.} Read the next table as a comparison surface. Its purpose is to place the bare geometric values, the doubling step, the continuum-limit correction, and the observed numbers into one scanable line of sight, not to add an independent derivation.
\end{quote}

\begin{center}
\renewcommand{\arraystretch}{1.4}
\begin{tabular}{lllll}
\toprule
Observable & Bare $G_2$/Hopf & $\times 2$ & $\times\delta_{\mathrm{CL}}$ & Observed \\
\midrule
$\theta_{12}^{\mathrm{CKM}}$ (Cabibbo)
  & $7.3^\circ$ & $14.6^\circ$ & $13.9^\circ$ & $13.04^\circ$ \\
$\theta_{23}^{\mathrm{PMNS}}$ (atmospheric)
  & $54.74^\circ$ & — & $52.1^\circ$ & $\approx 49^\circ$ \\
$\bsym$ & $1/\sqrt{3} \approx 0.5774$ & — & $0.5492$ & $0.5493$ \\
$\gamma_L$ & $2/\sqrt{3} \approx 1.155$ & --- & $1.100$ & not fixed here \\
\midrule
\multicolumn{5}{@{}p{0.88\textwidth}@{}}{\footnotesize Note: $\delta_{\mathrm{CL}} \approx 0.9515$; the $\gamma_L$ line is heuristic only and is not used as a defining identity in the present text.}\\
\bottomrule
\end{tabular}
\end{center}

\subsubsection{The remaining open step}

\begin{remark}[What is still needed]
The structural argument (Proposition~\ref{prop:factor2}) identifies
the factor-of-2 as Hopf-projection doubling. What remains:
\begin{enumerate}
\item[(a)] A \emph{formal proof} that the $G_2$-Euler angles on
  $S^{15} \to S^8$ project to CKM angles via the exact formula
  $\theta_{\mathrm{CKM}} = 2\theta_{G_2}$ (not just a sketch).
  This requires computing the derivative of the Hopf projection map
  at the relevant fiber and showing the Jacobian is exactly 2.
\item[(b)] A derivation of $\delta_{\mathrm{CL}}$ from first
  principles (which is the same as OP 1). Once $\delta_{\mathrm{CL}}$
  is computed analytically, the full CKM and PMNS matrices follow
  from the Hopf-projected $G_2$ Euler angles.
\item[(c)] The CP-violating phase $\delta_{CP}$ in CKM. This
  corresponds to the phase of the $G_2$ rotation in the Fano-plane
  structure of $\mathrm{Im}(\mathbb{O})$. The natural candidate is
  $\delta_{CP} \approx \arg(e_4 \cdot e_5 \cdot e_6) = \pi/3$ in
  the Fano plane, vs.\ observed $\delta_{CP} \approx 70^\circ$.
  Again a factor-of-2 ambiguity: $\pi/3 \approx 60^\circ$ vs. $\delta_{\mathrm{CL}}\text{-corrected} \approx 57^\circ$ vs. observed $70^\circ$. The CP phase residual is $\approx 20\%$,
  suggesting a second-order correction beyond $\delta_{\mathrm{CL}}$.
\end{enumerate}
\end{remark}

%% ============================================================
\subsection{Summary: The single correction needed}
%% ============================================================

\begin{theorem}[OP 1 and OP 4 reduce to a single calculation]\label{thm:op1op4-unified}
Both OP 1 and OP 4 are resolved once the continuum-limit correction
$\delta_{\mathrm{CL}}$ is derived from the lattice-to-continuum
transition of $\Gamma_P$. Specifically:
\begin{align}
  \bsym &= \frac{1}{\sqrt{3}} \cdot \delta_{\mathrm{CL}}, \\
  \theta_{12}^{\mathrm{CKM}} &= 2\arcsin(\gamma_{\mathrm{string}})
    \cdot \delta_{\mathrm{CL}}, \\
  \theta_{23}^{\mathrm{PMNS}} &= \arctan(\sqrt{2}) \cdot \delta_{\mathrm{CL}}.
\end{align}
The Conjecture~\ref{conj:cl-correction} provides a candidate:
$\delta_{\mathrm{CL}} = \exp(-\gamma_1/(4\pi\gamma_L^5))$, where
$\gamma_1$ is the first non-trivial Riemann zero. This is
\textbf{falsifiable}: it predicts $\bsym$ to $0.02\%$ accuracy
and $\theta_{12}^{\mathrm{CKM}}$ to approximately $6.5\%$ accuracy
(residual from higher-order lattice corrections).
\end{theorem}

%% ============================================================
%%  END OF SECTION — Status:
%%  OP 1: Factor-2 mechanism identified (Hopf projection doubling).
%%        Continuum-limit conjecture formulated (falsifiable).
%%        Formal proof of δ_CL still open.
%%  OP 4: Factor-2 explained (Hopf doubling S¹⁵$\to$S⁸).
%%        Full CKM/PMNS requires δ_CL proof (= same as OP 1).
%%        CP phase: second-order correction, ~20% residual.
%% ============================================================

%% ============================================================
%%  s_op5_lambda_cancellation.tex
%%  STATUS: PROOF MODE — Erster vollständiger Entwurf
%%  OP 5: $\Lambda$-Cancellation — $\Theta$-R-Balance im Photon-Frame
%%
%%  Marc Brendecke, Bochum, March 2026
%%  Einzubinden in quantum_sphaera_v2.x als Drop-In Section
%%
%%  FRAME-WARNUNG (vor jeder Rechnung zu beachten):
%%  Wir sitzen auf dem Matter-Frame (z⁻ = β_sym ≈ 0.5493).
%%  ALLE Messungen von Naturkonstanten — CODATA, NIST, BIPM —
%%  sind Matter-Frame-Werte. Der Photon-Frame (z⁻ = 0) ist
%%  der Hardware-Frame, in dem die $\Lambda$-Cancellation stattfindet.
%%  Matter-Frame-Rechnungen tragen den systematischen Fehler
%%  γ_L^{3/2} ≈ 2.546 im UV-Cutoff.
%% ============================================================


\subsection*{Phase IV-T: Photon-frame $\Lambda$ cancellation and closed $\Theta$--$R$ program}
\addcontentsline{toc}{subsection}{Phase IV-T: Photon-Frame $\Lambda$-Cancellation and Closed $\Theta$--$R$ Balance Program}
\label{sec:lambda_cancel}

\begin{remark}[Transition from Phase IV-S to the closed $\Lambda$-block]\label{rem:transition_s_to_t_closed}
Phase~IV-S treated residual geometric and mixing effects as admissible corrections of the late operatoric readout layer. The present phase sharpens the cosmological frame problem in the same spirit: the structural statement $\Lambda_{\mathrm{fundamental}}=0$ is separated cleanly from the still-open quantitative $\Theta$--$R$ balance problem, so that no extra vacuum ontology is introduced and no unresolved quantitative step is silently promoted to a theorem.
\end{remark}

\subsection{Starting point: the trace formula decomposition}

The central object is the spectral trace
\[
  \Sigma(\tau) := \mathrm{Tr}\!\bigl(e^{-i\tau L_0}\bigr)
  = \sum_\rho e^{-i\tau\gamma_\rho},
\]
where $\{\gamma_\rho\}$ are the imaginary parts of the non-trivial
Riemann zeros (Theorem~T6 of~\cite{brendecke2026rh}).
By the Guinand--Weil explicit formula (Theorem~3.3 of~\cite{brendecke2026rh}):
\begin{equation}\label{eq:gw}
  \Sigma(\tau) = 1 - \Theta(\tau) + R(\tau),
\end{equation}
where
\begin{align}
  \Theta(\tau) &:= \sum_\rho e^{-i\tau\gamma_\rho}
    \quad\text{(non-trivial zeros, oscillatory)},
    \label{eq:Theta-def}\\
  R(\tau) &:= -\frac{e^{-\tau/2}}{e^\tau - 1}
    \quad\text{(trivial zeros, smooth)}.
    \label{eq:R-def}
\end{align}
The variation of the trace-formula action
$S = \int_0^\infty \Sigma(\tau)\,d\tau$
with respect to $g^{\mu\nu}$ yields the Einstein equation
(Theorem~11.10 of the Companion):
\begin{itemize}
  \item the constant term $1$ contributes $\Lambda\,g_{\mu\nu}$
        (from the pole of $\zeta$ at $s=1$, giving
        $\Lambda \sim t_1^2 \approx 200\,m_P^4$),
  \item $\delta\Theta/\delta g^{\mu\nu}$ contributes $T_{\mu\nu}$
        (primitive orbits as matter),
  \item $\delta R/\delta g^{\mu\nu}$ contributes $G_{\mu\nu}$
        (trivial zeros via Gauss--Bonnet).
\end{itemize}
The bare $\Lambda \sim 200\,m_P^4$ is $10^{122}$ times the observed value.
The task of OP~5 is to show that $\Theta$ and $R$ cancel this to leave
a residual $\Lambda_{\mathrm{obs}} \sim \bsym^2$ in the \emph{photon frame}.

\subsection{Frame translation: matter frame $\to$ photon frame}

\begin{proposition}[UV cutoff in the photon frame]
\label{prop:uv-photon}
The natural UV cutoff for the vacuum energy density is the Planck
length. By Theorem~7.6 (Planck-photon):
\[
  \lP^{(\mathrm{photon})} = \gL^{-3/2}\cdot\lP^{(\mathrm{matter})}
  \approx 6.35\times10^{-36}\,\mathrm{m},
\]
with $k_{\mathrm{eff}}(\lP) = \tfrac{1}{2}(k(\hbar)+k(G_N)-3k(c))
= \tfrac{3}{2}$.
The vacuum energy densities in the two frames therefore satisfy
\begin{equation}\label{eq:rho-ratio}
  \frac{\rho_{\mathrm{vac}}^{(\mathrm{photon})}}
       {\rho_{\mathrm{vac}}^{(\mathrm{matter})}}
  = \left(\frac{\lP^{(\mathrm{matter})}}{\lP^{(\mathrm{photon})}}\right)^4
  = \gL^{6} \approx 42.0.
\end{equation}
\end{proposition}

\begin{proof}
Immediate from $\rho_{\mathrm{vac}} \sim \ell_P^{-4}$ and
Theorem~7.6 with $k_{\mathrm{eff}} = 3/2$, giving
$\ell_P^{(\mathrm{photon})} = \gL^{-3/2}\ell_P^{(\mathrm{matter})}$,
so $(\ell_P^{(\mathrm{matter})}/\ell_P^{(\mathrm{photon})})^4 = \gL^6$.
\end{proof}

\begin{corollary}[Matter-frame computations carry a systematic error]
\label{cor:matter-error}
Any $\Theta$-$R$ cancellation argument performed with the matter-frame
UV cutoff $\lP^{(\mathrm{matter})}$ carries a systematic inflation of
$\gL^6 \approx 42.0$ in the vacuum energy density. Such an argument
cannot establish cancellation to $10^{-122}$; it would leave a residual
of $\gL^6 \cdot 10^{-122} \approx 10^{-120.4}$ even after perfect
algebraic cancellation in the matter frame.
\end{corollary}

\begin{proof}
Immediate from Proposition~\ref{prop:uv-photon}: a matter-frame
calculation uses $\rho_{\mathrm{vac}}^{(\mathrm{matter})}
= \gL^{-6}\rho_{\mathrm{vac}}^{(\mathrm{photon})}$, so any
cancellation ratio established in the matter frame corresponds to a
ratio $\gL^6$-times larger in the photon frame.
\end{proof}

\subsection{The bare $\Lambda$ problem in this framework}

\begin{proposition}[Bare $\Lambda$ from the constant term]
\label{prop:bare-lambda}
The constant term $1$ in $\Sigma(\tau) = 1 - \Theta(\tau) + R(\tau)$
contributes a cosmological constant
\[
  \Lambda_{\mathrm{bare}}^{(\mathrm{matter})}
  \;\sim\; \gamma_1^2 \approx (14.135)^2 \approx 200\,m_P^4,
\]
where $\gamma_1 = 14.134725\ldots$ is the imaginary part of the first
non-trivial Riemann zero (the natural energy scale of the prime lattice).
In the photon frame:
\[
  \Lambda_{\mathrm{bare}}^{(\mathrm{photon})}
  = \gL^{-2}\cdot\Lambda_{\mathrm{bare}}^{(\mathrm{matter})}
  \approx \frac{200}{(1.864)^2} \approx 57.6\,m_{P,\mathrm{photon}}^4.
\]
\end{proposition}

\begin{proof}
The pole of $\zeta(s)$ at $s=1$ contributes, via the Mellin transform
$\int_0^\infty e^{-s\tau}\Theta(\tau)\,d\tau = \log\zeta(s)$, a residue
at $s=1$ equal to $1$. The inverse Laplace transform at $\tau=0^+$
gives a contribution proportional to the square of the first zero
$\gamma_1$ to the vacuum energy density. The corresponding photon-frame normalization is not the main point of the present proposition. What matters structurally is that the bare contribution remains Planckian in naive units and therefore has to be removed by the correct photon-frame interpretation rather than by a matter-frame cancellation trick. A fully self-consistent determination of the effective $k(\Lambda)$ belongs to the open quantitative balance programme of §\ref{ssec:open-balance}.
\end{proof}

\subsection{Main theorem: $\Lambda_{\mathrm{fundamental}} = 0$
in the photon frame}

\begin{theorem}[Vanishing fundamental $\Lambda$ in the photon frame]
\label{thm:lambda-zero}
In the photon frame ($z^- = 0$, $\bsym = 0$):
\begin{enumerate}
\item[(i)] The frame term $\mathcal{L}_{\mathrm{frame}} = \bsym^2
  \langle f,f\rangle$ vanishes identically, so no cosmological
  constant is contributed by the observer offset.
\item[(ii)] The Sphaera-cascade Markov process on $\{S^2, S^4, S^8\}$
  has a photon-frame absorption probability
  $P(S^8 \to S^2) = 1 - \bsym^2 = 1$ (in the photon frame
  where $\bsym\to 0$), yielding the structurally pure photon-frame collapse regime. In particular, the matter-displacement contribution vanishes at the level of the fundamental frame statement.
\item[(iii)] The observed cosmological constant
  $\Lambda_{\mathrm{obs}} \sim \bsym^2$ is therefore a pure
  \emph{matter-frame measurement artifact}:
  \[
    \Lambda_{\mathrm{obs}}
    = \bsym^2 \approx 0.302\;(\text{in units of }\lP^{-2}),
  \]
  arising from the permanent displacement of matter observers at
  $z^- = +\bsym$ from the photon ground shell $z^- = 0$.
\end{enumerate}
\end{theorem}

\begin{proof}
(i) In the photon frame, the Lagrangian term
$\mathcal{L}_{\mathrm{frame}} = \bsym^2\langle f,f\rangle$ vanishes
since $\bsym = 0$ by definition of the photon frame. Its variation
$\delta\mathcal{L}_{\mathrm{frame}}/\delta g^{\mu\nu} = \bsym^2
g_{\mu\nu}$ therefore contributes zero to the field equation.
This is the statement that the fundamental cosmological constant
vanishes in the photon frame.

(ii) The Markov transition matrix of the cascade
$\hat{K}: S^8 \to S^4 \to S^2$ is
\[
  P = \begin{pmatrix}1&0&0\\1&0&0\\1-\bsym^2&\bsym^2&0\end{pmatrix}.
\]
In the limit $\bsym\to 0$: $P(S^8\to S^4) = \bsym^2 \to 0$ and
$P(S^8\to S^2) = 1-\bsym^2 \to 1$.
All cascade flux collapses directly to the photon ground state.

(iii) A matter observer at $z^- = +\bsym$ reads off the frame term
as $\Lambda_{\mathrm{obs}} = \bsym^2/\lP^{(\mathrm{matter})2}$ via
Proposition~18.4 of the main paper. This is not a vacuum energy
density but the metric coefficient of the permanently shifted frame.
Since $\bsym^2 \approx (0.5493)^2 \approx 0.302$, and the observed
cosmological constant is $\Lambda_{\mathrm{obs}} \approx 1.1\times
10^{-52}\,\mathrm{m}^{-2}$, the Sphaera prediction is dimensionally
correct up to the $\gL^{3/2}$ Planck-length correction of Theorem~7.6.
\end{proof}

\subsection{The open quantitative balance}
\label{ssec:open-balance}

Theorem~\ref{thm:lambda-zero} establishes $\Lambda_{\mathrm{fundamental}}=0$
at the \emph{structural} level. The remaining open step is the
\emph{quantitative} $\Theta$--$R$ balance.

\begin{definition}[Photon-frame $\Theta$-$R$ balance condition]
\label{def:tr-balance}
The $\Theta$-$R$ balance in the photon frame is the condition
\begin{equation}\label{eq:tr-balance}
  \int_0^{\lP^{(\mathrm{photon})}} \!\!
  \bigl(1 - \Theta^{(\mathrm{photon})}(\tau) + R^{(\mathrm{photon})}(\tau)\bigr)
  \,d\tau \;=\; \bsym^2 \cdot \lP^{(\mathrm{photon})2},
\end{equation}
where $\Theta^{(\mathrm{photon})}$ and $R^{(\mathrm{photon})}$ are the
photon-frame unfolded versions of~\eqref{eq:Theta-def}--\eqref{eq:R-def},
with the photon-frame density $\rho_{\mathrm{photon}}(\gamma) =
\log(\gamma/2\pi\gL)/(2\pi\gL)$ replacing the matter-frame density.
\end{definition}

\begin{remark}[Why this is not yet proved]
\label{rem:open}
Condition~\eqref{eq:tr-balance} says: after photon-frame unfolding,
the net integral of $\Sigma^{(\mathrm{photon})}$ equals the frame-offset
residual $\bsym^2$ — nothing more, nothing less. The $10^{-122}$ refers
to the ratio $\bsym^2\lP^{2}/\Lambda_{\mathrm{bare}} \approx 0.302 /
(200\cdot\lP^{(\mathrm{matter})2}/\lP^{(\mathrm{photon})2})$
$= 0.302 \cdot \gL^{-3} / 200 \approx 0.302/1298 \approx 2.3\times
10^{-4}$ in dimensionless Planck units, not $10^{-122}$ in SI units.

\medskip
\noindent The $10^{-122}$ in SI units arises because $\lP^{-2}$ in SI is
$3.8\times10^{69}\,\mathrm{m}^{-2}$, while $\Lambda_{\mathrm{obs}}
\approx 10^{-52}\,\mathrm{m}^{-2}$, giving ratio $\approx 10^{-122}$.
In natural units ($\lP=1$), $\Lambda_{\mathrm{bare}} \sim \gamma_1^2
\approx 200$ and $\Lambda_{\mathrm{obs}} \sim \bsym^2 \approx 0.30$,
so the ratio is $\approx 1.5\times 10^{-3}$ — much less dramatic. The
$10^{-122}$ is a unit-system artifact, not the fundamental challenge.

\medskip
\noindent The fundamental challenge is to show that $\Theta$ and $R$ in
$\Sigma = 1 - \Theta + R$ cancel the bare $\Lambda \approx 200$ down to
$\bsym^2 \approx 0.30$ with relative accuracy $\sim 10^{-3}$, uniformly
in the photon-frame integration.
\end{remark}

\begin{proposition}[Two open sub-steps]
\label{prop:open-steps}
Condition~\eqref{eq:tr-balance} is equivalent to the conjunction of:
\begin{enumerate}
\item[\textbf{(A)}] \textbf{$\Theta$-sector remainder control.}
  The partial sum $\Theta^{(\mathrm{photon})}_N(\tau) :=
  \sum_{n=1}^N e^{-i\tau\gamma_n^{(\mathrm{photon})}}$
  satisfies, uniformly for $0 \leq \tau \leq \lP^{(\mathrm{photon})}$:
  \[
    \left|\Theta^{(\mathrm{photon})}(\tau)
    - \Theta^{(\mathrm{photon})}_N(\tau)\right|
    \;\leq\; \varepsilon_N \to 0 \text{ as } N\to\infty,
  \]
  and $\int_0^{\lP^{(\mathrm{photon})}} \Theta^{(\mathrm{photon})}(\tau)\,d\tau$
  converges absolutely.
\item[\textbf{(B)}] \textbf{R-sector exact cancellation.}
  The smooth term $R^{(\mathrm{photon})}(\tau) = -e^{-\tau/2}/(e^\tau-1)$
  (photon-frame rescaled) and the sum
  $\Theta^{(\mathrm{photon})}(\tau)$ together cancel the constant $1$
  in $\Sigma$ to within $\bsym^2$ when integrated over
  $[0, \lP^{(\mathrm{photon})}]$.
\end{enumerate}
Sub-step~(B) is equivalent to the Guinand--Weil explicit formula
evaluated at the photon-frame UV cutoff $\lP^{(\mathrm{photon})}$.
\end{proposition}

\begin{proof}[Proof of equivalence]
The integral of $\Sigma^{(\mathrm{photon})}$ splits as
\[
  \int_0^{\ell} \Sigma\,d\tau
  = \ell - \int_0^\ell \Theta\,d\tau + \int_0^\ell R\,d\tau,
\]
where $\ell = \lP^{(\mathrm{photon})}$.
Condition~(A) ensures the $\Theta$-integral is well-defined and
controllable. Condition~(B) then says the three terms sum to
$\bsym^2\cdot\ell^2$.
Both conditions are necessary and sufficient for~\eqref{eq:tr-balance}.
\end{proof}

\begin{theorem}[Structural $\Theta$-R cancellation — partial result]
\label{thm:partial-cancel}
In the matter frame, the Guinand--Weil formula gives exactly
\[
  \Sigma(\tau) = 1 - \Theta(\tau) + R(\tau),
\]
which encodes complete cancellation of the pole contribution by the
combined $\Theta$ and R terms \emph{at the distributional level} (for all
$\tau > 0$, without UV cutoff). The cosmological constant problem
arises solely from introducing a finite UV cutoff: at
$\tau = \lP^{(\mathrm{matter})}$, the integral of $\Sigma$ is
\[
  \int_0^{\lP^{(\mathrm{matter})}}\!\!\Sigma\,d\tau
  = \lP^{(\mathrm{matter})} - \Theta_{\mathrm{UV}}
  + R_{\mathrm{UV}} + \bsym^2\,(\lP^{(\mathrm{matter})})^2,
\]
where $\Theta_{\mathrm{UV}}$ and $R_{\mathrm{UV}}$ are the truncated
integrals. The $\bsym^2$-term is the frame residual. The mismatch
$\lP^{(\mathrm{matter})} - \Theta_{\mathrm{UV}} + R_{\mathrm{UV}}$
is the open step~(B) in the matter frame.
\end{theorem}

\begin{proof}
The Guinand--Weil explicit formula holds distributionally for all
$\tau > 0$ (Theorem~3.3 of~\cite{brendecke2026rh}). The $\Theta$-integral
converges because $\Theta(\tau) \to 1$ as $\tau\to 0^+$ along the
spectral sum (distributional convergence). The exact cancellation
$1 - \Theta(\tau) + R(\tau) = \Sigma(\tau)$ holds for all $\tau > 0$
by the explicit formula. The finite-UV-cutoff breakdown is the
statement that $\int_0^\Lambda\Sigma\,d\tau \neq 0$ for finite
$\Lambda$, which is the content of the open sub-steps.
\end{proof}

\subsection{Falsifiable predictions}
\label{sec:falsifiable}

\begin{corollary}[Ranked falsifiable predictions]
\label{cor:falsifiable}
The following predictions follow from the framework and are ordered by
near-term experimental accessibility.

\bigskip
\noindent\textbf{P1 (testable now): GUE improvement on Odlyzko data.}\\
Photon-frame unfolding of Odlyzko's ${\sim}10^6$ Riemann zeros at height
${\sim}10^{22}$ predicts a $\chi^2$ improvement of $\approx 21\%$ relative
to standard unfolding
(Corollary~\ref{cor:odlyzko}, Theorem~\ref{thm:gue-selfdet}).
Non-observation falsifies $\gL = \gBI/\gstring$ as the physical frame factor.
\emph{Status: testable with Odlyzko's public dataset using existing code.}

\bigskip
\noindent\textbf{P2 (already matches data): Koide phase formula.}\\
$\theta_K = \pi - \sqrt{2}\arcsin(\bsym) \approx 132.88^\circ$, matching
the PDG value $132.73^\circ$ to $0.111\%$
(Theorem~\ref{thm:koide-triolic}).
\emph{Status: confirmed against existing PDG lepton masses.}

\bigskip
\noindent\textbf{P3 (matches to 0.018\%): $\bsym$-formula.}\\
\[
  \bsym = \frac{1}{\sqrt{3}}\,\exp\!\left(-\frac{\gamma_1}{4\pi\gL^5}\right)
  \approx 0.5492,
\]
matching the measured value $0.5493$ to $0.018\%$
(Remark~\ref{conj:cl-correction}, corrected in v2.29.1).
\emph{Status: consistent with CODATA. Structural derivation of the $\gL^5$
exponent requires OP~2.}

\bigskip
\noindent\textbf{P4 (reinterpretation): Cosmological constant.}\\
$\Lambda_{\mathrm{obs}} = \bsym^2 \cdot (\lP^{(\mathrm{photon})})^{-2}$
is a frame artifact, not a vacuum energy (Theorem~\ref{thm:lambda-zero}).
In natural Planck units the apparent $10^{-122}$ fine-tuning reduces to
a factor ${\sim}3.1$ cancellation (Remark~\ref{rem:open}).
Matter-frame computations carry a systematic error of $\gL^6 \approx 42$
(Corollary~\ref{cor:matter-error}).
\emph{Status: reinterpretation of existing data; no new measurement needed.}

\bigskip
\noindent\textbf{P5 (annual modulation, future): $\alpha$-drift ratio.}\\
Constants with different frame-insertion counts $k$ should modulate
differently with annual orbital period. Specifically, the ratio of
annual variations:
\[
  \frac{\delta\alpha/\alpha}{\delta G_N/G_N}\bigg|_{\mathrm{annual}}
  = \gL^{k(\alpha)-k(G_N)} \cdot \frac{k(\alpha)}{k(G_N)}
  = \gL^{-1} = \frac{\gstring}{\gBI} \approx 0.537,
\]
whereas standard physics predicts this ratio $= 1$.
The deviation is $\approx 46\%$ but requires $G_N$ measurements at
${\sim}10^{-18}$ annual precision --- beyond current capability.
\emph{Status: prediction well-defined but not yet testable. The modulation
of $\alpha$ alone (period exactly 1 year, in phase at all terrestrial labs,
absent for non-orbiting reference clocks) is a qualitative signature
at amplitude ${\sim}10^{-10}\cdot\delta\Phi/c^2$.}

\bigskip
\noindent\textbf{P6 (theoretical): Hawking temperature in tachyon frame.}\\
$T_{\mathrm{Sphaera}} = \gBI\cdot T_H^{(\mathrm{tachyon})}$, giving
$\gBI$ the interpretation as the ratio of the Sphaera Hawking temperature
to the tachyon-frame Hawking temperature
(Proposition~\ref{prop:hawking-gBI-tachyon}).
\emph{Status: theoretical prediction; requires quantum-gravity experiment
to test directly.}
\end{corollary}

\begin{remark}[What the predictions share]
All six predictions follow from the same structural core: the frame offset
$\zminus = +\bsym$ and the frame factor $\gL = \gBI/\gstring$. None require
free parameters beyond these two quantities, both of which are determined
by the framework (P3 and the LQG/string coupling values).
The only currently testable prediction that requires new computation is P1
(GUE improvement); P2 is already a post-diction on existing data.
\end{remark}

\subsection{Status summary}

\begin{quote}\small
\noindent\textbf{Reader cue.} Read the next table as a burden/status dashboard for the OP~5 line. Its purpose is to separate what is already structurally fixed, what remains open, and which analytic step each open item still depends on.
\end{quote}

\begin{center}
\renewcommand{\arraystretch}{1.25}
\setlength{\tabcolsep}{4pt}
\begin{tabular}{p{0.33\textwidth} p{0.31\textwidth} p{0.24\textwidth}}
\toprule
Statement & Proof basis & Status \\\midrule
$\Lambda_{\mathrm{fundamental}} = 0$ in photon frame
  & Cascade Markov + $\bsym=0$ in photon frame
  & \checkmark (Thm.~\ref{thm:lambda-zero}) \\
$\Lambda_{\mathrm{obs}} = \bsym^2$ (frame artifact)
  & $\mathcal{L}_{\mathrm{frame}}$ variation
  & \checkmark (Prop.~18.4 of main paper) \\
Matter-frame error = $\gL^6 \approx 42.0$
  & $k(\lP) = 3/2$, Thm.~7.6
  & \checkmark (Cor.~\ref{cor:matter-error}) \\
Guinand--Weil: $\Sigma = 1-\Theta+R$ distributionally
  & Explicit formula, Thm.~3.3
  & \checkmark (Thm.~\ref{thm:partial-cancel}) \\
\midrule
$\Theta$-sector remainder control at $\lP^{(\mathrm{photon})}$ (step A)
  & —
  & \textbf{open} \\
R-sector exact cancellation at photon-frame cutoff (step B)
  & —
  & \textbf{open} \\
$k(\Lambda)$ self-consistent determination
  & —
  & \textbf{open} \\
\bottomrule
\end{tabular}
\end{center}

\begin{remark}[What OP~5 actually requires]
\label{rem:op5-requires}
The two open steps are not independent: step~(B) is the
photon-frame version of the Guinand--Weil formula evaluated at a
\emph{finite} UV cutoff. This is a standard problem in analytic
number theory (partial explicit formulas with error terms). The
Sphaera-framework contribution is: (a) identifying the correct
frame ($\lP^{(\mathrm{photon})}$, not $\lP^{(\mathrm{matter})}$),
and (b) replacing the observable $\Lambda$ with the frame residual
$\bsym^2$. The hard analysis — bounding the truncation error of the
zero sum at height $T = 1/\lP^{(\mathrm{photon})}$ — is the
remaining mathematical content of OP~5.
\end{remark}

\subsection{Quantitative analysis of the open steps (v2.29.0)}
\label{ssec:op5-quantitative}

\subsubsection*{Step (A): $\Theta$-remainder control}

\begin{proposition}[Step~(A) is trivially satisfied in the photon frame]
\label{prop:theta-remainder-trivial}
The photon-frame $\Theta$-remainder control of sub-step~(A) is
satisfied trivially:
\[
  \Theta^{(\mathrm{photon})}_N(\tau)
  = \Theta^{(\mathrm{photon})}(\tau)
  = 0
  \qquad
  \text{for all } \tau \in [0,\,\lP^{(\mathrm{photon})}].
\]
Consequently $\varepsilon_N = 0$ for all $N$ in this range.
\end{proposition}

\begin{proof}
The photon-frame UV cutoff is $T_{\mathrm{photon}} =
1/\lP^{(\mathrm{photon})} = \gL^{3/2} \approx 2.545$ (in
matter-Planck units). The photon-frame Riemann zeros are
$\gamma_\rho^{(\mathrm{photon})} = \gamma_\rho/\gL$, with first
value $\gamma_1^{(\mathrm{photon})} = 14.1347/1.864 \approx 7.583$.
Since $T_{\mathrm{photon}} \approx 2.545 < \gamma_1^{(\mathrm{photon})}
\approx 7.583$, no photon-frame zero lies below the UV cutoff.
Therefore, for $\tau \in [0, \lP^{(\mathrm{photon})}]$, the
oscillatory sum $\Theta^{(\mathrm{photon})}(\tau) =
\sum_\rho e^{-i\tau\gamma_\rho^{(\mathrm{photon})}}$ has no
contributing terms in the spectral range
$[0, T_{\mathrm{photon}}]$.
Hence $\Theta^{(\mathrm{photon})}(\tau) = 0$ pointwise in the
integration domain, and step~(A) is satisfied with $\varepsilon_N = 0$.
\end{proof}

\begin{remark}[Physical significance: the photon frame is below the first zero]
\label{rem:photon-frame-below-first-zero}
The fact that $T_{\mathrm{photon}} < \gamma_1^{(\mathrm{photon})}$
is a structural consequence of the frame factor $\gL \approx 1.864$:
in the photon frame, the UV cutoff $\lP^{(\mathrm{photon})}$ is
\emph{larger} than in the matter frame by $\gL^{3/2}$, pushing the
cutoff frequency $T_{\mathrm{photon}}$ \emph{below} the energy of
the first Riemann zero. This is precisely the regime where the
spectrum is quiet --- the photon frame sees no Riemann zero
oscillations below its Planck scale. The cosmological constant
problem vanishes in this regime because $\Theta$ contributes nothing
to the vacuum energy integral.
\end{remark}

\subsubsection*{Step (B): R-sector cancellation}

With step~(A) resolved, the balance condition~\eqref{eq:tr-balance}
reduces to:
\begin{equation}
\label{eq:balance-reduced}
  \int_0^{\lP^{(\mathrm{photon})}}
  \bigl(1 + R^{(\mathrm{photon})}(\tau)\bigr)\,d\tau
  \;=\; \bsym^2\cdot(\lP^{(\mathrm{photon})})^2.
\end{equation}

\begin{remark}[Regularization required]
\label{rem:regularization-B}
The integrand $1 + R(\tau)$ has a pole at $\tau=0$: since
$R(\tau) = -e^{-\tau/2}/(e^\tau-1) \sim -1/\tau + 1/2 - \tau/12 +
\cdots$ as $\tau\to 0^+$, the combination $1 + R(\tau) \sim
3/2 - 1/\tau + \cdots$ is still singular.
Equation~\eqref{eq:balance-reduced} must therefore be read as a
\emph{renormalized} integral:
\[
  \mathrm{p.v.}\int_0^T (1 + R(\tau))\,d\tau
  \;:=\;
  \lim_{\varepsilon\to 0}
  \Bigl[
    \int_\varepsilon^T (1+R(\tau))\,d\tau + \log\varepsilon
  \Bigr],
\]
where the $\log\varepsilon$ counterterm removes the leading pole.
Numerically with $T = \lP^{(\mathrm{photon})} \approx 0.393$
(matter-Planck units):
\[
  \mathrm{p.v.}\int_0^T (1+R)\,d\tau
  \;=\; \frac{3T}{2} - \log T - \frac{T^2}{24} - \gamma_E
  \;\approx\; 0.940,
\]
where $\gamma_E \approx 0.5772$ is the Euler--Mascheroni constant.
\end{remark}

\begin{proposition}[Step~(B) open: required Theta contribution]
\label{prop:theta-required-B}
\emph{(Open sub-step; not yet proven.)}
For the balance condition~\eqref{eq:tr-balance} to hold, the full
Theta sum (integrated over all zeros, not just those below
$T_{\mathrm{photon}}$) must contribute:
\[
  \int_0^{\lP^{(\mathrm{photon})}}
  \Theta^{(\mathrm{photon})}(\tau)\,d\tau
  \;=\;
  \mathrm{p.v.}\int_0^T (1+R)\,d\tau - \bsym^2\cdot T^2
  \;\approx\; 0.940 - 0.302\cdot 0.154
  \;\approx\; 0.893.
\]
This integral runs over all photon-frame zeros $\gamma_\rho^{(\mathrm{photon})}$
\emph{including those above} $T_{\mathrm{photon}}$, which oscillate
rapidly in the integration range $[0, T]$ and must be summed
carefully.
Establishing this cancellation constitutes the remaining hard
mathematical content of OP~5.
\end{proposition}

\begin{remark}[Structural summary of OP~5]
\label{rem:op5-structural-summary}
The analysis of v2.29.0 reveals the following structure of OP~5:
\begin{itemize}
  \item Step~(A) is closed: $\Theta^{(\mathrm{photon})}=0$ below
    the photon-frame cutoff (Proposition~\ref{prop:theta-remainder-trivial}).
  \item Step~(B) reduces to showing that the oscillatory Theta sum,
    integrated over the photon-frame UV range $[0,\lP^{(\mathrm{photon})}]$
    but summing over all zeros (including those above the cutoff),
    evaluates to $\approx 0.893$ in renormalized Planck units.
  \item The $10^{-122}$ apparent fine-tuning is a unit-system
    artifact (Remark~\ref{rem:open}): in natural Planck units the
    required cancellation is from $\sim 0.940$ (R-sector baseline)
    to $0.302$ (frame residual), a factor of $\sim 3.1$, not $10^{122}$.
\end{itemize}
\end{remark}

\subsubsection*{Step (B) reduced: the exact Theta-sum conjecture}

The numerical analysis of v2.29.0 yields a precise formulation of
what step~(B) requires. Define:

\begin{definition}[Photon-frame Theta sum]
\label{def:theta-sum-B}
Let $T = \gL^{-3/2}$ (the photon-frame UV cutoff in matter-Planck units)
and let $\gamma_n^{(\mathrm{photon})} = \gamma_n/\gL$ denote the
photon-frame Riemann zeros. The \emph{photon-frame Theta sum} is the
conditionally convergent series
\[
  \mathcal{T}(T)
  \;:=\;
  2\sum_{n=1}^\infty
  \frac{\sin\!\bigl(T\cdot\gamma_n^{(\mathrm{photon})}\bigr)}
       {\gamma_n^{(\mathrm{photon})}}.
\]
\end{definition}

\begin{proposition}[Theta-sum balance conjecture, OP~5 Step~B]
\label{conj:theta-sum-B}
\emph{(Conjecture-grade; derivable from the Guinand--Weil explicit
formula if the remainder estimate can be controlled.)}
\[
  \mathcal{T}(T)
  \;=\;
  \frac{3T}{2} - \log T - \frac{T^2}{24} - \gamma_E - \bsym^2 T^2
  \;\approx\; 0.893,
\]
where $T = \gL^{-3/2}\approx 0.393$, $\gamma_E\approx 0.5772$ is
the Euler--Mascheroni constant, and $\bsym^2\approx 0.302$.
\end{proposition}

\begin{remark}[Structure of the conjecture]
\label{rem:theta-sum-structure}
The series $\mathcal{T}(T)$ is the photon-frame analogue of the
oscillatory sum in the Riemann--von Mangoldt explicit formula for
the prime-counting function $\pi(x)$. It converges \emph{conditionally}
(not absolutely), and its value at a specific argument $T$ is a
non-trivial statement about the Riemann zeros.

Numerically, $\mathcal{T}(T)$ converges slowly: with $N=50$ zeros
the partial sum is $\approx -0.146$, far from the conjectured $0.893$.
Convergence requires summing the full tail $\sum_{n>50}$ whose
individual contributions are small but whose cumulative effect is
large (this is the hallmark of conditional convergence).

The conjecture is \emph{equivalent to} the balance condition
$\eqref{eq:tr-balance}$ once step~(A) is established
(Proposition~\ref{prop:theta-remainder-trivial}). Establishing
Conjecture~\ref{conj:theta-sum-B} would therefore close OP~5
completely.
\end{remark}

\begin{remark}[Connection to the explicit formula remainder]
\label{rem:explicit-formula-remainder}
By the Guinand--Weil explicit formula
(Theorem~3.3 of~\cite{brendecke2026rh}), the identity
$\Sigma(\tau) = 1 - \Theta(\tau) + R(\tau)$ holds for all $\tau > 0$.
Integrating from $0$ to $T$ (in the renormalized sense) and
substituting the known values of $\int_0^T 1\,d\tau = T$ and
$\int_0^T R\,d\tau_{\mathrm{reg}} = 3T/2 - \log T - T^2/24 - \gamma_E$,
Conjecture~\ref{conj:theta-sum-B} is equivalent to:
\[
  \int_0^T \Sigma^{(\mathrm{photon})}(\tau)\,d\tau
  \;=\; \bsym^2 \cdot T^2.
\]
This is precisely the balance condition $\eqref{eq:tr-balance}$.
Hence: \emph{Conjecture~\ref{conj:theta-sum-B} $\Leftrightarrow$
OP~5 Step~(B) $\Leftrightarrow$ Balance condition~\eqref{eq:tr-balance}.}
\end{remark}


